% Les matières plastiques — PCT 3ème — Cours signé OnBuch+
% Programme officiel MINESEC. Compiler avec : tectonic N17-matieres-plastiques.tex
\newcommand{\DOCMATIERE}{PCT}
\newcommand{\DOCNIVEAU}{3ème}
\newcommand{\DOCMODULE}{Module 3 : Chimie appliquée}
\newcommand{\DOCLECON}{N17}
\newcommand{\DOCTITRE}{Les matières plastiques}
% OnBuch+ — préambule « cours signés » ENRICHI (compilé avec tectonic / XeLaTeX)
% Système visuel « Pop » d'OnBuch+ : fond crème (jamais blanc), bordures encre
% épaisses, ombres « dures » décalées, titres Archivo Black en capitales.
% Version MATHÉMATIQUES : graphes (pgfplots/tikz), boîtes pédagogiques
% (définition, propriété, méthode, attention, le savais-tu, exercice résolu,
% exercice, TP, bilan), page de garde et signature OnBuch+ ancrée en bas.
%
% Macros attendues dans main.tex AVANT \input{preamble} :
%   \DOCMATIERE  \DOCNIVEAU  \DOCMODULE  \DOCLECON (numéro)  \DOCTITRE
\documentclass[11pt]{article}

\usepackage{fontspec}
\usepackage[french]{babel}
\usepackage[a4paper,margin=2cm,top=3.4cm,bottom=2.8cm,headheight=1.4cm,footskip=1.3cm]{geometry}
\usepackage{amsmath,amssymb,amsfonts}
\usepackage{mathtools} % \xmapsto, \coloneqq... souvent utilisés par les modèles
\usepackage{cancel} % simplifications barrées (bilans de forces)
% \abs{z} (module, valeur absolue) et \norm{u} (norme), utilisés par les modèles
\providecommand{\abs}{}\let\abs\relax\DeclarePairedDelimiter{\abs}{\lvert}{\rvert}
\providecommand{\norm}{}\let\norm\relax\DeclarePairedDelimiter{\norm}{\lVert}{\rVert}
\usepackage[table]{xcolor} % [table] : fournit aussi \rowcolors (alternance de lignes)
\usepackage{pagecolor}
\usepackage{tikz}
\usetikzlibrary{arrows.meta,positioning,calc,shapes.geometric,shapes.misc,shapes.arrows,decorations.pathmorphing,decorations.pathreplacing,decorations.markings,patterns,shadows,mindmap,trees,fit,backgrounds,matrix,intersections,angles,quotes}
\usepackage{pgfplots}
\usepackage[version=4]{mhchem} % équations nucléaires \ce{^{235}_{92}U}
\usepackage[european,siunitx]{circuitikz} % schémas électriques
\pgfplotsset{compat=1.18}
\usepgfplotslibrary{fillbetween,groupplots} % aires entre courbes (calcul intégral)
\usepackage{enumitem}
\usepackage{fancyhdr}
% [most] charge toutes les bibliothèques tcolorbox (listings, minted, raster,
% documentation...) et gonfle inutilement la mémoire TeX sur un document long
% (~300 boîtes) : on ne charge que ce qui est réellement utilisé.
\usepackage[skins,breakable]{tcolorbox}
\usepackage{graphicx}
\usepackage{colortbl}
\usepackage{booktabs}
\usepackage{tabularx}
\usepackage{multirow}
\usepackage{diagbox} % case d'en-tête diagonale des tableaux à double entrée
% Colonnes extensibles centrée / gauche / droite (C, L, R), souvent utilisées par les modèles
\newcolumntype{C}{>{\centering\arraybackslash}X}
\newcolumntype{L}{>{\raggedright\arraybackslash}X}
\newcolumntype{R}{>{\raggedleft\arraybackslash}X}
\usepackage{array}
\usepackage{multicol}
\usepackage{siunitx}
% parse-numbers=false : accepte aussi \SI{2,5 \times 10^{-3}}{...}
\sisetup{locale=FR,output-decimal-marker={,},group-minimum-digits=4,parse-numbers=false}
\DeclareSIUnit\ohm{\ensuremath{\symup{\Omega}}} % U+2126 absent de la police
\DeclareSIUnit\division{div} % réglages d'oscilloscope (ms/div, V/div)
\DeclareSIUnit\tr{tr} % tours (vitesse de rotation, ex. \SI{1500}{\tr\per\minute})
\DeclareSIUnit\molar{M} % concentration molaire (ex. \SI{3}{\molar}, \SI{1}{\micro\molar})
\DeclareSIUnit\an{an} % année (le modèle confond parfois avec \year, un compteur TeX) — ex. \SI{5}{\kilo\meter\per\an}
\DeclareSIUnit\FCFA{FCFA} % franc CFA (ex. \SI{500}{\FCFA\per\kilo\gram})
\DeclareSIUnit\degreeBrix{\textdegree Brix} % teneur en sucre d'un sirop/jus (réfractométrie)
\DeclareSIUnit\month{mois} % le modèle utilise parfois \month, un compteur TeX, comme unité de durée
\DeclareSIUnit\microliter{\micro\liter} % le modèle écrit parfois \microliter en un seul token
\DeclareSIUnit\million{millions} % dénombrement (ex. \SI{15}{\million\per\mL} spermatozoïdes)
\usepackage{qrcode}
% \degree : commande fréquemment utilisée par les modèles pour un angle ou une
% température en dehors de \SI{}{} (non fournie par les paquets chargés).
\providecommand{\degree}{^{\circ}}
% \qed (fin de démonstration), utilisé par les modèles sans amsthm
\providecommand{\qed}{\hfill$\square$}
% \barre : double barre des valeurs interdites dans un tableau de variations (tkz-tab)
\providecommand{\barre}{\ensuremath{\|}}
% environnement proof (amsthm non chargé) utilisé par les modèles
\ifdefined\proof\else\newenvironment{proof}[1][Démonstration]{\par\noindent\textit{#1.}\ }{\qed\par}\fi
\usepackage{lastpage}
\usepackage{hyperref}
\hypersetup{hidelinks}

% ============================================================
% Palette « Pop » OnBuch+ — mêmes hex que la classe P (plus_ui.dart)
% ============================================================
\definecolor{poporange}{HTML}{F59321}
\definecolor{popdark}{HTML}{E07A0C}
\definecolor{popcream}{HTML}{FFF6E8}
\definecolor{popink}{HTML}{1C1714}
\definecolor{poppurple}{HTML}{7A5AE0}
\definecolor{popgreen}{HTML}{2E9E6A}
\definecolor{poppink}{HTML}{E0457B}
\definecolor{popblue}{HTML}{2F7DE1}
\definecolor{popgold}{HTML}{C98A12}
\definecolor{popmuted}{HTML}{8A7F76}
\definecolor{popline}{HTML}{EADFD2}
\definecolor{popgray}{HTML}{8A7F76}
\colorlet{popgrey}{popgray}
% Alias tolérants (le modèle nomme parfois les couleurs différemment)
\colorlet{popwhite}{white}
\colorlet{popblack}{popink}
\colorlet{popred}{poppink}
\colorlet{popyellow}{popgold}
% Teintes claires (fonds de tableaux/encadrés) — le modèle nomme parfois
% les couleurs avec un suffixe L (ex. popblueL) sans les avoir définies.
\colorlet{poporangeL}{poporange!20}
\colorlet{popdarkL}{popdark!20}
\colorlet{poppurpleL}{poppurple!20}
\colorlet{popgreenL}{popgreen!20}
\colorlet{poppinkL}{poppink!20}
\colorlet{popblueL}{popblue!20}
\colorlet{popgoldL}{popgold!20}
\colorlet{popredL}{popred!20}
\colorlet{popgrayL}{popgray!20}
% Teintes claires pour fonds de tableaux / zones
\colorlet{popblueL}{popblue!10!white}
\colorlet{poporangeL}{poporange!14!white}
\colorlet{popgreenL}{popgreen!12!white}
\colorlet{poppurpleL}{poppurple!12!white}
\colorlet{poppinkL}{poppink!12!white}
\colorlet{popgoldL}{popgold!14!white}

\pagecolor{popcream}

% ============================================================
% Typographie
% ============================================================
\defaultfontfeatures{Ligatures=TeX}
\setmainfont{PlusJakartaSans-Medium.ttf}[Path=../../../fonts/,BoldFont=PlusJakartaSans-Bold.ttf]
% Maths en sans-serif (Fira Math) avec chiffres et lettres droites de Plus
% Jakarta, pour que les formules se fondent dans le texte.
\usepackage[math-style=ISO,bold-style=ISO]{unicode-math}
\setmathfont{FiraMath-Regular.otf}[Path=../../../fonts/]
\setmathfont{PlusJakartaSans-Medium.ttf}[Path=../../../fonts/,range={up/num,up/{Latin,latin},\mathup}]
\usepackage{icomma} % virgule décimale sans espace en mode math
% Caractères mathématiques Unicode absents des polices de texte (Plus Jakarta,
% Archivo Black) : rendus par leur équivalent mathématique, en texte comme en
% formule — sinon ils s'affichent en carré vide (ex. « Arithmétique dans ℤ »).
\usepackage{newunicodechar}
\newunicodechar{ℝ}{\ensuremath{\mathbb{R}}}
\newunicodechar{ℤ}{\ensuremath{\mathbb{Z}}}
\newunicodechar{ℂ}{\ensuremath{\mathbb{C}}}
\newunicodechar{ℕ}{\ensuremath{\mathbb{N}}}
\newunicodechar{ℚ}{\ensuremath{\mathbb{Q}}}
\newunicodechar{𝔻}{\ensuremath{\mathbb{D}}}
\newunicodechar{α}{\ensuremath{\alpha}}
\newunicodechar{β}{\ensuremath{\beta}}
\newunicodechar{γ}{\ensuremath{\gamma}}
\newunicodechar{δ}{\ensuremath{\delta}}
\newunicodechar{ε}{\ensuremath{\varepsilon}}
\newunicodechar{θ}{\ensuremath{\theta}}
\newunicodechar{λ}{\ensuremath{\lambda}}
\newunicodechar{μ}{\ensuremath{\mu}}
\newunicodechar{σ}{\ensuremath{\sigma}}
\newunicodechar{τ}{\ensuremath{\tau}}
\newunicodechar{φ}{\ensuremath{\varphi}}
\newunicodechar{ψ}{\ensuremath{\psi}}
\newunicodechar{ω}{\ensuremath{\omega}}
\newunicodechar{Σ}{\ensuremath{\Sigma}}
% majuscules produites par \MakeUppercase dans les titres (α-aminés -> Α)
\newunicodechar{Α}{\ensuremath{\alpha}}
\newunicodechar{Β}{\ensuremath{\beta}}
\newunicodechar{Γ}{\ensuremath{\Gamma}}
\newunicodechar{Δ}{\ensuremath{\Delta}}
\newunicodechar{Ε}{\ensuremath{\varepsilon}}
\newunicodechar{Θ}{\ensuremath{\Theta}}
\newunicodechar{Λ}{\ensuremath{\Lambda}}
\newunicodechar{Μ}{\ensuremath{\mu}}
\newunicodechar{Τ}{\ensuremath{\tau}}
\newunicodechar{Φ}{\ensuremath{\Phi}}
\newunicodechar{Ψ}{\ensuremath{\Psi}}
\newunicodechar{Ω}{\ensuremath{\Omega}}
\newunicodechar{↦}{\ensuremath{\mapsto}}
\newunicodechar{∈}{\ensuremath{\in}}
\newunicodechar{∉}{\ensuremath{\notin}}
\newunicodechar{∧}{\ensuremath{\wedge}}
\newunicodechar{⇔}{\ensuremath{\Leftrightarrow}}
\newunicodechar{⟺}{\ensuremath{\Longleftrightarrow}}
\newunicodechar{⇒}{\ensuremath{\Rightarrow}}
\newunicodechar{⟹}{\ensuremath{\Longrightarrow}}
\newunicodechar{∘}{\ensuremath{\circ}}
\newunicodechar{∪}{\ensuremath{\cup}}
\newunicodechar{∩}{\ensuremath{\cap}}
\newunicodechar{≡}{\ensuremath{\equiv}}
\newunicodechar{⊂}{\ensuremath{\subset}}
\newunicodechar{⊥}{\ensuremath{\perp}}
\newunicodechar{∃}{\ensuremath{\exists}}
\newunicodechar{∀}{\ensuremath{\forall}}
\newunicodechar{∛}{\ensuremath{\sqrt[3]{\,}}}
\newunicodechar{∜}{\ensuremath{\sqrt[4]{\,}}}
\newunicodechar{⃗}{\ensuremath{{}^{\to}}}
\newunicodechar{ⁿ}{\ifmmode^{n}\else\textsuperscript{\ensuremath{n}}\fi}
\newunicodechar{ˣ}{\ifmmode^{x}\else\textsuperscript{\ensuremath{x}}\fi}
\newunicodechar{ᵇ}{\ifmmode^{b}\else\textsuperscript{\ensuremath{b}}\fi}
\newunicodechar{ʳ}{\ifmmode^{r}\else\textsuperscript{\ensuremath{r}}\fi}
\newunicodechar{ᵘ}{\ifmmode^{u}\else\textsuperscript{\ensuremath{u}}\fi}
\newunicodechar{ʸ}{\ifmmode^{y}\else\textsuperscript{\ensuremath{y}}\fi}
\newunicodechar{ᵃ}{\ifmmode^{a}\else\textsuperscript{\ensuremath{a}}\fi}
\newunicodechar{ᵉ}{\ifmmode^{e}\else\textsuperscript{\ensuremath{e}}\fi}
\newunicodechar{ᵏ}{\ifmmode^{k}\else\textsuperscript{\ensuremath{k}}\fi}
\newunicodechar{ᵖ}{\ifmmode^{p}\else\textsuperscript{\ensuremath{p}}\fi}
\newunicodechar{ᴺ}{\ifmmode^{N}\else\textsuperscript{\ensuremath{N}}\fi}
\newunicodechar{ᶜ}{\ifmmode^{c}\else\textsuperscript{\ensuremath{c}}\fi}
\newunicodechar{ʰ}{\ifmmode^{h}\else\textsuperscript{\ensuremath{h}}\fi}
\newunicodechar{ᵠ}{\ifmmode^{q}\else\textsuperscript{\ensuremath{q}}\fi}
\newunicodechar{ₙ}{\ifmmode_{n}\else\textsubscript{\ensuremath{n}}\fi}
\newunicodechar{ₐ}{\ifmmode_{a}\else\textsubscript{\ensuremath{a}}\fi}
\newunicodechar{ᵢ}{\ifmmode_{i}\else\textsubscript{\ensuremath{i}}\fi}
\newunicodechar{ᵣ}{\ifmmode_{r}\else\textsubscript{\ensuremath{r}}\fi}
\newunicodechar{ₖ}{\ifmmode_{k}\else\textsubscript{\ensuremath{k}}\fi}
\newunicodechar{ᵦ}{\ifmmode_{\beta}\else\textsubscript{\ensuremath{\beta}}\fi}
\newunicodechar{ₓ}{\ifmmode_{x}\else\textsubscript{\ensuremath{x}}\fi}
\newfontfamily\popdisplay{ArchivoBlack-Regular.ttf}[Path=../../../fonts/]
\newfontfamily\poplogo{SpaceGrotesk-Bold.ttf}[Path=../../../fonts/]
\newfontfamily\popsemi{PlusJakartaSans-SemiBold.ttf}[Path=../../../fonts/]
\newfontfamily\popxbold{PlusJakartaSans-ExtraBold.ttf}[Path=../../../fonts/]
\newfontfamily\popxboldls{PlusJakartaSans-ExtraBold.ttf}[Path=../../../fonts/,LetterSpace=3.0]

\setlist[enumerate]{leftmargin=1.7em,itemsep=5pt,topsep=4pt}
\setlist[itemize]{leftmargin=1.7em,itemsep=4pt,topsep=4pt,label={\color{poporange}\textbullet}}
\setlength{\parindent}{0pt}
\setlength{\parskip}{5pt}
\renewcommand{\arraystretch}{1.25}

% pgfplots : style Pop par défaut
\pgfplotsset{
  every axis/.append style={axis line style={popink,line width=1pt},tick style={popink},
    grid style={popline,line width=0.5pt},label style={font=\small},tick label style={font=\footnotesize}},
  popcurve/.style={poporange,line width=1.8pt,no markers,smooth},
}
% Même style disponible dans un \draw TikZ simple (hors pgfplots)
\tikzset{popcurve/.style={poporange,line width=1.8pt}}
\tikzset{popgrid/.style={popline,very thin}}
\colorlet{popcurve}{poporange} % le nom du style est parfois utilisé comme couleur

% ============================================================
% Pill / squiggle
% ============================================================
\newcommand{\poppill}[3]{%
  \tikz[baseline=-0.55ex]\node[draw=popink,line width=1.2pt,rounded corners=2.6pt,%
    fill=#2,inner xsep=6.5pt,inner ysep=3pt]{%
    \popxboldls\fontsize{7}{7}\selectfont\color{#3}\MakeUppercase{#1}};%
}
\newcommand{\popsquiggle}[1]{%
  \tikz[baseline=-0.4ex]{
    \draw[#1,line width=2.6pt,line cap=round]
      plot [smooth] coordinates {(0,0.09) (0.34,0.01) (0.68,0.21) (1.02,0.01) (1.36,0.10)};
  }%
}
% Mot-clé surligné (marqueur orange sous le texte)
\newcommand{\cle}[1]{{\popxbold\color{popdark}#1}}

% ============================================================
% En-tête / pied
% ============================================================
\pagestyle{fancy}
\fancyhf{}
\renewcommand{\headrulewidth}{0pt}
\renewcommand{\footrulewidth}{0pt}
\lhead{\raisebox{-0.15ex}{\poplogo\fontsize{13}{13}\selectfont\color{popink}OnBuch\color{poporange}+}}
\rhead{\poppill{\DOCMATIERE\ \textbullet\ \DOCNIVEAU\ \textbullet\ Leçon \DOCLECON}{white}{popink}}
\lfoot{\popxbold\fontsize{7.6}{7.6}\selectfont\color{popmuted}\MakeUppercase{Cours signé OnBuch+}\ \textbullet\ {\popsemi\DOCTITRE}}
\rfoot{\tikz[baseline=-0.6ex]\node[rounded corners=8.5pt,fill=popink,minimum height=17pt,minimum width=17pt,inner xsep=5pt,inner sep=0]{\popxbold\fontsize{8}{8}\selectfont\color{popcream}\thepage\,/\,\pageref*{LastPage}};}

% ============================================================
% Titres
% ============================================================
\newcommand{\chaptitre}[2]{%
  \begin{tcolorbox}[enhanced,colback=poporange,colframe=popink,boxrule=2.6pt,arc=11pt,
      drop shadow=popink,left=15pt,right=15pt,top=12pt,bottom=11pt,before skip=3pt,after skip=18pt]
    \poppill{#2}{popink}{popcream}\par\vspace{9pt}
    {\raggedright\hyphenpenalty=10000\exhyphenpenalty=10000\popdisplay\fontsize{22}{24}\selectfont\color{popcream}\MakeUppercase{#1}\par}
  \end{tcolorbox}
}
% Section numérotée : pastille orange avec le numéro + titre Archivo
\newcounter{popsec}
\newcommand{\coursec}[1]{%
  \stepcounter{popsec}\setcounter{popsub}{0}%
  \par\vspace{16pt}\noindent
  \begin{minipage}{\linewidth}
  \tikz[baseline=-0.7ex]\node[draw=popink,line width=1.6pt,fill=poporange,rounded corners=4pt,minimum size=22pt,inner sep=0]{\popdisplay\fontsize{12}{12}\selectfont\color{popcream}\thepopsec};%
  \hspace{8pt}{\popdisplay\fontsize{15}{17}\selectfont\color{popink}\MakeUppercase{#1}}\par
  \vspace{4pt}\noindent{\color{poporange}\rule{36pt}{3.6pt}}
  \end{minipage}\par\nopagebreak\vspace{8pt}
}
\newcounter{popsub}
\newcommand{\courssub}[1]{%
  \stepcounter{popsub}%
  \par\vspace{9pt}\noindent{\popxbold\fontsize{11.5}{13}\selectfont\color{popink}{\color{poporange}\thepopsec.\thepopsub}\hspace{6pt}#1}\par\nopagebreak\vspace{4pt}
}

% ============================================================
% Boîtes pédagogiques — même recette : fond blanc, bordure épaisse colorée,
% ombre dure encre, badge plein. Titre optionnel [#1].
% ============================================================
\tcbset{popbase/.style={breakable,enhanced,colback=white,boxrule=2.3pt,arc=9pt,
  shadow={3pt}{-3pt}{0pt}{popink},left=14pt,right=14pt,top=11pt,bottom=11pt,before skip=11pt,after skip=15pt,
  fontupper=\popsemi\fontsize{10.6}{14.5}\selectfont\color{popink},
  fontlower=\popsemi\fontsize{10.6}{14.5}\selectfont\color{popink}}}
% Coche dessinée (le glyphe ✓ n'existe pas dans les polices du document)
\newcommand{\popcheck}[1]{\tikz[baseline=-0.5ex]\draw[#1,line width=1.6pt,line cap=round,line join=round](-0.13,0.0)--(-0.04,-0.09)--(0.13,0.1);}
\providecommand{\coche}{\popcheck{popgreen}}
\providecommand{\resultat}[1]{\par\smallskip\noindent\fcolorbox{popgreen}{popgreenL}{\parbox{\dimexpr\linewidth-2\fboxsep-2\fboxrule\relax}{\textbf{Résultat.} #1}}\par\smallskip}
\newcommand{\popboxhead}[3]{\poppill{#1}{#2}{white}\ifx\relax#3\relax\else\hspace{6pt}{\popxbold\fontsize{10.6}{12}\selectfont\color{popink}#3}\fi\par\vspace{7pt}}

\newtcolorbox{aretenir}[1][]{popbase,colframe=popgreen,before upper={\popboxhead{À retenir}{popgreen}{#1}}}
\newtcolorbox{exemplebox}[1][]{popbase,colframe=poporange,before upper={\popboxhead{Exemple}{poporange}{#1}}}
\newtcolorbox{definition}[1][]{popbase,colframe=popblue,before upper={\popboxhead{Définition}{popblue}{#1}}}
\newtcolorbox{propriete}[1][]{popbase,colframe=poppurple,before upper={\popboxhead{Propriété}{poppurple}{#1}}}
\newtcolorbox{methode}[1][]{popbase,colframe=popgold,before upper={\popboxhead{Méthode}{popgold}{#1}}}
\newtcolorbox{attention}[1][]{popbase,colframe=poppink,before upper={\popboxhead{Attention}{poppink}{#1}}}
\newtcolorbox{experience}[1][]{popbase,colframe=popblue,colback=popblueL,before upper={\popboxhead{Expérience}{popblue}{#1}}}
% « Le savais-tu ? » : carte violette pleine (contexte, culture, Cameroun)
\newtcolorbox{savaistu}[1][]{popbase,colback=poppurple,colframe=popink,
  fontupper=\popsemi\fontsize{10.4}{14.2}\selectfont\color{white},
  before upper={\poppill{Le savais-tu\,?}{white}{poppurple}\ifx\relax#1\relax\else\hspace{6pt}{\popxbold\color{white}#1}\fi\par\vspace{7pt}}}
% Exercice résolu : énoncé en haut, \tcblower puis la solution sur fond crème
\newcounter{popexo}\newcounter{popexr}
\newtcolorbox{exoresolu}[1][]{popbase,colframe=poporange,colbacklower=popcream,
  segmentation style={popink,line width=1.2pt,dashed},
  before upper={\stepcounter{popexr}\popboxhead{Exercice résolu \thepopexr}{poporange}{#1}},
  before lower={\poppill{Solution}{popink}{popcream}\par\vspace{6pt}}}
% Exercice (non résolu en place) — la correction est dans la partie Corrigés
\newtcolorbox{exercice}[1][]{popbase,colframe=popink,
  before upper={\stepcounter{popexo}\popboxhead{Exercice \thepopexo}{popink}{#1}}}
% Corrigé
\newtcolorbox{corrige}[1][]{popbase,colframe=popgreen,colback=popgreenL,
  before upper={\popboxhead{Corrigé}{popgreen}{#1}}}
% Objectifs / prérequis / bilan
\newtcolorbox{objectifs}{popbase,colframe=popink,colback=poporangeL,
  before upper={\popboxhead{Objectifs de la leçon}{popink}{}}}
\newtcolorbox{prerequis}{popbase,colframe=popink,colback=popblueL,
  before upper={\popboxhead{Prérequis}{popblue}{}}}
\newtcolorbox{bilan}{popbase,colframe=popink,colback=white,boxrule=2.8pt,
  before upper={\popboxhead{Fiche bilan}{poporange}{L'essentiel à connaître pour l'examen}}}
% Figure encadrée avec légende
\newtcolorbox{popfigure}[1][]{enhanced,breakable=false,colback=white,colframe=popline,boxrule=1.4pt,arc=8pt,
  left=10pt,right=10pt,top=10pt,bottom=8pt,before skip=10pt,after skip=12pt,halign=center,
  lower separated=false,halign lower=center,
  fontlower=\popsemi\fontsize{9}{11.5}\selectfont\color{popmuted},
  #1}
\newcommand{\legende}[1]{\par\vspace{6pt}{\popsemi\fontsize{9}{11.5}\selectfont\color{popmuted}#1}}

% ============================================================
% Signature OnBuch+
% ============================================================
\newcommand{\poplogoinline}[1]{{\poplogo\fontsize{#1}{#1}\selectfont\color{popink}OnBuch\color{poporange}+}}
\newcommand{\signatureonbuch}{%
  \begin{tcolorbox}[enhanced,colback=popink,colframe=popink,boxrule=2.2pt,arc=10pt,
      drop shadow=poporange,left=14pt,right=14pt,top=10pt,bottom=10pt,before skip=0pt,after skip=0pt]
    \tikz[baseline=-0.7ex]\node[circle,fill=popgreen,draw=popcream,line width=1.3pt,minimum size=22pt,inner sep=0]{\popcheck{white}};%
    \hspace{9pt}\begin{minipage}[c]{0.62\linewidth}
      {\popxbold\fontsize{12}{13}\selectfont\color{popcream}Signé\ \textbullet\ {\poplogo OnBuch\color{poporange}+}}\par\vspace{2pt}
      {\raggedright\popsemi\fontsize{8}{10.5}\selectfont\color{popline}\MakeUppercase{Équipe pédagogique OnBuch+}\\\MakeUppercase{Conforme au programme officiel MINESEC}\par}
    \end{minipage}\hfill
    {\poplogo\fontsize{22}{22}\selectfont\color{popcream}OnBuch\color{poporange}+}
  \end{tcolorbox}%
}

% ============================================================
% Page de garde
% #1 titre · #2 matière · #3 niveau · #4 module · #5 n° leçon · #6 durée ·
% #7 description / objectifs (texte ou itemize)
% ============================================================
\newcommand{\pagegarde}[7]{%
  \thispagestyle{empty}
  \newgeometry{a4paper,margin=1.7cm,top=1.6cm,bottom=1.4cm}%
  \begin{tikzpicture}[remember picture,overlay]
    \fill[poporange!18] ([xshift=-3.2cm,yshift=-3.6cm]current page.north east) circle (4.6cm);
    \fill[poppurple!14] ([xshift=2.4cm,yshift=7.6cm]current page.south west) circle (3.8cm);
  \end{tikzpicture}%
  {\poplogo\fontsize{30}{30}\selectfont\color{popink}OnBuch\color{poporange}+}\hfill
  \raisebox{0.6ex}{\poppill{Cours signé \textbullet\ Premium}{popink}{popcream}}\par
  \vspace{1pt}\popsquiggle{poppurple}\par
  \vspace{8pt}
  \begin{tcolorbox}[enhanced,colback=poporange,colframe=popink,boxrule=2.8pt,arc=13pt,
      drop shadow=popink,left=18pt,right=18pt,top=12pt,bottom=13pt,after skip=10pt]
    \poppill{#2 \textbullet\ #3}{popink}{popcream}\hspace{5pt}\poppill{#4}{white}{popink}\par\vspace{8pt}
    {\popdisplay\fontsize{14}{14}\selectfont\color{popink}LEÇON #5}\par\vspace{4pt}
    {\raggedright\hyphenpenalty=10000\exhyphenpenalty=10000\popdisplay\fontsize{25}{27}\selectfont\color{popcream}\MakeUppercase{#1}\par}
    \vspace{8pt}
    {\popxbold\fontsize{9.5}{9.5}\selectfont\color{popink}\MakeUppercase{Durée conseillée : #6}}
  \end{tcolorbox}
  \begin{tcolorbox}[enhanced,colback=white,colframe=popink,boxrule=2.2pt,arc=10pt,
      drop shadow=popink,left=16pt,right=16pt,top=10pt,bottom=10pt,after skip=8pt]
    \poppill{Dans cette leçon}{poppurple}{white}\par\vspace{5pt}
    \popsemi\fontsize{9.3}{12.6}\selectfont\color{popink}#7
  \end{tcolorbox}
  \noindent\begin{minipage}[t]{0.487\linewidth}
    \begin{tcolorbox}[enhanced,colback=popgreen,colframe=popink,boxrule=2.2pt,arc=10pt,
        drop shadow=popink,left=11pt,right=11pt,top=10pt,bottom=10pt,height=3.1cm,valign=center]
      \tcbox[colback=white,colframe=popink,boxrule=1.3pt,arc=5pt,boxsep=0pt,nobeforeafter,
        left=3pt,right=3pt,top=3pt,bottom=3pt,valign=center]{\qrcode[height=1.9cm]{https://play.google.com/store/apps/details?id=com.luvvix.onbuch&hl=fr}}%
      \hspace{8pt}\begin{minipage}[c]{0.5\linewidth}
        {\popdisplay\fontsize{11}{12.5}\selectfont\color{white}\MakeUppercase{Télécharge OnBuch}}\par\vspace{4pt}
        {\raggedright\popsemi\fontsize{8.4}{11}\selectfont\color{white}Gratuit \textbullet\ Google Play\\Tuteur IA, annales, concours, résultats\par}
      \end{minipage}
    \end{tcolorbox}
  \end{minipage}\hfill
  \begin{minipage}[t]{0.487\linewidth}
    \begin{tcolorbox}[enhanced,colback=poppurple,colframe=popink,boxrule=2.2pt,arc=10pt,
        drop shadow=popink,left=11pt,right=11pt,top=10pt,bottom=10pt,height=3.1cm,valign=center]
      \tikz[baseline=-0.7ex]\node[circle,fill=white,draw=popink,line width=1.3pt,minimum size=1.6cm,inner sep=0]{\popdisplay\fontsize{24}{24}\selectfont\color{poppurple}+};%
      \hspace{8pt}\begin{minipage}[c]{0.6\linewidth}
        {\popdisplay\fontsize{11}{12.5}\selectfont\color{white}\MakeUppercase{Passe à OnBuch+}}\par\vspace{4pt}
        {\raggedright\popsemi\fontsize{8.4}{11}\selectfont\color{white}Tous les cours signés, Tuteur IA illimité, fiches PDF — onglet OnBuch+ de l'app\par}
      \end{minipage}
    \end{tcolorbox}
  \end{minipage}
  \vspace{8pt}
  \popcontenu
  \vfill
  \signatureonbuch
  \restoregeometry
  \newpage
}

% Rangée « Ce cours contient » (page de garde)
\newcommand{\poptile}[3]{% #1 couleur #2 gros texte #3 légende
  \begin{tcolorbox}[enhanced,colback=#1,colframe=popink,boxrule=2pt,arc=9pt,shadow={2.5pt}{-2.5pt}{0pt}{popink},
      left=6pt,right=6pt,top=9pt,bottom=9pt,halign=center,valign=center,height=1.75cm,width=\linewidth,nobeforeafter]
    {\popdisplay\fontsize{12.5}{13}\selectfont\color{white}\MakeUppercase{#2}}\par\vspace{3pt}
    {\centering\popsemi\fontsize{7.8}{9.5}\selectfont\color{white}#3\par}
  \end{tcolorbox}}
\newcommand{\popcontenu}{%
  \par\noindent{\popxboldls\fontsize{7.5}{7.5}\selectfont\color{popmuted}CE COURS SIGNÉ CONTIENT}\par\vspace{7pt}
  \noindent\begin{minipage}[t]{0.235\linewidth}\poptile{popblue}{Cours}{complet, illustré, pas à pas}\end{minipage}\hfill
  \begin{minipage}[t]{0.235\linewidth}\poptile{poporange}{Méthodes}{et exercices résolus}\end{minipage}\hfill
  \begin{minipage}[t]{0.235\linewidth}\poptile{poppink}{Exercices}{type examen + corrigés}\end{minipage}\hfill
  \begin{minipage}[t]{0.235\linewidth}\poptile{popgreen}{Fiche bilan}{l'essentiel à retenir}\end{minipage}\par}

% Sommaire
\newtcolorbox{sommaire}{enhanced,colback=white,colframe=popink,boxrule=2.2pt,arc=10pt,
  drop shadow=popink,left=18pt,right=18pt,top=13pt,bottom=13pt,before skip=6pt,after skip=20pt,
  before upper={\poppill{Sommaire}{poppurple}{white}\par\vspace{9pt}\popsemi\fontsize{10.6}{14.5}\selectfont\color{popink}}}

% Signature de fin de cours — poussée tout en bas de la dernière page
\newcommand{\findecours}{%
  \par\vfill\nopagebreak
  \begin{center}\popsquiggle{poporange}\end{center}\vspace{4pt}
  \signatureonbuch
}
\AtBeginDocument{\def\llbracket{\mathopen{[\mkern-3.5mu[}}\def\rrbracket{\mathclose{]\mkern-3.5mu]}}} % intervalles d'entiers (glyphe ⟦ absent de la police)
% Glyphes absents de Fira Math : redéfinis à partir de glyphes disponibles
\AtBeginDocument{%
  \def\setminus{\mathbin{\backslash}}\let\smallsetminus\setminus
  \def\vdots{\mathord{\vbox{\baselineskip3.2pt\lineskiplimit0pt\kern4pt\hbox{.}\hbox{.}\hbox{.}}}}%
  \def\ddots{\mathinner{\mkern1mu\raise6.5pt\hbox{.}\mkern2mu\raise3.6pt\hbox{.}\mkern2mu\raise0.7pt\hbox{.}\mkern1mu}}%
  \def\triangle{\mathord{\scalebox{0.9}{$\symup{\Delta}$}}}%
  \def\nabla{\mathord{\raisebox{\depth}{\scalebox{1}[-1]{$\symup{\Delta}$}}}}%
  \def\checkmark{\popcheck{popdark}}%
  \def\star{\mathbin{\ast}}\def\bigstar{\mathord{\ast}}%
  \def\diamond{\mathbin{\scalebox{0.6}{$\Diamond$}}}%
  \def\bigcup{\mathop{\mathchoice{\raisebox{-0.3em}{\scalebox{1.7}{$\cup$}}}{\scalebox{1.25}{$\cup$}}{\cup}{\cup}}\displaylimits}%
  \def\bigcap{\mathop{\mathchoice{\raisebox{-0.3em}{\scalebox{1.7}{$\cap$}}}{\scalebox{1.25}{$\cap$}}{\cap}{\cap}}\displaylimits}%
  \def\lesseqqgtr{\mathrel{\lessgtr}}\def\bigoplus{\mathop{\oplus}\displaylimits}%
}
\providecommand{\ssi}{si et seulement si} % abréviation fréquente
\AtBeginDocument{\providecommand{\wideparen}{\overparen}} % arc de cercle

%% FIGFIX-BEGIN v4 — correctifs globaux des figures (docs/onbuch-generation/tools/figfix)
\usepackage{adjustbox}\usepackage{etoolbox}
% toute figure TikZ/pgfplots plus large que la ligne est réduite à la largeur disponible
\BeforeBeginEnvironment{tikzpicture}{\begin{adjustbox}{max width=\linewidth}}
\AfterEndEnvironment{tikzpicture}{\end{adjustbox}}
% légendes pgfplots lisibles par-dessus les courbes
\pgfplotsset{every axis/.append style={legend style={fill=white,fill opacity=0.92,text opacity=1,draw=black!15,inner sep=3pt}}}
% étiquettes d'arêtes (syntaxe edge["..."]) avec fond blanc
\tikzset{every edge quotes/.append style={fill=white,inner sep=1.2pt}}
%% FIGFIX-END
\begin{document}

\pagegarde{Les matières plastiques}{PCT}{3ème}{Module 3 : Chimie appliquée}{N17}{1 séance (1 h chacune)}{Cette leçon présente les matières plastiques : leur origine, leur définition, leur classification (thermoplastiques et thermodurcissables) ainsi que leurs usages. Elle aborde enfin la gestion des déchets plastiques (recyclage, pollution) et les gestes éco-citoyens adaptés au contexte camerounais.
\vspace{4pt}
\begin{multicols}{2}\small
\begin{itemize}
\item À la fin de cette leçon, je sais définir une matière plastique et un polymère.
\item Je sais citer les matières premières des plastiques (pétrole, gaz naturel).
\item Je sais distinguer expérimentalement un thermoplastique d'un thermodurcissable.
\item Je sais identifier les plastiques courants (PE, PVC, PET, PS, PP) et leurs usages quotidiens.
\item Je sais interpréter les symboles de recyclage présents sur les emballages.
\item Je sais expliquer pourquoi les plastiques posent un problème de pollution (non biodégradabilité).
\end{itemize}\end{multicols}}

\begin{sommaire}
\begin{multicols}{2}
\begin{enumerate}
\item Origine et définition des matières plastiques
\item Classification : thermoplastiques et thermodurcissables
\item Les plastiques courants : identification et usages
\item Avantages et limites des plastiques
\item Pollution plastique au Cameroun et dans le monde
\item Gestion des déchets plastiques : les 3R et le recyclage
\item Activité d'intégration
\item Méthodes et exercices résolus
\item Exercices
\item Corrigés des exercices
\item Fiche bilan
\end{enumerate}
\end{multicols}
\end{sommaire}

\begin{objectifs}
\begin{itemize}
\item À la fin de cette leçon, je sais définir une matière plastique et un polymère.
\item Je sais citer les matières premières des plastiques (pétrole, gaz naturel).
\item Je sais distinguer expérimentalement un thermoplastique d'un thermodurcissable.
\item Je sais identifier les plastiques courants (PE, PVC, PET, PS, PP) et leurs usages quotidiens.
\item Je sais interpréter les symboles de recyclage présents sur les emballages.
\item Je sais expliquer pourquoi les plastiques posent un problème de pollution (non biodégradabilité).
\item Je sais proposer des solutions de gestion des déchets plastiques : réduire, réutiliser, recycler, valoriser.
\item Je sais appliquer les consignes de sécurité lors du chauffage ou de la combustion de plastiques.
\end{itemize}
\end{objectifs}

\begin{prerequis}
\begin{itemize}
\item Notions de chimie organique de base : le carbone et les chaînes carbonées (hydrocarbures vus en début d'année).
\item Le pétrole et ses produits de raffinage (essence, gazole, bitume).
\item Notion de mélange et de corps pur, changements d'état (fusion, solidification).
\item Règles générales de sécurité au laboratoire (lunettes, hotte, manipulation de la flamme).
\end{itemize}
\end{prerequis}

\newpage

\coursec{Origine et définition des matières plastiques}

\courssub{Une situation du quotidien : le plastique, omniprésent et problématique}

À Yaoundé comme à Douala, après une forte pluie, les caniveaux débordent. Regarde autour de toi : des \textbf{sachets plastiques} noirs, des \textbf{bouteilles d'eau} vides, des \textbf{gobelets} jonchent le sol et bouchent les regards. L'eau stagne, les moustiques prolifèrent, le risque de paludisme et de choléra grimpe. Pourtant, ce même matériau plastique est indispensable : il embouteille l'eau \textsc{société des eaux minérales du Cameroun} (\textsc{semc}), il forme les \textbf{chaises} monoblocs de nos cours, les \textbf{seaux} et \textbf{bassines} de nos douches, les \textbf{tuyaux} d'adduction d'eau de la \textsc{camwater}, les \textbf{câbles} électriques \textsc{eneo}, le \textbf{carénage} de nos motos-taxi \emph{benskineurs} et même le \textbf{tableau de bord} des taxis.

\medskip
\noindent
\textbf{Problématique :} D'où vient cette matière aux mille visages ? Pourquoi certains plastiques ramollissent sous le soleil de mars alors que d'autres noircissent et sentent le brûlé sans fondre ? Comment les produire, les utiliser et surtout les jeter sans étouffer notre environnement ? C'est à ces questions que cette première section répond.

\courssub{De la matière première au monomère : le pétrole, source carbonée}

\begin{experience}[Origine industrielle des plastiques : la filière pétrochimique]
\textbf{Observation (documentaire).} Au Cameroun, la \textsc{sonara} (Société Nationale de Raffinage) à Limbé transforme le pétrole brut en essence, gasoil, kérosène, butane/propane (gaz domestique) et \textbf{naphta}. Ce naphta est la matière première clé de la pétrochimie.

\textbf{Interprétation.} Le pétrole brut est un mélange d'hydrocarbures (molécules formées seulement de carbone C et d'hydrogène H). Le \textbf{raffinage} sépare ces hydrocarbures par \textbf{distillation fractionnée} selon leur point d'ébullition. La coupe légère, le \textbf{naphta} (C$_5$ à C$_{10}$), est envoyée vers un \textbf{vapocraqueur} (four à $\approx$ 800--900\,°C sous vapeur d'eau). La chaleur casse les grosses molécules en plus petites : c'est le \textbf{craquage}. On obtient ainsi des \textbf{monomères} légers : éthylène (C$_2$H$_4$), propylène (C$_3$H$_6$), butadiène, benzène, toluène, xylènes.

\textbf{Conclusion.} Les plastiques sont \textbf{issus du pétrole} (et du gaz naturel) via le raffinage puis le vapocraquage. Les monomères ainsi produits sont les ``briques'' de base qui, assemblées, formeront les polymères.
\end{experience}

\begin{popfigure}
\begin{tikzpicture}[scale=0.9, every node/.style={font=\small}]
% Couleurs
\colorlet{petrol}{popdark}
\colorlet{refine}{popblue}
\colorlet{crack}{poporange}
\colorlet{mono}{popgreen}
\colorlet{poly}{poppurple}
\colorlet{obj}{poppink}

% Boîtes
\node[draw=petrol, fill=petrol!10, rounded corners, thick, minimum width=3.5cm, minimum height=1.2cm, align=center] (pet) at (0,0) {\textbf{Pétrole brut}\\(gisement / import)};
\node[draw=refine, fill=refine!10, rounded corners, thick, minimum width=3.5cm, minimum height=1.2cm, align=center] (raf) at (5,0) {\textbf{Raffinage}\\(Distillation fractionnée)};
\node[draw=crack, fill=crack!10, rounded corners, thick, minimum width=3.5cm, minimum height=1.2cm, align=center] (cra) at (10,0) {\textbf{Vapocraquage}\\(Craquage thermique $\approx$ 850\,°C)};
\node[draw=mono, fill=mono!10, rounded corners, thick, minimum width=3.5cm, minimum height=1.2cm, align=center] (mon) at (15,0) {\textbf{Monomères}\\(Éthylène, Propylène,\\ Benzène, Butadiène...)};
\node[draw=poly, fill=poly!10, rounded corners, thick, minimum width=3.5cm, minimum height=1.2cm, align=center] (pol) at (10,-3) {\textbf{Polymérisation}\\(Réaction en chaîne,\\ catalyseurs Ziegler-Natta)};
\node[draw=obj, fill=obj!10, rounded corners, thick, minimum width=3.5cm, minimum height=1.2cm, align=center] (objn) at (10,-6) {\textbf{Objets plastiques}\\(Granulés $\to$ Extrusion,\\ Injection, Soufflage)};

% Flèches horizontales
\draw[->, thick, petrol] (pet.east) -- node[above] {Pétrole brut} (raf.west);
\draw[->, thick, refine] (raf.east) -- node[above] {Naphta} (cra.west);
\draw[->, thick, crack] (cra.east) -- node[above] {Monomères purs} (mon.west);

% Flèches verticales
\draw[->, thick, mono] (mon.south) -- node[right] {Éthylène, Propylène...} (pol.north);
\draw[->, thick, poly] (pol.south) -- node[right] {Polymères (granulés)} (objn.north);

% Produits secondaires raffinage
\node[draw=popmuted, fill=popcream, rounded corners, font=\tiny, align=center] (ess) at (5,2.2) {Essence, Gasoil,\\ Kérosène, Gaz (GPL)};
\draw[->, dashed, popmuted] (raf.north) -- (ess.south);

% Exemples monomères
\node[draw=popmuted, fill=popcream, rounded corners, font=\tiny, align=center] (exm) at (15,2.2) {C$_2$H$_4$, C$_3$H$_6$,\\ C$_6$H$_6$, C$_4$H$_6$};
\draw[->, dashed, popmuted] (mon.north) -- (exm.south);

% Exemples objets
\node[draw=popmuted, fill=popcream, rounded corners, font=\tiny, align=center] (exo) at (15,-6) {Sachets (PE), Bouteilles (PET),\\(PP), Tuyaux (PVC), Mousse (PS)};
\draw[->, dashed, popmuted] (objn.east) -- (exo.west);
\end{tikzpicture}
\legende{La filière pétrochimique : du pétrole brut à l'objet plastique. Au Cameroun, le naphta de la SONARA alimente cette chaîne.}
\end{popfigure}

\courssub{Qu'est-ce qu'une matière plastique ? Définition rigoureuse}

\begin{definition}[Matière plastique]
Une \cle{matière plastique} est un \textbf{matériau} constitué principalement d'un \cle{polymère} (macromolécule organique de masse molaire très élevée, formée par la répétition d'unités identiques ou voisines appelées \cle{monomères}) auquel on incorpore des \cle{additifs} (colorants, plastifiants, charges minérales, stabilisants UV, retardateurs de flamme) pour modifier ses propriétés physiques, mécaniques, thermiques ou esthétiques.
\end{definition}

\begin{aretenir}
\begin{itemize}
\item \textbf{Polymère} = ``nombreuses parties'' (grec \emph{polus} = nombreux, \emph{meros} = partie). C'est une chaîne géante (macromolécule) faite de l'enchaînement de milliers de monomères.
\item \textbf{Monomère} = ``une seule partie''. Petite molécule réactive (souvent insaturée : double liaison C=C) capable de s'enchaîner.
\item \textbf{Additifs} = ingrédients « secrets » (5 à 50\,\% en masse) qui transforment le polymère brut en matériau utilisable (souplesse, couleur, résistance au feu, au soleil, coût réduit).
\end{itemize}
\end{aretenir}

\courssub{La polymérisation : construire la chaîne, maillon par maillon}

\begin{methode}[Comprendre la polymérisation (démarche qualitative)]
\begin{enumerate}
\item \textbf{Identifier le monomère :} une petite molécule insaturée (double liaison C=C). Exemple : l'\textbf{éthylène} (éthène), CH$_2$=CH$_2$, gaz à l'odeur sucrée, produit par vapocraquage.
\item \textbf{Ouvrir la double liaison :} sous l'effet de la chaleur, de la pression et d'un \textbf{catalyseur}, la liaison $\pi$ se rompt. Chaque carbone garde un électron libre : c'est un \textbf{radical libre} (ou un site actif ionique/de coordination).
\item \textbf{Enchaîner :} le radical libre attaque un autre monomère, forme une liaison simple C--C, et le radical libre se déplace au bout de la chaîne naissante.
\item \textbf{Répéter des milliers de fois :} la chaîne grandit jusqu'à terminaison (couplage de deux radicaux, transfert, etc.). On obtient le \textbf{polyéthylène} (PE).
\item \textbf{Écrire l'équation globale :} $n$ monomères $\to$ 1 macromolécule. Le $n$ (degré de polymérisation) vaut 10\,000 à 100\,000 pour le PE courant.
\end{enumerate}
\end{methode}

\begin{popfigure}
\begin{tikzpicture}[scale=0.7, every node/.style={font=\small}]
% Monomères initiaux
\foreach \i in {0,1,2,3} {
  \begin{scope}[xshift=\i*3.5cm]
    \draw[thick, popblue] (0,0) -- (1,0) node[midway, above=2pt] {C} -- (2,0) node[midway, above=2pt] {C};
    \draw[thick, popblue] (0,0) -- (-0.3,-0.3) (0,0) -- (0.3,-0.3);
    \draw[thick, popblue] (1,0) -- (0.7,-0.3) (1,0) -- (1.3,-0.3);
    \draw[thick, popblue] (2,0) -- (1.7,-0.3) (2,0) -- (2.3,-0.3);
    \node[popblue] at (1,0.5) {H$_2$C=CH$_2$};
  \end{scope}
}
% Flèche réaction
\draw[->, very thick, poporange] (14.5,0.5) -- node[above, align=center]{\textbf{Polymérisation}\\(chaleur, pression, catalyseur)} (17,0.5);

% Polymère resulting (simplified representation)
\begin{scope}[xshift=18.5cm]
\foreach \j in {0,1,2,3,4,5} {
  \draw[thick, poppurple] (\j*1.2,0) -- (\j*1.2+1.2,0);
  \draw[thick, poppurple] (\j*1.2,0) -- (\j*1.2-0.15,-0.2) (\j*1.2,0) -- (\j*1.2+0.15,-0.2);
  \draw[thick, poppurple] (\j*1.2+1.2,0) -- (\j*1.2+1.05,-0.2) (\j*1.2+1.2,0) -- (\j*1.2+1.35,-0.2);
}
\node[poppurple] at (3.6,0.6) {\textbf{--[CH$_2$--CH$_2$]$_n$--}};
\node[poppurple, align=center] at (3.6,-1.2) {Chaîne de \textbf{polyéthylène (PE)}\\(masse molaire $\sim$ 10$^5$--10$^6$ g/mol)};
\end{scope}
\end{tikzpicture}
\legende{Schéma de principe : polymérisation de l'éthylène en polyéthylène. La double liaison C=C s'ouvre pour former des liaisons simples C--C entre monomères. Les atomes d'hydrogène ne sont pas tous représentés pour la clarté.}
\end{popfigure}

\begin{propriete}[Équation de polymérisation de l'éthylène]
L'équation bilan s'écrit avec la notation des unités répétitives :
\[
\ce{n CH2=CH2 ->[- Catalyseur, T, P] -[CH2-CH2]_n-}
\]
\begin{itemize}
\item $n$ = degré de polymérisation (nombre d'unités répétitives), typiquement $10^4$ à $10^5$.
\item L'unité répétitive est \ce{-[CH2-CH2]-}. Elle \textbf{n'a pas} de double liaison : la chaîne est \textbf{saturée}.
\item La masse molaire du polymère $M_{\text{pol}} \approx n \times M_{\text{monomère}}$ ($M_{\text{C2H4}} = 28$ g/mol).
\end{itemize}
\end{propriete}

\begin{exemplebox}[Le polyéthylène (PE) : le plastique le plus produit au monde]
\begin{itemize}
\item \textbf{Origine :} éthylène (gaz de craquage du naphta).
\item \textbf{Structure :} chaîne carbonée linéaire (PEHD -- haute densité) ou ramifiée (PEBD -- basse densité).
\item \textbf{Propriétés :} inerte chimiquement, hydrophobe, bon isolant électrique, thermoplastique (fond à $\approx$ 110--130\,°C).
\item \textbf{Usages au Cameroun :} sachets « \emph{pain chaud} » (PEBD), bouteilles de lait/jus, bidons d'huile, jerricanes d'eau 20\,L, tuyaux d'irrigation, géomembranes (PEHD).
\end{itemize}
\end{exemplebox}

\courssub{Plastiques naturels et plastiques synthétiques : une distinction nécessaire}

\begin{definition}[Plastique naturel vs synthétique]
\begin{itemize}
\item \textbf{Plastique naturel :} polymère existant dans la nature, utilisable directement ou après modification légère. Exemples : \textbf{caoutchouc naturel} (latex d'hévéa, polyisopène \emph{cis}-1,4), \textbf{ambre} (résine fossilisée), \textbf{corne}, \textbf{écaille de tortue}, \textbf{cellulose} (bois, coton), \textbf{amidon} (maïs, manioc).
\item \textbf{Plastique synthétique :} polymère obtenu par synthèse chimique à partir de monomères d'origine fossile (pétrole, gaz, charbon) ou, de plus en plus, biosourcés (éthanol de canne à sucre $\to$ éthylène « vert » $\to$ PE « vert »).
\end{itemize}
\end{definition}

\begin{savaistu}[Étymologie et histoire]
Le mot \textbf{« plastique »} vient du grec \emph{plastikos} (πλαστικός), signifiant « qui peut être modelé, façonné », lui-même de \emph{plassein} (πλάσσειν) : « modeler, former ». Le premier plastique \textbf{synthétique} industriel est la \textbf{Bakélite} (phénol-formol), inventée en 1907 par Leo Baekeland : c'est un \textbf{thermodurcissable} (voir section 2). Avant lui, on utilisait le \textbf{celluloïd} (nitrate de cellulose + camphre, 1869) pour les boules de billard, le cinéma et les peignes — très inflammable !
\end{savaistu}

\begin{attention}[Pièges de vocabulaire à éviter absolument]
\begin{enumerate}
\item \textbf{Ne pas confondre « polymère » et « plastique ».} Le polymère est la molécule géante (le principe actif). Le plastique est le \textbf{matériau fini} = polymère + additifs. Un même polymère (ex. PVC) donne des plastiques rigides (tuyaux, fenêtres) ou souples (tuyaux d'arrosage, nappes, câbles) selon les plastifiants ajoutés.
\item \textbf{« Plastique » $\neq$ « Matière plastique » au sens large.} En langage courant, « plastique » désigne souvent les thermoplastiques d'emballage. En science des matériaux, on distingue généralement trois grandes familles de matériaux polymères : les \textbf{plastiques} (rigides ou souples, mis en forme par moulage), les \textbf{élastomères} (caoutchoucs, élasticité réversible) et les \textbf{fibres synthétiques} (polyester, nylon, filage). Tous sont des polymères + additifs, mais leurs usages et propriétés macroscopiques diffèrent.
\item \textbf{Ne pas écrire « molécule de plastique ».} On dit : « macromolécule de polymère », « chaîne polymère », « granulé de matière plastique ».
\end{enumerate}
\end{attention}

\courssub{Exercice résolu : du monomère au polymère, calcul de masse molaire}

\begin{exoresolu}[De l'éthylène au polyéthylène : ordre de grandeur]
\textbf{Énoncé.} On réalise une polymérisation d'éthylène $\ce{C2H4}$ ($M = 28,0$ g/mol) qui conduit à un polyéthylène de degré de polymérisation moyen $n = 50\,000$.
\begin{enumerate}
\item Écrire l'équation de la réaction de polymérisation.
\item Calculer la masse molaire moyenne $M_{\text{PE}}$ du polyéthylène obtenu.
\item Un granulé de PE pèse 2,5 mg. Combien de chaînes polymères contient-il environ ?
\end{enumerate}

\tcblower
\textbf{Solution détaillée.}
\begin{enumerate}
\item \textbf{Équation :} $\ce{n CH2=CH2 -> -[CH2-CH2]_n-}$ (la double liaison s'ouvre, pas d'atome perdu, pas de sous-produit : polymérisation par \textbf{addition}).
\item \textbf{Masse molaire du polymère :}
\[
M_{\text{PE}} = n \times M_{\text{monomère}} = 50\,000 \times 28,0\ \text{g/mol} = 1\,400\,000\ \text{g/mol} = \SI{1,4e6}{g/mol}.
\]
C'est une masse molaire \textbf{très grande} (typique des plastiques usuels : $10^5$ à $10^6$ g/mol).
\item \textbf{Nombre de chaînes dans un granulé :}
Masse d'un granulé $m = \SI{2,5}{mg} = \SI{2,5e-3}{g}$.
Quantité de matière de polymère : $n_{\text{PE}} = \frac{m}{M_{\text{PE}}} = \frac{2,5 \times 10^{-3}}{1,4 \times 10^6} \approx \SI{1,79e-9}{mol}$.
Nombre de chaînes $N = n_{\text{PE}} \times N_A$ ($N_A \approx 6,02 \times 10^{23}\ \text{mol}^{-1}$).
\[
N \approx 1,79 \times 10^{-9} \times 6,02 \times 10^{23} \approx \mathbf{1,1 \times 10^{15}\ \text{chaînes}}.
\]
Un minuscule granulé contient plus de \textbf{1 000 milliards de chaînes} polymères ! Cela explique les propriétés macroscopiques (solidité, viscosité fondue) issues de l'enchevêtrement de ces géants moléculaires.
\end{enumerate}
\end{exoresolu}

\courssub{Bilan de la section 1 : ce qu'il faut retenir}

\begin{aretenir}[Synthèse : Origine et définition]
\begin{itemize}
\item \textbf{Origine :} Pétrole (raffiné à la \textsc{sonara}) $\xrightarrow{\text{vapocraquage}}$ Monomères (éthylène, propylène...) $\xrightarrow{\text{polymérisation}}$ Polymères $\xrightarrow{\text{+ additifs + mise en forme}}$ Objets plastiques.
\item \textbf{Matière plastique} = \textbf{Polymère} (macromolécule = répétition de monomères) + \textbf{Additifs}.
\item \textbf{Polymérisation :} réaction d'enchaînement de monomères (souvent par ouverture de double liaison C=C) sans sous-produit (polymérisation par addition).
\item \textbf{Exemple type :} $n\,\ce{CH2=CH2} \to \ce{-[CH2-CH2]_n-}$ (Polyéthylène).
\item \textbf{Naturels vs Synthétiques :} Caoutchouc, ambre, cellulose = naturels ; PE, PP, PVC, PS, PET = synthétiques (majorité actuelle).
\end{itemize}
\end{aretenir}

\begin{exercice}[Entraînement : vocabulaire et origine]
\begin{enumerate}
\item Complète : Le naphta, fraction légère du pétrole, est transformé en monomères par \underline{\hspace{4cm}}.
\item Vrai ou Faux ? Justifie.
\begin{itemize}
\item[a)] Une matière plastique est une substance pure.
\item[b)] Le polyéthylène est un polymère naturel.
\item[c)] La polymérisation de l'éthylène libère de l'eau.
\end{itemize}
\item Le mot « plastique » vient du grec \emph{plastikos}. Que signifie-t-il ? Donne un exemple d'objet qui illustre ce sens au quotidien.
\item Schématise (schéma simplifié au feutre) l'enchaînement de 4 monomères de propylène $\ce{CH2=CH-CH3}$ pour former un bout de chaîne de polypropylène (PP). Indique l'unité répétitive.
\end{enumerate}
\end{exercice}

\begin{corrige}[Exercice : Entraînement]
\begin{enumerate}
\item \textbf{Vapocraquage} (craquage thermique à haute température sous vapeur d'eau).
\item \textbf{Vrai/Faux :}
\begin{itemize}
\item[a)] \textbf{Faux.} C'est un \textbf{mélange} (matériau) : polymère + additifs (colorants, plastifiants...).
\item[b)] \textbf{Faux.} C'est un polymère \textbf{synthétique} issu de l'éthylène (pétrole).
\item[c)] \textbf{Faux.} C'est une polymérisation par \textbf{addition} : ouverture de la double liaison, \textbf{pas de sous-produit} (ni eau, ni autre).
\end{itemize}
\item \emph{Plastikos} = « qui peut être modelé ». Exemple : une \textbf{chaise monobloc} moulée par injection de PP fondu dans un moule ; un \textbf{sachet} soufflé à partir d'une baudruche de PE fondu.
\item \textbf{Schéma attendu :}
$\ce{CH2=CH-CH3 ->[-] -[CH2-CH(CH3)]_n-}$
Unité répétitive : \ce{-[CH2-CH(CH3)]-} (groupement \ce{-CH3} porté par un carbone sur deux). Chaîne saturée (liaisons simples uniquement).
\end{enumerate}
\end{corrige}

\coursec{Classification : thermoplastiques et thermodurcissables}

Nous venons de voir que les matières plastiques sont des polymères synthétiques, issus principalement du pétrole. Mais tous les plastiques ne se comportent pas de la même manière face à la chaleur. C'est ce comportement thermique qui va nous permettre de les classer en deux grandes familles aux propriétés radicalement différentes. Comprendre cette distinction est essentiel pour choisir le bon matériau pour un objet donné (une bouteille d'eau ou une prise électrique ?) et pour savoir comment le recycler en fin de vie.

\courssub{Une expérience révélatrice : le test du chauffage}

Pour découvrir la différence fondamentale entre les plastiques, réalisons (ou imaginons) une expérience simple au laboratoire, sous hotte aspirante et avec des lunettes de protection.

\begin{experience}[Comportement à la chaleur de deux plastiques courants]
\textbf{Matériel :} brûleur Bunsen (ou chalumeau de laboratoire), pince à creuset, deux petits échantillons : un morceau de sachet plastique fin (polyéthylène, PE) et un morceau de prise électrique cassée (bakélite, résine phéno-formolique), spatule métallique, carreau de céramique.
\textbf{Protocole :}
\begin{enumerate}
    \item Placer l'échantillon de sachet PE sur le carreau de céramique. Le chauffer \emph{doucement} avec la flamme bleue du brûleur (pointe de la flamme) en le maintenant avec la pince. Observer.
    \item Laisser refroidir l'échantillon. Tenter de le plier.
    \item Recommencer l'opération avec l'échantillon de bakélite (prise électrique). Le chauffer de la même manière. Observer l'odeur, l'aspect, la forme.
    \item Laisser refroidir. Tenter de le plier ou de le gratter.
\end{enumerate}
\textbf{Observations :}
\begin{itemize}
    \item \textbf{Sachet PE :} Il ramollit rapidement, devient mou comme de la cire, fond en gouttes visqueuses. En refroidissant, il redevient dur et solide. On peut le refaire fondre plusieurs fois.
    \item \textbf{Bakélite :} Elle ne ramollit pas. Elle noircit à la surface, dégage une odeur forte et désagréable (phénol), se craquèle et se décompose en poudre charbonneuse. Elle ne fond \emph{jamais}. Une fois refroidie, elle reste dure et cassante.
\end{itemize}
\textbf{Interprétation :} Ces deux plastiques réagissent différemment à l'apport d'énergie thermique. Le premier change d'état physique (solide $\leftrightarrow$ liquide) de façon \textbf{réversible}. Le second subit une transformation chimique irréversible (décomposition, carbonisation).
\legende{Expérience comparative : le polyéthylène fond et se solidifie à volonté ; la bakélite se dégrade sans fondre.}
\end{experience}

\begin{attention}[Sécurité impérative]
Ne \emph{jamais} chauffer un plastique dans une flamme nue sans hotte aspirante ni protection respiratoire adaptée. La combustion ou la pyrolyse des plastiques dégage des gaz toxiques (monoxyde de carbone, cyanure d'hydrogène pour les plastiques azotés, chlorure d'hydrogène pour le PVC, dioxines, furanes). Ne \emph{jamais} brûler des déchets plastiques à l'air libre (cours, poubelles) : c'est un danger grave pour la santé et l'environnement.
\end{attention}

\courssub{Les deux grandes familles : définitions et structure microscopique}

L'expérience nous conduit à définir deux catégories majeures.

\begin{definition}[Thermoplastique]
Un \textbf{thermoplastique} est un polymère qui \textbf{ramollit à la chaleur} (transition vitreuse ou fusion) et \textbf{durcit au refroidissement}, de façon \textbf{réversible} et reproductible indéfiniment (tant que la température ne dépasse pas la température de dégradation).
\end{definition}

\begin{definition}[Thermodurcissable]
Un \textbf{thermodurcissable} (ou \emph{thermoréticulable}) est un polymère qui, lors de sa fabrication (moulage), \textbf{durcit définitivement} par réticulation chimique. Une fois durci, il \textbf{ne fond plus} : chauffé, il se \textbf{dégrade} (carbonisation) sans passer par l'état liquide.
\end{definition}

Pour comprendre \emph{pourquoi} ces comportements sont si différents, il faut regarder à l'échelle microscopique, au niveau des chaînes polymères.

\begin{propriete}[Structure microscopique et comportement macroscopique]
\begin{itemize}
    \item \textbf{Thermoplastiques :} Les chaînes polymères sont \textbf{linéaires} ou \textbf{ramifiées}, mais \textbf{indépendantes} les unes des autres. Elles sont enchevêtrées comme un plat de spaghettis. Les forces intermoléculaires (forces de Van der Waals, liaisons hydrogène) maintiennent l'ensemble. À chaud, l'agitation thermique vainc ces forces : les chaînes glissent les unes sur les autres $\rightarrow$ le matériau coule (fusion). Au froid, les forces reprennent le dessus $\rightarrow$ le matériau durcit. C'est un changement d'état \textbf{physique}, réversible.
    \item \textbf{Thermodurcissables :} Lors du moulage (souvent sous chaleur et pression), les chaînes pré-polymères (résines) réagissent entre elles pour former des \textbf{liaisons covalentes fortes} (ponts) : c'est la \textbf{réticulation}. On obtient un \textbf{réseau tridimensionnel (3D) unique}, une seule macromolécule géante. Il n'y a plus de chaînes mobiles. Chauffer apporte de l'énergie qui casse les liaisons covalentes du squelette carboné $\rightarrow$ \textbf{décomposition chimique} irréversible (pyrolyse).
\end{itemize}
\end{propriete}

\begin{popfigure}
\begin{tikzpicture}[scale=0.9, every node/.style={font=\small}]
    % Thermoplastique - Gauche
    \begin{scope}[xshift=-4.5cm]
        \node[align=center, font=\bfseries\large, text=popblue] at (0,4.8) {Thermoplastique};
        \node[align=center, font=\small] at (0,4.3) {Chaînes linéaires / ramifiées};
        \node[align=center, font=\small] at (0,4.0) {Indépendantes (forces Van der Waals)};
        
        % Froid
        \node[align=center, font=\bfseries\normalsize, text=popdark] at (-1.5,3.0) {À froid (solide)};
        \foreach \i/\y in {1/2.5, 2/2.1, 3/1.7, 4/1.3, 5/0.9} {
            \draw[popblue, thick, decorate, decoration={snake, amplitude=0.5mm, segment length=3mm}] 
                (-2.5, \y) -- (2.5, \y);
        }
        \draw[popblue!50, very thick] (-3,0.5) -- (3,0.5);
        
        % Flèche chauffage
        \draw[poporange, very thick, ->] (0,0.2) -- (0,-0.5) node[midway, right] {Chaleur};
        
        % Chaud
        \node[align=center, font=\bfseries\normalsize, text=popdark] at (-1.5,-1.3) {À chaud (fondu)};
        \foreach \i/\y in {1/-1.8, 2/-2.2, 3/-2.6, 4/-3.0, 5/-3.4} {
            \draw[poporange, thick] 
                (-2.5, \y) .. controls (-1.5,\y+0.2) and (1.5,\y-0.2) .. (2.5, \y);
        }
        \draw[poporange!50, very thick] (-3,-3.8) -- (3,-3.8);
        
        % Flèche refroidissement
        \draw[popgreen, very thick, ->] (0,-4.1) -- (0,-4.8) node[midway, right] {Refroidissement};
        
        % Retour froid
        \node[align=center, font=\bfseries\normalsize, text=popdark] at (-1.5,-5.7) {Retour à froid (solide)};
        \foreach \i/\y in {1/-6.2, 2/-6.6, 3/-7.0, 4/-7.4, 5/-7.8} {
            \draw[popblue, thick, decorate, decoration={snake, amplitude=0.5mm, segment length=3mm}] 
                (-2.5, \y) -- (2.5, \y);
        }
    \end{scope}
    
    % Thermodurcissable - Droite
    \begin{scope}[xshift=5.5cm]
        \node[align=center, font=\bfseries\large, text=poporange] at (0,4.8) {Thermodurcissable};
        \node[align=center, font=\small] at (0,4.3) {Réseau 3D réticulé};
        \node[align=center, font=\small] at (0,4.0) {Liaisons covalentes (ponts)};
        
        % Avant moulage (résine)
        \node[align=center, font=\bfseries\normalsize, text=popdark] at (1.5,3.0) {Résine liquide (avant moulage)};
        \foreach \x in {-1.5,-0.5,0.5,1.5} {
            \foreach \y in {2.5,2.1,1.7,1.3} {
                \fill[poporange!70] (\x, \y) circle (0.08);
            }
        }
        \draw[poporange!50, very thick] (-2,0.9) -- (2,0.9);
        
        % Flèche moulage/chaleur
        \draw[poporange, very thick, ->] (0,0.6) -- (0,-0.1) node[midway, right] {Moulage + Chaleur};
        
        % Après moulage (réseau)
        \node[align=center, font=\bfseries\normalsize, text=popdark] at (1.5,-1.1) {Objet final (réseau 3D rigide)};
        % Dessiner un réseau ordonné
        \foreach \x in {-1.5,-0.5,0.5,1.5} {
            \foreach \y in {-1.6,-2.6,-3.6} {
                \fill[poporange!90] (\x, \y) circle (0.06);
            }
        }
        \foreach \x in {-1.5,-0.5,0.5,1.5} {
            \foreach \y in {-1.6,-2.6} {
                \draw[poporange!70, thick] (\x, \y) -- (\x, \y-1.0);
                \ifdim\x pt<1.5pt \draw[poporange!70, thick] (\x, \y) -- (\x+1.0, \y); \fi
            }
        }
        % Diagonales pour montrer 3D
        \draw[poporange!50, thick] (-1.5,-1.6) -- (-0.5,-2.6);
        \draw[poporange!50, thick] (-0.5,-1.6) -- (0.5,-2.6);
        \draw[poporange!50, thick] (0.5,-1.6) -- (1.5,-2.6);
        \draw[poporange!50, thick] (-1.5,-2.6) -- (-0.5,-3.6);
        
        \draw[poporange!50, very thick] (-2,-4.1) -- (2,-4.1);
        
        % Flèche surchauffe
        \draw[popdark, very thick, ->] (0,-4.4) -- (0,-5.1) node[midway, right] {Surchauffe};
        
        % Dégradation
        \node[align=center, font=\bfseries\normalsize, text=popdark] at (1.5,-5.8) {Dégradation (carbonisation)};
        \foreach \i in {1,...,12} {
            \pgfmathsetmacro{\rx}{-1.5 + 3*rnd}
            \pgfmathsetmacro{\ry}{-6.3 + 0.8*rnd}
            \fill[popmuted!80!black] (\rx, \ry) circle (0.05);
        }
        \node[text=popdark, font=\small] at (0,-6.8) {Irréversible : poudre noire, gaz toxiques};
    \end{scope}
\end{tikzpicture}
\legende{Modélisation microscopique : les chaînes indépendantes du thermoplastique glissent à chaud (réversible) ; le réseau réticulé du thermodurcissable est figé définitivement et casse sous l'effet de la chaleur.}
\end{popfigure}

\courssub{Exemples courants et identification}

Voici les plastiques que tu croises chaque jour au Cameroun, classés par famille.

\begin{center}
\begin{tabularx}{\linewidth}{|l|X|X|}\hline
\rowcolor{popblueL}
\textbf{Famille} & \textbf{Exemples d'objets (Contexte Cameroun)} & \textbf{Code / Symbole} \\ \hline
\textbf{Thermoplastiques} 
(Recyclables par refonte) 
& Sachets « \emph{500F} » ou « \emph{1000F} » (PEBD), bouteilles d'eau/jus (PET), bouchons, bassines, seaux, chaises en plastique (PP), tuyaux PVC (évacuation), gaines électriques (PVC), pots de yaourt (PS), emballages alimentaires rigides. 
& \ce{PET} (1), \ce{PE-HD} (2), \ce{PVC} (3), \ce{PE-BD} (4), \ce{PP} (5), \ce{PS} (6) \\ \hline
\textbf{Thermodurcissables} 
(Non recyclables par refonte) 
& Prises électriques, interrupteurs, boîtiers de compteurs ENEO (bakélite), poignées de casseroles/poêles/marmite (phénolique / mélamine), plaques de circuit imprimé (époxy), colle forte (résine époxy bi-composant), vernis, laques, pneus (caoutchouc vulcanisé \emph{proche parent}). 
& Bakélite (PF), Résines Époxy (EP), Mélamine (MF), Urée-formol (UF) \\ \hline
\end{tabularx}
\end{center}

\begin{exemplebox}[Pourquoi la poignée de la casserole est-elle en bakélite ?]
Imagine une casserole en aluminium sur le feu à bois ou au gaz. La température dépasse \SI{200}{\celsius}. Si la poignée était en polypropylène (PP, thermoplastique, fusion $\approx$ \SI{160}{\celsius}), elle fondrait et tu te brûlerais gravement. La bakélite (thermodurcissable) ne fond \emph{jamais} ; elle résiste à la chaleur sans se déformer (jusqu'à sa dégradation vers \SI{300}{\celsius}). C'est un choix de \textbf{sécurité} dicté par la classification.
\end{exemplebox}

\courssub{Méthode d'identification et conséquences sur le recyclage}

\begin{methode}[Identifier la famille d'un plastique inconnu (Test au chauffage doux)]
\textbf{Attention :} À réaliser \emph{uniquement} en milieu scolaire (hotte, lunettes, petite quantité) ou par un professionnel. Ne pas inhaler les fumées.
\begin{enumerate}
    \item Prendre un tout petit éclat du plastique (taille tête d'allumette).
    \item Le déposer sur une spatule métallique ou un carreau de céramique.
    \item Approcher \emph{doucement} de la flamme bleue d'un brûleur (pointe de flamme), sans mettre dans la flamme.
    \item \textbf{Observation :}
    \begin{itemize}
        \item Ramollissement $\rightarrow$ fusion $\rightarrow$ goutte visqueuse $\rightarrow$ durcissement au refroidissement : \textbf{Thermoplastique}.
        \item Pas de ramollissement, noircissement, odeur de brûlé/phénol, craquellement, poudre : \textbf{Thermodurcissable}.
    \end{itemize}
    \item Confirmer si possible par le code de recyclage (triangle + chiffre) moulé sur l'objet.
\end{enumerate}
\end{methode}

\begin{aretenir}[Conséquence majeure : le recyclage]
\begin{itemize}
    \item \textbf{Thermoplastiques :} \textbf{Recyclables mécaniquement par refonte.} On les broie, lave, fond et regranule pour faire de nouveaux objets (bouteilles $\rightarrow$ polaires, bassines $\rightarrow$ tuyaux). C'est la filière active au Cameroun (collecte par les \emph{« ramasseurs »}, usines de granulation à Douala/Yaoundé).
    \item \textbf{Thermodurcissables :} \textbf{NON recyclables par refonte.} Impossible de les refondre. Leur valorisation est difficile : broyage fin comme charge (remplissage) dans d'autres matériaux, ou valorisation énergétique (incinération en cimenterie avec filtres performants), ou mise en décharge contrôlée. \textbf{Séparer les deux familles à la source est crucial.}
\end{itemize}
\end{aretenir}

\begin{savaistu}[Le caoutchouc : un cas particulier]
Le caoutchouc naturel (latex d'hévéa, cultivé au Sud-Cameroun) est un polymère naturel \emph{thermoplastique} à l'état brut : il est collant à chaud, dur et cassant à froid. Pour en faire des pneus ou des semelles, on le \textbf{vulcanise} (chauffage avec du soufre) : cela crée des ponts soufre entre les chaînes $\rightarrow$ il devient \textbf{thermodurcissable} (élastomère réticulé). C'est pour cela qu'un pneu usagé ne fond pas et pose un énorme problème de recyclage (pyrolyse, broyage pour sols sportifs, cimenterie).
\end{savaistu}

\courssub{Exercice résolu : analyse d'objets du quotidien}

\begin{exoresolu}[Classer et justifier]
\textbf{Énoncé :} Tu vides ton sac d'école et tu trouves : une bouteille d'eau vide (1), un stylo à bille cassé (corps transparent) (2), une prise multiple « \emph{multiprise} » (3), le couvercle d'une boîte de beurre de karité (4). Classe ces objets en thermoplastiques ou thermodurcissables. Justifie chaque choix par le comportement attendu à la chaleur et l'intérêt pour l'usage.

\tcblower
\textbf{Solution détaillée :}
\begin{enumerate}
    \item \textbf{Bouteille d'eau (PET) : Thermoplastique.} Elle doit être moulée par soufflage (chauffage $\rightarrow$ forme $\rightarrow$ refroidissement). Si on la chauffe, elle se déforme (ramollit). Recyclable (code 1).
    \item \textbf{Corps de stylo (PS ou SAN) : Thermoplastique.} Moulé par injection. Fond s'il tombe près d'une flamme. Recyclable (code 6 ou 7).
    \item \textbf{Prise multiple (Bakélite / Urée-formol) : Thermodurcissable.} Doit supporter l'échauffement par effet Joule ($P = R I^2$) sans fondre, pour éviter court-circuit et incendie. Ne fond pas, noircit si surchauffe. Non recyclable par refonte.
    \item \textbf{Couvercle beurre karité (PP) : Thermoplastique.} Souvent en polypropylène (PP, code 5) qui supporte \SI{100}{\celsius} (stérilisation) sans fondre (Tf $\approx$ \SI{165}{\celsius}), mais fondra au feu nu. Recyclable.
\end{enumerate}
\textbf{Bilan :} 3 thermoplastiques (recyclables), 1 thermodurcissable (déchet ultime ou valorisation énergétique).
\end{exoresolu}

\begin{exercice}[À toi de jouer !]
\begin{enumerate}
    \item Un électricien doit choisir un matériau pour fabriquer un nouveau modèle de domino (connecteur électrique). Il hésite entre du polyéthylène haute densité (PEHD) et de la résine mélamine-formol. Lequel conseillerais-tu ? Justifie par deux arguments techniques.
    \item Dans un centre de tri à Douala, un opérateur reçoit un sac mélangeant des bouteilles PET, des bouchons PP, des gaines PVC et des morceaux de vieilles prises électriques. Explique pourquoi il est \emph{impératif} de séparer les prises électriques du reste avant d'envoyer le tout en extrusion (refonte).
    \item Dessine schématiquement (schéma type « spaghettis » vs « filet 3D ») la différence de structure microscopique entre un thermoplastique et un thermodurcissable. Légende.
\end{enumerate}
\end{exercice}

\begin{corrige}[Exercice n°17]
\begin{enumerate}
    \item \textbf{Résine mélamine-formol (thermodurcissable).} Arguments : (1) Ne fond pas, résiste à l'échauffement par effet Joule dans le domino (sécurité incendie). (2) Rigidité diélectrique élevée, ne se déforme pas sous le serrage des vis.
    \item Si les prises (thermodurcissables) partent en extrusion avec les bouteilles (thermoplastiques) : elles ne fondront pas $\rightarrow$ bouchent la filière, dégradent la vis de l'extrudeuse, créent des points noirs/défauts dans les granulés recyclés, dégagent des gaz toxiques (phénol, formaldéhyde) dans l'atelier. Le lot de granulés est perdu.
    \item Schéma attendu : Gauche : lignes ondulées parallèles non connectées (flèches de glissement). Droite : grille/réseau avec nœuds de liaison (ponts covalents) bloquant tout mouvement. Légende : « Chaînes libres $\rightarrow$ fusion réversible » vs « Réseau réticulé $\rightarrow$ dégradation irréversible ».
\end{enumerate}
\end{corrige}

\begin{attention}[Pièges BEPC à éviter]
\begin{itemize}
    \item Confondre « thermodurcissable » et « résistant à la chaleur ». Un thermoplastique comme le PP résiste à \SI{100}{\celsius} (eau bouillante) mais fond au feu. Un thermodurcissable ne fond \emph{jamais}.
    \item Dire « le thermodurcissable brûle ». Préciser : « il se \textbf{décompose} (pyrolyse/carbonisation) sans fondre ». La combustion est une réaction avec l'oxygène ; la pyrolyse est une décomposition thermique à l'abri de l'air (ou en début de chauffe).
    \item Oublier que \textbf{tous} les plastiques finissent par se dégrader s'on les chauffe trop (température de dégradation $T_d$). La différence est : \emph{avant} $T_d$, le thermoplastique a \emph{fondue} (état liquide), le thermodurcissable est resté solide.
    \item Croire que « recyclable » = « biodégradable ». Faux. Le PET est recyclable (refonte) mais pas biodégradable (des siècles dans la nature).
\end{itemize}
\end{attention}

\courssub{Tableau de synthèse comparatif}

\begin{center}
\begin{tabularx}{\linewidth}{|>{\bfseries}l|X|X|}\hline
\rowcolor{popblueL}
\textbf{Critère de comparaison} & \textbf{Thermoplastiques} & \textbf{Thermodurcissables} \\ \hline
\textbf{Comportement à la chaleur} & Ramollissement $\rightarrow$ Fusion (état liquide visqueux) $\rightarrow$ Solidification au refroidissement. \textbf{Réversible}. & Pas de ramollissement. Noircissement $\rightarrow$ Décomposition (pyrolyse) $\rightarrow$ Gaz + Charbon. \textbf{Irréversible}. \\ \hline
\textbf{Structure microscopique} & Chaînes linéaires ou ramifiées, \textbf{indépendantes}. Liaisons intermoléculaires faibles (Van der Waals). & \textbf{Réseau 3D réticulé}. Liaisons covalentes (ponts) entre chaînes. Une seule macromolécule géante. \\ \hline
\textbf{Mise en forme industrielle} & Injection, extrusion, soufflage, thermoformage (cycle : chauffer $\rightarrow$ former $\rightarrow$ refroidir). Répétable. & Moulage par compression/transfert/injection \textbf{réactive} (chauffer $\rightarrow$ réaction chimique réticulation $\rightarrow$ démoulage). \textbf{Une seule fois}. \\ \hline
\textbf{Exemples d'objets (Cameroun)} & Sachets, bouteilles, bassines, chaises, tuyaux PVC, gaines, bouchons, jouets, emballages alimentaires. & Prises, interrupteurs, poignées casseroles, boutons four, circuits imprimés, colles époxy, mélamine (vaisselle incassable). \\ \hline
\textbf{Recyclage matière (refonte)} & \textbf{OUI} (broyage $\rightarrow$ lavage $\rightarrow$ extrusion $\rightarrow$ granulés). Filière active. & \textbf{NON} (ne fond pas). Broyage comme charge seulement, ou valorisation énergétique (cimenterie). \\ \hline
\textbf{Risque incendie} & Fondent et gouttent (propagent le feu). Dégagent CO, $\ce{CO2}$, $\ce{H2O}$ (et $\ce{HCl}$ pour PVC). & Ne gouttent pas (auto-extinguibles souvent). Dégagent gaz très toxiques (cyanures, phénol, formaldéhyde). \\ \hline
\end{tabularx}
\end{center}

\begin{popfigure}
\begin{tikzpicture}[scale=0.85, every node/.style={font=\small}]
    % Titre
    \node[font=\bfseries\large, text=popdark] at (0,7.5) {Identification visuelle rapide : les deux familles dans ton quotidien};
    
    % Colonne Thermoplastiques
    \begin{scope}[xshift=-5.5cm]
        \node[align=center, font=\bfseries\large, text=popblue] at (0,6.5) {THERMOPLASTIQUES};
        \node[align=center, draw=popblue, thick, rounded corners, fill=popblue!5, inner sep=4pt] at (0,5.8) {Se ramollissent \& fondent \\ \textbf{Réversible} $\rightarrow$ \textbf{Recyclables}};
        
        % Objets
        \begin{scope}[yshift=4.5]
            % Sachet
            \draw[popblue, thick, fill=popblue!10, rounded corners=2pt] (-1.8,0) rectangle (1.8,1.2);
            \draw[popblue, thick] (-1.8,0.8) -- (1.8,0.8);
            \draw[popblue, thick] (-0.6,0.8) .. controls (-0.6,1.3) and (0.6,1.3) .. (0.6,0.8);
            \node at (0,0.4) {\tiny Sachet PE (500F)};
            \node at (0,-0.3) {\tiny Code \textbf{2} ou \textbf{4}};
            
            % Bouteille
            \draw[popblue, thick, fill=popblue!10, rounded corners=2pt] (-1.8,-1.8) rectangle (1.8,-0.6);
            \draw[popblue, thick] (-0.5,-0.6) -- (-0.5,-0.2) -- (0.5,-0.2) -- (0.5,-0.6);
            \draw[popblue, thick] (-1.2,-0.6) .. controls (-1.2,-0.3) and (1.2,-0.3) .. (1.2,-0.6);
            \node at (0,-1.2) {\tiny Bouteille PET (Eau/Jus)};
            \node at (0,-1.9) {\tiny Code \textbf{1}};
            
            % Bouchon / Basse
            \draw[popblue, thick, fill=popblue!10, rounded corners=2pt] (-1.8,-3.6) rectangle (1.8,-2.4);
            \draw[popblue, thick] (-1.0,-2.4) arc (180:0:1.0);
            \node at (0,-3.0) {\tiny Bouchon / Bassine PP};
            \node at (0,-3.7) {\tiny Code \textbf{5}};
        \end{scope}
    \end{scope}
    
    % Flèche séparation
    \draw[popdark, very thick, <->] (-2.5,2) -- (2.5,2) node[midway, above, font=\bfseries\small] {SÉPARER À LA SOURCE};
    
    % Colonne Thermodurcissables
    \begin{scope}[xshift=5.5cm]
        \node[align=center, font=\bfseries\large, text=poporange] at (0,6.5) {THERMODURCISSABLES};
        \node[align=center, draw=poporange, thick, rounded corners, fill=poporange!5, inner sep=4pt] at (0,5.8) {Ne fondent pas \\ Se dégradent \\ \textbf{Irréversible} $\rightarrow$ \textbf{Déchets ultimes}};
        
        % Objets
        \begin{scope}[yshift=4.5]
            % Prise électrique
            \draw[poporange, thick, fill=poporange!10, rounded corners=2pt] (-1.8,0) rectangle (1.8,1.2);
            \foreach \x in {-1.0, 0, 1.0} {
                \fill[popdark] (\x, 0.3) circle (0.12);
            }
            \node at (0,0.4) {\tiny Prise / Interrupteur};
            \node at (0,-0.3) {\tiny Bakélite (PF)};
            
            % Poignée casserole
            \draw[poporange, thick, fill=poporange!10, rounded corners=2pt] (-1.8,-1.8) rectangle (1.8,-0.6);
            \draw[poporange, thick, fill=popdark] (-1.5,-1.8) rectangle (-1.0,-2.5);
            \draw[poporange, thick, fill=popdark] (1.0,-1.8) rectangle (1.5,-2.5);
            \node at (0,-1.2) {\tiny Poignée casserole / marmite};
            \node at (0,-1.9) {\tiny Phénolique / Mélamine};
            
            % Circuit imprimé
            \draw[poporange, thick, fill=popgreen!10, rounded corners=2pt] (-1.8,-3.6) rectangle (1.8,-2.4);
            \draw[popgreen!80!black, very thick] (-1.5,-3.0) -- (1.5,-3.0);
            \draw[popgreen!80!black, very thick] (-1.0,-2.6) -- (-1.0,-3.4);
            \draw[popgreen!80!black, very thick] (0,-2.6) -- (0,-3.4);
            \draw[popgreen!80!black, very thick] (1.0,-2.6) -- (1.0,-3.4);
            \node at (0,-3.0) {\tiny Circuit imprimé (PCB)};
            \node at (0,-3.7) {\tiny Résine Époxy (FR4)};
        \end{scope}
    \end{scope}
\end{tikzpicture}
\legende{Repérage visuel et codes de recyclage : les thermoplastiques (bleu) portent un code 1 à 6 dans un triangle ; les thermodurcissables (orange) n'en ont généralement pas (ou code 7 « autres ») et ne vont pas dans le bac de recyclage classique.}
\end{popfigure}

Cette classification binaire structure toute la filière des matériaux polymères. Dans la prochaine section, nous étudierons les \textbf{plastiques courants} un par un (PE, PP, PVC, PET, PS, etc.) avec leurs propriétés physiques, leurs usages précis et leurs codes d'identification, avant d'aborder le défi majeur de la \textbf{pollution} et de la \textbf{gestion des déchets} au Cameroun.

\coursec{Les plastiques courants : identification et usages}

Dans la section précédente, nous avons appris à classer les matières plastiques en deux grandes familles : les \cle{thermoplastiques}, qui se ramollissent à la chaleur et peuvent être refondus, et les \cle{thermodurcissables}, qui durcissent définitivement après leur première mise en forme. Mais cette classification ne suffit pas au quotidien. Quand tu tiens une bouteille d'eau minérale, un sachet de « kash », une chaise en plastique ou un tuyau d'arrosage, comment savoir de quel plastique il s'agit exactement ? Et surtout, comment expliquer que le fabricant ait choisi \emph{ce} plastique-là et pas un autre ? C'est toute l'ambition de cette section : apprendre à reconnaître les plastiques courants et à justifier leurs usages par leurs propriétés.

\courssub{Pourquoi existe-t-il plusieurs plastiques ?}

Observe autour de toi, à la maison ou au marché. Le sachet transparent qui enveloppe le pain est souple et léger. La chaise en plastique du salon est rigide et solide. Le tuyau d'arrosage est flexible mais résistant. La prise électrique murale est dure et ne fond pas, même quand elle chauffe un peu. Tous ces objets sont « en plastique », pourtant ils ne se ressemblent pas du tout !

L'explication est simple : il n'existe pas \emph{un} plastique, mais \emph{des} plastiques. Chacun est fabriqué à partir d'un \cle{polymère} différent, c'est-à-dire d'une très longue molécule formée par la répétition d'une même petite unité (le \cle{monomère}). Selon la nature de cette unité et la façon dont les chaînes sont arrangées (linéaires, ramifiées, réticulées), le matériau obtenu sera souple ou rigide, transparent ou opaque, léger ou dense, fusible ou infusible.

\begin{aretenir}[Idée directrice]
Chaque plastique possède un \textbf{nom}, une \textbf{abréviation} (sigle international) et un ensemble de \textbf{propriétés} qui déterminent ses \textbf{usages}. Au BEPC, il faut savoir associer un objet courant à son plastique et justifier ce choix par au moins une propriété.
\end{aretenir}

\courssub{Panorama des plastiques usuels}

\courssub{Le polyéthylène (PE) : le roi des sachets et des bidons}

Le \cle{polyéthylène} (sigle \cle{PE}) est le plastique le plus produit au monde. Il est obtenu par polymérisation de l'éthylène, un gaz issu du pétrole (raffiné à la SONARA). Il existe en deux grandes variétés selon la densité, qui dépend du degré de ramification des chaînes :

\begin{itemize}
\item le \cle{PEBD} (polyéthylène basse densité) : chaînes très ramifiées, \textbf{souple}, \textbf{léger}, \textbf{transparent}. C'est le plastique des \textbf{sachets} d'emballage, des sacs de courses, des films alimentaires (film étirable) ;
\item le \cle{PEHD} (polyéthylène haute densité) : chaînes peu ramifiées, plus \textbf{rigide} et plus \textbf{résistant}. On le trouve dans les \textbf{bidons} d'huile, les jerricans, les bouchons, les caisses, certains tuyaux.
\end{itemize}

Ses propriétés expliquent ses usages : il est \textbf{léger}, \textbf{imperméable à l'eau}, \textbf{insoluble} dans la plupart des liquides courants et \textbf{bon marché}. C'est un thermoplastique : le PEBD fond vers \SI{110}{\celsius}, le PEHD vers \SI{130}{\celsius}, ce qui permet de les recycler facilement par fusion.

\courssub{Le polypropylène (PP) : le plastique des chaises et des marmites}

Le \cle{polypropylène} (sigle \cle{PP}) ressemble au polyéthylène, mais il est plus \textbf{rigide}, plus \textbf{résistant à la chaleur} (température de ramollissement Vicat vers \SI{150}{\celsius}, il supporte l'eau bouillante à \SI{100}{\celsius} sans se déformer) et plus \textbf{dur}. C'est pourquoi on l'utilise pour fabriquer :

\begin{itemize}
\item les \textbf{chaises et tables} en plastique si répandues au Cameroun ;
\item les \textbf{marmites}, bassines, seaux et cuvettes ;
\item les boîtes alimentaires allant au four à micro-ondes ;
\item les cordes et les sacs tissés (sacs de riz, d'engrais, de ciment).
\end{itemize}

\begin{exemplebox}[Pourquoi le PP pour une marmite ?]
Une marmite en plastique doit supporter de l'eau chaude sans se ramollir. Le PE fondrait ou se déformerait vers \SI{110}{\celsius} ; le PP, lui, résiste à la chaleur de l'eau bouillante (\SI{100}{\celsius}) et au-delà. Le choix du plastique se justifie donc par la propriété : \emph{résistance à la chaleur}.
\end{exemplebox}

\courssub{Le polychlorure de vinyle (PVC) : tuyaux et câbles électriques}

Le \cle{polychlorure de vinyle} (sigle \cle{PVC}) contient des atomes de chlore dans ses chaînes, ce qui lui donne des propriétés particulières : il est \textbf{résistant à la corrosion}, \textbf{difficilement inflammable} (auto-extinguible) et \textbf{excellent isolant électrique}. Selon les additifs ajoutés (notamment des \emph{plastifiants} comme les phtalates), il peut être :

\begin{itemize}
\item \textbf{rigide} (PVC non plastifié) : tuyaux de plomberie et d'évacuation des eaux usées, gouttières, cadres de fenêtres ;
\item \textbf{souple} (PVC plastifié) : gaine des \textbf{câbles électriques}, tuyaux d'arrosage, nappes de table, bottes.
\end{itemize}

\begin{attention}[Danger : ne jamais brûler du PVC]
Le PVC brûlé dégage des vapeurs \textbf{chlorées très toxiques} (notamment du chlorure d'hydrogène HCl, qui forme de l'acide chlorhydrique au contact de l'humidité des poumons), dangereuses pour la santé et pour l'environnement. Il ne faut jamais jeter de gaines de câbles ou de tuyaux PVC dans un feu de déchets. C'est une consigne de sécurité à retenir absolument.
\end{attention}

\courssub{Le polyéthylène téréphtalate (PET) : le plastique des bouteilles}

Le \cle{PET} (polyéthylène téréphtalate) est le plastique des \textbf{bouteilles} d'eau minérale et de boissons gazeuses. Pourquoi lui ? Parce qu'il est à la fois :

\begin{itemize}
\item \textbf{transparent} comme le verre, ce qui permet de voir le contenu ;
\item \textbf{très léger}, ce qui facilite le transport (une bouteille PET vide de \SI{1,5}{L} pèse environ \SI{30}{g}, contre plus de \SI{400}{g} pour le verre) ;
\item \textbf{incassable}, contrairement au verre ;
\item \textbf{bon barrage aux gaz}, ce qui conserve le gaz carbonique des boissons gazeuses et empêche l'oxydation de l'eau ;
\item \textbf{non toxique} au contact des aliments (apte au contact alimentaire).
\end{itemize}

\begin{savaistu}[Du plastique aux vêtements]
Le PET des bouteilles recyclées peut être fondu et transformé en \textbf{fibres textiles} (polyester) ! Avec les bouteilles collectées, on fabrique des vêtements polaires, des tapis, des sacs, et même le rembourrage des couettes. Il faut environ 25 bouteilles en PET de \SI{1,5}{L} pour fabriquer un pull polaire. Recycler, c'est donner une seconde vie à la matière.
\end{savaistu}

\courssub{Le polystyrène (PS) : gobelets et emballages}

Le \cle{polystyrène} (sigle \cle{PS}) se présente sous deux formes :

\begin{itemize}
\item le PS \textbf{rigide et transparent} (PS cristal) : gobelets jetables, couverts en plastique, boîtiers de CD, pots de yaourt ;
\item le PS \textbf{expansé} (PSE), blanc et plein de bulles d'air (98\% air) : barquettes alimentaires, emballages protecteurs des appareils électroménagers, glacières. L'air emprisonné en fait un excellent \textbf{isolant thermique}.
\end{itemize}

Il est très léger et bon marché, mais \textbf{cassant} (peu résistant aux chocs pour le PS cristal) et difficile à recycler (faible densité, volume important pour le PSE). Les barquettes en polystyrène jetées dans la nature sont l'une des sources majeures de pollution plastique dans nos villes (caniveaux bouchés, paysage dégradé).

\courssub{La bakélite : la pionnière thermodurcissable}

La \cle{bakélite} (phénol-formaldéhyde) est historiquement le \textbf{premier plastique entièrement synthétique}, inventé en 1907 par le chimiste Léo Baekeland. C'est un \textbf{thermodurcissable} : une fois cuite et mise en forme (par condensation irréversible), elle ne fond plus jamais, même chauffée fortement. Elle se dégrade (carbonise) avant de fondre. Elle est dure, brune ou noire, et surtout \textbf{excellente isolante électrique} et \textbf{résistante à la chaleur}.

C'est pourquoi on la retrouve dans l'\textbf{appareillage électrique} : prises de courant, interrupteurs, douilles de lampes, fiches, poignées de casseroles et de fers à repasser. Imagine une prise en polyéthylène : elle fondrait à la première surchauffe !

\begin{exemplebox}[Pourquoi la bakélite pour une prise électrique ?]
Une prise doit (1) empêcher le courant de passer vers l'utilisateur : la bakélite est \emph{isolante électrique} ; (2) supporter l'échauffement sans fondre : la bakélite est \emph{thermodurcissable}, donc infusible. Deux propriétés justifient un usage.
\end{exemplebox}

\courssub{Le code d'identification des plastiques}

\courssub{Un problème pratique : comment trier ce qu'on ne reconnaît pas ?}

Pour recycler les plastiques, il faut les \textbf{trier par famille} : on ne peut pas fondre ensemble du PET et du PVC, le mélange donnerait un matériau de mauvaise qualité (incompatibilité chimique, dégradation du PVC). Mais à l'œil nu, un bidon en PEHD et un bidon en PP se ressemblent beaucoup ! Pour résoudre ce problème, les industriels ont adopté en 1988 un système international de marquage (norme ISO 11469 / SPI).

\begin{definition}[Code d'identification des plastiques]
Le \cle{code d'identification des plastiques} est un symbole formé d'un \textbf{triangle de trois flèches} (symbole du recyclage, dit « boucle de Möbius ») entourant un \textbf{chiffre de 1 à 7}, souvent accompagné d'un sigle. Ce chiffre indique la \textbf{nature chimique} du plastique afin de faciliter le \textbf{tri} et le \textbf{recyclage}.
\end{definition}

\begin{aretenir}[Les sept codes à connaître]
\begin{itemize}
\item \textbf{1 = PET} : bouteilles d'eau et de boissons ;
\item \textbf{2 = PEHD} : bidons, jerricans, flacons de produits ménagers ;
\item \textbf{3 = PVC} : tuyaux, gaines de câbles, nappes ;
\item \textbf{4 = PEBD} : sachets, sacs, films d'emballage ;
\item \textbf{5 = PP} : chaises, marmites, boîtes alimentaires ;
\item \textbf{6 = PS} : gobelets, barquettes, couverts jetables ;
\item \textbf{7 = autres} : tous les autres plastiques (polycarbonate, bakélite, mélanges, bioplastiques\ldots).
\end{itemize}
\end{aretenir}

\begin{popfigure}
\begin{center}
\begin{tikzpicture}[scale=1.1]
% Triangle de recyclage (trois flèches stylisées)
\draw[line width=3pt, popgreen, rounded corners=4pt, -{Triangle[length=3.5mm,width=2.5mm]}]
  (0,0) -- (2.5,0);
\draw[line width=3pt, popgreen, rounded corners=4pt, -{Triangle[length=3.5mm,width=2.5mm]}]
  (2.8,0.3) -- (1.4,2.7);
\draw[line width=3pt, popgreen, rounded corners=4pt, -{Triangle[length=3.5mm,width=2.5mm]}]
  (1.1,2.7) -- (-0.3,0.3);
\node[font=\Large\bfseries, popdark] at (1.25,1.1) {1 à 7};
\node[font=\small\itshape, popmuted] at (1.25,-0.5) {symbole du code};
% Tableau des codes à droite
\node[anchor=west, font=\small] at (3.8,2.5) {\textbf{1} PET -- bouteilles};
\node[anchor=west, font=\small] at (3.8,1.9) {\textbf{2} PEHD -- bidons, jerricans};
\node[anchor=west, font=\small] at (3.8,1.3) {\textbf{3} PVC -- tuyaux, câbles};
\node[anchor=west, font=\small] at (3.8,0.7) {\textbf{4} PEBD -- sachets, films};
\node[anchor=west, font=\small] at (3.8,0.1) {\textbf{5} PP -- chaises, marmites};
\node[anchor=west, font=\small] at (3.8,-0.5) {\textbf{6} PS -- gobelets, barquettes};
\node[anchor=west, font=\small] at (3.8,-1.1) {\textbf{7} autres plastiques};
\draw[popline, thin] (3.4,2.8) -- (3.4,-1.4);
\end{tikzpicture}
\end{center}
\legende{Le code d'identification des plastiques : un triangle de trois flèches entourant un chiffre de 1 à 7, moulé ou imprimé sur l'objet, qui indique la famille du plastique.}
\end{popfigure}

\begin{methode}[Identifier le plastique d'un emballage]
\begin{enumerate}
\item \textbf{Retourner} l'objet : le symbole se trouve presque toujours \emph{sous} la bouteille, le bidon ou le gobelet, ou sur une face peu visible.
\item \textbf{Chercher le triangle} de trois flèches, moulé en relief ou imprimé.
\item \textbf{Lire le chiffre} au centre (de 1 à 7) et, si présent, le sigle gravé en dessous (PET, PE-HD, PP\ldots).
\item \textbf{Consulter le tableau} des codes pour nommer le plastique.
\item \textbf{Conclure} en reliant le plastique identifié aux propriétés utiles de l'objet.
\end{enumerate}
\end{methode}

\begin{exemplebox}[Deux identifications types]
\begin{itemize}
\item Une \textbf{bouteille d'eau minérale} porte sous son fond le code \textbf{1} : c'est du \textbf{PET}, choisi pour sa transparence et sa légèreté.
\item Un \textbf{tuyau d'arrosage} porte le code \textbf{3} : c'est du \textbf{PVC} souple, choisi pour sa flexibilité et sa résistance à la corrosion par l'eau.
\end{itemize}
\end{exemplebox}

\courssub{Activité d'observation : chasse aux codes en classe}

\begin{experience}[Relever les codes d'identification sur des emballages]
\textbf{Matériel} : emballages plastiques propres apportés par les élèves (bouteille d'eau, bidon d'huile, sachet, gobelet, flacon de savon liquide, chaise ou bassine\ldots), loupe éventuelle, tableau à compléter.

\textbf{Protocole} :
\begin{enumerate}
\item Former des groupes de trois ou quatre élèves et répartir les objets.
\item Pour chaque objet, chercher le triangle de flèches (souvent sous le fond) et noter le chiffre et le sigle.
\item Compléter le tableau : objet, code observé, plastique correspondant, propriété utile.
\item Comparer les résultats entre groupes et discuter les objets sans symbole visible.
\end{enumerate}

\textbf{Observations attendues} : la bouteille d'eau porte le code 1 (PET), le bidon d'huile le code 2 (PEHD), le flacon souple le code 4 (PEBD) ou 5 (PP), le gobelet le code 6 (PS). Certains objets anciens ou de fabrication artisanale ne portent aucun code.

\textbf{Interprétation} : le marquage est systématique sur les emballages industriels car il est indispensable au tri en usine de recyclage. L'absence de code signale souvent un objet thermodurcissable ou un mélange (code 7), difficilement recyclable.

\textbf{Conclusion} : grâce au code d'identification, n'importe qui peut reconnaître la famille d'un plastique et le diriger vers la bonne filière de tri.
\end{experience}

\courssub{Tableau de synthèse : plastique, propriété, objet courant}

\begin{center}
\begin{tabularx}{\linewidth}{|l|l|X|X|}
\hline
\rowcolor{popblueL} \textbf{Code} & \textbf{Plastique} & \textbf{Propriétés utiles} & \textbf{Objets courants} \\
\hline
1 & PET & transparent, léger, barrage aux gaz & bouteilles d'eau et de boissons gazeuses \\
\hline
2 & PEHD & rigide, résistant, imperméable & bidons d'huile, jerricans, flacons \\
\hline
3 & PVC & isolant électrique, résistant à la corrosion & tuyaux, gaines de câbles, nappes \\
\hline
4 & PEBD & souple, léger, bon marché & sachets, sacs, films d'emballage \\
\hline
5 & PP & rigide, résistant à la chaleur & chaises, marmites, bassines, boîtes \\
\hline
6 & PS & léger, isolant thermique (expansé) & gobelets, barquettes, couverts jetables \\
\hline
7 & autres (bakélite\ldots) & infusible, isolant électrique & prises, interrupteurs, poignées \\
\hline
\end{tabularx}
\end{center}

\begin{attention}[Pièges fréquents au BEPC]
\begin{itemize}
\item Ne confonds pas \textbf{PEHD} (code 2, rigide) et \textbf{PEBD} (code 4, souple) : ce sont deux variétés du polyéthylène, mais leurs usages diffèrent.
\item Le chiffre dans le triangle n'est \textbf{pas} un numéro de recyclage garanti : il indique seulement la \emph{nature} du plastique. Le code 7, par exemple, regroupe des plastiques rarement recyclés.
\item Justifier un usage demande de citer une \textbf{propriété précise} (« le PVC est isolant électrique »), et non une phrase vague (« parce que c'est du plastique »).
\item Ne pas confondre \textbf{fusible} (thermoplastique) et \textbf{infusible} (thermodurcissable).
\end{itemize}
\end{attention}

\begin{exoresolu}[Associer un objet à son plastique et justifier]
Énoncé. Voici quatre objets : (a) une bouteille de boisson gazeuse, (b) la gaine d'un câble électrique, (c) une chaise en plastique, (d) une prise de courant. Pour chacun, indique le plastique utilisé, son code d'identification, et justifie le choix par une propriété.
\tcblower
\textbf{Solution détaillée.}

(a) \textbf{Bouteille de boisson gazeuse} : plastique \textbf{PET}, code \textbf{1}. Justification : le PET est transparent (on voit la boisson), léger (transport facile) et bon barrage aux gaz (le gaz carbonique ne s'échappe pas).

(b) \textbf{Gaine de câble électrique} : plastique \textbf{PVC} souple, code \textbf{3}. Justification : le PVC est un excellent isolant électrique, il protège l'utilisateur du courant, et il résiste à la corrosion et à l'humidité.

(c) \textbf{Chaise en plastique} : plastique \textbf{PP}, code \textbf{5}. Justification : le polypropylène est rigide et résistant, il supporte le poids d'une personne sans se casser ni se déformer.

(d) \textbf{Prise de courant} : plastique \textbf{bakélite}, code \textbf{7}. Justification : la bakélite est thermodurcissable (elle ne fond pas en cas d'échauffement) et isolante électrique, ce qui garantit la sécurité.

\textbf{Méthode à retenir} : pour chaque objet, pose-toi deux questions : \emph{que doit faire l'objet ?} et \emph{quelle propriété du plastique le permet ?} La réponse relie alors naturellement l'objet à son plastique.
\end{exoresolu}

\begin{exercice}[À toi de jouer]
\begin{enumerate}
\item Tu retournes un jerrican d'huile végétale et tu lis le chiffre 2 dans un triangle. Quel est le plastique ? Cite deux de ses propriétés utiles ici.
\item Pourquoi ne fabrique-t-on pas les gobelets jetables en bakélite ? Donne au moins deux raisons.
\item Un élève affirme : « Tous les plastiques portant un triangle peuvent être recyclés ensemble dans la même usine. » Corrige cette affirmation.
\end{enumerate}
\end{exercice}

\begin{corrige}[Exercice « À toi de jouer »]
\begin{enumerate}
\item Le chiffre 2 correspond au \textbf{PEHD} (polyéthylène haute densité). Propriétés utiles : il est \textbf{rigide} (le jerrican garde sa forme) et \textbf{imperméable} (l'huile ne traverse pas la paroi) ; il est aussi résistant aux chocs.
\item La bakélite est \textbf{chère} et \textbf{dure à mouler} en objets très fins ; elle est aussi \textbf{lourde} comparée au polystyrène. Pour un gobelet jetable, on veut un plastique très léger et très bon marché : c'est le PS (code 6).
\item Faux : le triangle indique seulement la \textbf{nature} du plastique. Les différentes familles (PET, PVC, PP\ldots) doivent être \textbf{triées séparément} avant recyclage, car elles fondent à des températures différentes et ne se mélangent pas correctement (incompatibilité).
\end{enumerate}
\end{corrige}

\begin{aretenir}[L'essentiel de la section]
\begin{itemize}
\item Les plastiques courants sont : \textbf{PE} (sachets, bidons), \textbf{PP} (chaises, marmites), \textbf{PVC} (tuyaux, câbles), \textbf{PET} (bouteilles), \textbf{PS} (gobelets, emballages) et la \textbf{bakélite} (appareillage électrique).
\item Chaque usage se justifie par une \textbf{propriété} : légèreté, imperméabilité, isolation électrique, résistance à la chaleur ou à la corrosion.
\item Le \textbf{code d'identification} est un triangle de flèches avec un chiffre de 1 à 7 : 1 = PET, 2 = PEHD, 3 = PVC, 4 = PEBD, 5 = PP, 6 = PS, 7 = autres. Pour lire le code, on retourne l'objet.
\item Le PVC ne doit \textbf{jamais être brûlé} : ses vapeurs chlorées sont très toxiques.
\end{itemize}
\end{aretenir}

Maintenant que tu sais reconnaître chaque plastique et expliquer pourquoi il a été choisi, une question s'impose : ces matériaux si pratiques sont-ils sans inconvénients ? C'est ce que nous examinerons dans la section suivante, consacrée aux avantages et aux limites des plastiques.

\coursec{Avantages et limites des plastiques}

Dans la section précédente, nous avons appris à identifier les plastiques courants (PET, PE, PVC, PP, PS) et à associer chacun à ses usages : bouteilles d'eau, sachets, tuyaux, câbles électriques, chaises en plastique. Une question se pose alors naturellement : si les plastiques sont partout autour de nous, c'est bien qu'ils possèdent des qualités remarquables. Mais alors, pourquoi entend-on aussi souvent dire que le plastique est un problème ? Pour répondre honnêtement, il faut peser le pour et le contre : c'est l'objet de cette section. Tu vas découvrir que le plastique est à la fois un matériau extraordinaire et un déchet redoutable, et que tout est une question de gestion.

\courssub{Les avantages des matières plastiques}

\textbf{Pourquoi le plastique a-t-il conquis le monde ?}

Observe autour de toi : la chaise sur laquelle tu es assis, le stylo avec lequel tu écris, le seau de la maison, le téléphone, les câbles électriques, la bassine, les sandales... Il y a soixante ans, presque tous ces objets étaient en bois, en métal, en verre ou en céramique. Aujourd'hui, ils sont en plastique. Ce n'est pas un hasard : le plastique possède une combinaison de propriétés qu'aucun autre matériau ne réunit à lui seul.

\begin{definition}[Propriétés avantageuses des plastiques]
Les matières plastiques présentent généralement les avantages suivants :
\begin{itemize}
\item la \cle{légèreté} : leur \cle{masse volumique} est faible (comprise entre \SI{0,9}{g/cm^3} et \SI{1,4}{g/cm^3}), bien inférieure à celle des métaux ou du verre ;
\item le \cle{faible coût} : leur fabrication en grande série est bon marché ;
\item la \cle{facilité de mise en forme} : chauffés, ils se moulent en objets de toutes les formes ;
\item la \cle{résistance à l'eau} et à la plupart des produits chimiques : ils ne rouillent pas, ne pourrissent pas ;
\item les \cle{propriétés isolantes} : ils ne conduisent ni l'électricité ni (bien) la chaleur.
\end{itemize}
\end{definition}

\begin{experience}[Le plastique est-il vraiment léger ?]
\textbf{Matériel :} une bouteille en plastique vide de \SI{1,5}{L} (type PET), une bouteille en verre vide de contenance voisine, une balance de ménage, une bassine d'eau.

\textbf{Protocole :}
\begin{enumerate}
\item Peser séparément la bouteille en plastique et la bouteille en verre.
\item Déposer la bouteille en plastique vide et \emph{bouchée} à la surface de l'eau.
\item (Facultatif) Déposer un bouchon en PE (polyéthylène) et un fragment de PVC dans l'eau.
\end{enumerate}

\textbf{Observations :} la bouteille en plastique pèse environ \SI{30}{g} à \SI{40}{g}, la bouteille en verre environ \SI{400}{g} à \SI{600}{g}, soit plus de dix fois plus. La bouteille en plastique bouchée flotte. Le bouchon en PE flotte, le fragment de PVC coule.

\textbf{Interprétation :} La masse volumique du PET (\approx \SI{1,38}{g/cm^3}) est supérieure à celle de l'eau (\SI{1,00}{g/cm^3}) : le matériau seul coule. Mais la bouteille \emph{vide et bouchée} contient de l'air : sa masse volumique \emph{apparente} devient inférieure à celle de l'eau, donc elle flotte. Le PE (\approx \SI{0,95}{g/cm^3}) a une masse volumique inférieure à l'eau : il flotte toujours. Le PVC (\approx \SI{1,4}{g/cm^3}) coule.

\textbf{Conclusion :} À objet identique, le plastique est beaucoup plus léger que le verre ou le métal. Attention : \cle{tous les plastiques ne flottent pas} (PE, PP flottent ; PET, PVC, PS coulent), mais leur masse volumique reste 3 à 8 fois plus faible que celle de l'acier.
\end{experience}

\textbf{Des avantages qui changent la vie quotidienne.}

Chaque avantage a des conséquences concrètes et positives :

\begin{itemize}
\item \textbf{Conservation des aliments.} Grâce à leur étanchéité à l'eau et à l'air, les emballages plastiques protègent les aliments de l'humidité, des microbes et des insectes. Le riz, l'huile, l'eau minérale ou les arachides se conservent plus longtemps. Cela réduit le gaspillage alimentaire.
\item \textbf{Allègement des véhicules.} Dans une moto ou une voiture, de nombreuses pièces (réservoir, pare-chocs, tableau de bord, rétroviseurs) sont en plastique. Un véhicule plus léger consomme moins d'essence : c'est une économie pour le portefeuille et moins de pollution de l'air.
\item \textbf{Sécurité électrique.} Le plastique est un excellent isolant électrique. Les fils électriques de la maison sont gainés de PVC, les prises et les interrupteurs sont en plastique : c'est ce qui nous protège des électrocutions. Sans plastique, l'électricité domestique serait bien plus dangereuse.
\item \textbf{Hygiène et médecine.} Seringues à usage unique, poches de perfusion, gants : le plastique bon marché a rendu possible le matériel médical jetable, qui limite les contaminations.
\item \textbf{Accessibilité.} Le faible coût du plastique rend des objets utiles (seaux, chaises, chaussures, cahiers protégés) accessibles aux familles modestes.
\end{itemize}

\begin{exemplebox}[Le transport de l'eau]
Pour transporter 1000 bouteilles d'eau de \SI{1,5}{L}, un camion transporte environ \SI{40}{kg} d'emballage si les bouteilles sont en plastique, contre plus de \SI{500}{kg} si elles sont en verre. Le plastique permet donc de transporter plus d'eau avec moins de carburant : c'est l'avantage de la légèreté.
\end{exemplebox}

\courssub{Les limites des matières plastiques}

\textbf{Première limite : une ressource d'origine fossile.}

Nous l'avons vu au début de la leçon : les plastiques sont fabriqués à partir du \cle{pétrole}, une ressource fossile formée en des millions d'années. Or les réserves de pétrole ne sont pas inépuisables.

\begin{aretenir}[Une ressource non renouvelable]
Le pétrole, matière première des plastiques, est une ressource \cle{non renouvelable} : la nature en fabrique infiniment plus lentement que l'homme ne le consomme. Chaque objet plastique jeté après quelques minutes d'usage, c'est un peu de pétrole gaspillé.
\end{aretenir}

\textbf{Deuxième limite : les plastiques ne sont pas biodégradables.}

Voici le cœur du problème. Une peau de banane jetée dans la cour disparaît en quelques mois : des micro-organismes (bactéries, champignons) la digèrent. On dit qu'elle est \cle{biodégradable}. Un sachet plastique, lui, reste là, année après année.

\begin{propriete}[Non biodégradabilité des plastiques (admise)]
La plupart des plastiques ne sont \cle{pas biodégradables} : aucun micro-organisme de la nature n'est capable de dégrader leurs très longues chaînes carbonées (macromolécules), qui n'existaient pas dans la nature avant l'invention du plastique. Un plastique abandonné dans l'environnement persiste donc pendant des dizaines, voire des centaines d'années.
\end{propriete}

\begin{exemplebox}[Une bouteille qui survit à plusieurs générations]
Une bouteille en PET jetée dans la nature met environ \cle{450 ans} à disparaître. Une bouteille jetée aujourd'hui dans un ravin de Yaoundé existerait encore, sous forme de fragments, vers l'an 2475 ! Aucun de tes arrière-arrière-petits-enfants ne la verrait disparaître.
\end{exemplebox}

Comparons maintenant les durées de décomposition de différents déchets courants, à partir de données documentaires.

\begin{popfigure}
\begin{center}
\begin{tikzpicture}
\begin{axis}[
    width=0.95\linewidth,
    height=7cm,
    ybar,
    bar width=14pt,
    ymode=log,
    log origin=infty,
    ylabel={Durée de décomposition (années, échelle logarithmique)},
    symbolic x coords={Peau de banane, Papier, Sac plastique, Bouteille PET, Bouteille en verre},
    xtick=data,
    x tick label style={rotate=20, anchor=east, font=\small},
    ymin=0.05, ymax=2000000,
    ytick={0.1,1,10,100,1000,1000000},
    yticklabels={{0,1 an},{1 an},{10 ans},{100 ans},{1000 ans},{1 million d'années}},
    nodes near coords,
    nodes near coords style={font=\scriptsize, rotate=90, anchor=west},
    grid=major,
    major grid style={popline!50},
    axis line style={popdark},
]
\addplot[fill=popgreen, draw=popdark] coordinates {(Peau de banane,0.25)};
\addplot[fill=popgold, draw=popdark] coordinates {(Papier,0.5)};
\addplot[fill=poporange, draw=popdark] coordinates {(Sac plastique,200)};
\addplot[fill=popblue, draw=popdark] coordinates {(Bouteille PET,450)};
\addplot[fill=poppurple, draw=popdark] coordinates {(Bouteille en verre,1000000)};
\end{axis}
\end{tikzpicture}
\end{center}
\legende{Durées approximatives de décomposition de divers déchets dans la nature (échelle logarithmique) : peau de banane $\approx$ 3 mois (0,25 an) ; papier $\approx$ 6 mois (0,5 an) ; sac plastique $\approx$ 100 à 450 ans ; bouteille PET $\approx$ 450 ans ; bouteille en verre $\approx$ 1 million d'années. L'écart entre les déchets organiques et les plastiques est immense.}
\end{popfigure}

\begin{attention}[Enterré ne veut pas dire disparu !]
Beaucoup croient qu'en enterrant un plastique dans le sol, on s'en débarrasse. C'est faux : \cle{enterrer un plastique ne le fait pas disparaître}. Sous l'effet du soleil (UV), de la chaleur et des frottements mécaniques, il se \cle{fragmente} lentement en morceaux de plus en plus petits appelés \cle{microplastiques} (moins de \SI{5}{mm}). Invisibles à l'œil nu, ces microplastiques contaminent les sols, les rivières, puis l'océan, et finissent dans la chaîne alimentaire.
\end{attention}

\textbf{Troisième limite : la combustion à l'air libre dégage des substances toxiques.}

Dans beaucoup de quartiers, on brûle les ordures, y compris les plastiques, à l'air libre. Cette pratique est dangereuse.

\begin{propriete}[Toxicité des fumées de plastique (admise)]
La combustion des plastiques à l'air libre, à basse température et en présence de peu d'oxygène, est incomplète : elle dégage des fumées contenant des \cle{substances toxiques}, notamment des \cle{dioxines} (surtout avec le PVC), des gaz irritants (chlorure d'hydrogène, formaldéhyde) et des suies. Ce savoir est admis ; il est justifié par des observations sanitaires : les populations exposées régulièrement à ces fumées souffrent davantage de maladies respiratoires, et les dioxines sont reconnues cancérigènes.
\end{propriete}

\begin{attention}[Ne jamais brûler les plastiques]
Il est interdit et dangereux de brûler des sachets, bouteilles ou vieux câbles électriques dans la cour ou au bord de la route. Les fumées noires et âcres sont toxiques pour celui qui les respire, pour les voisins et pour les cultures voisines. La bonne solution est le tri et le recyclage, que nous étudierons dans la dernière section.
\end{attention}

\begin{savaistu}[Des microplastiques jusque dans nos assiettes]
Des microplastiques ont été retrouvés dans les poissons pêchés sur le littoral camerounais (Kribi, Limbé, Douala) ainsi que dans le sel de mer récolté dans plusieurs régions du monde. Le plastique jeté dans la rue finit en fragments dans l'océan, est avalé par les poissons... et revient dans notre assiette. Jeter un sachet par terre, c'est un peu le jeter dans sa propre marmite, avec quelques années de retard.
\end{savaistu}

\courssub{Bilan : un bon matériau, une mauvaise gestion}

Faisons le point. D'un côté, le plastique est léger, bon marché, facile à façonner, étanche, isolant : il conserve nos aliments, allège nos véhicules et nous protège de l'électricité. De l'autre, il est fabriqué à partir d'une ressource non renouvelable, il persiste des siècles dans la nature, se fragmente en microplastiques et dégage des substances toxiques quand on le brûle.

\begin{aretenir}[Le vrai coupable, c'est la mauvaise gestion]
Le problème ne vient pas du plastique lui-même, mais de sa \cle{mauvaise gestion après usage} : jeter n'importe où, enterrer, brûler à l'air libre. Un plastique utilisé longtemps, réutilisé puis recyclé est un matériau utile ; un plastique jeté dans la nature après cinq minutes d'usage est un poison pour l'environnement. La solution passe par les 3R (Réduire, Réutiliser, Recycler), que nous détaillerons dans la dernière section.
\end{aretenir}

\begin{exoresolu}[Qui pollue le plus ?]
\textbf{Énoncé.} Un élève affirme : « Le verre est plus écologique que le plastique, donc je peux jeter ma bouteille en verre dans la brousse sans problème. »
\begin{enumerate}
\item À l'aide du diagramme des durées de décomposition, comparer le sort d'une bouteille en verre et d'une bouteille en PET abandonnées dans la nature.
\item L'affirmation de l'élève est-elle correcte ? Justifier.
\end{enumerate}
\tcblower
\textbf{Solution détaillée.}
\begin{enumerate}
\item D'après le diagramme, une bouteille en PET met environ 450 ans à se décomposer, tandis qu'une bouteille en verre peut persister environ un million d'années dans la nature : le verre persiste donc bien plus longtemps que le plastique.
\item L'affirmation est \textbf{incorrecte}. Certes, le verre est un matériau naturel (à base de sable) et non toxique, mais il ne se dégrade pratiquement jamais et ses éclats blessent les hommes et les animaux. La règle est la même pour tous les déchets : aucun déchet, plastique ou non, ne doit être abandonné dans la nature. Il faut le déposer dans une poubelle, puis le trier pour le réutiliser ou le recycler.
\end{enumerate}
\end{exoresolu}

\begin{exercice}[À toi de jouer]
\begin{enumerate}
\item Citer trois avantages des plastiques et, pour chacun, donner un exemple d'application dans la vie courante au Cameroun.
\item Expliquer pourquoi les micro-organismes ne parviennent pas à dégrader les plastiques.
\item Ta petite sœur veut brûler un tas de sachets plastiques derrière la maison. Rédige deux phrases pour l'en dissuader, en citant le danger précis.
\end{enumerate}
\end{exercice}

\begin{corrige}[Exercice « À toi de jouer »]
\begin{enumerate}
\item Exemples de réponses : légèreté (seaux et bassines faciles à porter, pièces de moto allégées) ; faible coût (chaises, sandales et emballages accessibles à tous) ; propriétés isolantes (gaine des câbles électriques, prises et interrupteurs qui protègent de l'électrocution) ; résistance à l'eau (conservation du riz et de l'huile, bouteilles d'eau étanches).
\item Les plastiques sont formés de très longues chaînes carbonées (macromolécules) qui n'existaient pas dans la nature : aucun micro-organisme ne possède les enzymes capables de couper et de digérer ces chaînes. Les plastiques sont donc non biodégradables et persistent des siècles.
\item Exemple : « Ne brûle pas ces sachets : leur combustion à l'air libre dégage des fumées toxiques contenant des dioxines, qui provoquent des maladies respiratoires et sont cancérigènes. Mets-les plutôt dans la poubelle pour qu'ils soient triés et recyclés. »
\end{enumerate}
\end{corrige}

\coursec{Pollution plastique au Cameroun et dans le monde}

Dans la section précédente, nous avons vu que les plastiques possèdent des qualités remarquables : ils sont légers, solides, bon marché et... presque indestructibles. C'est précisément cette dernière qualité qui devient un terrible défaut lorsque le plastique devient un déchet. Un sachet utilisé dix minutes pour transporter des beignets peut mettre plusieurs centaines d'années à disparaître de la nature. Où vont donc tous ces sachets, ces bouteilles et ces emballages que nous jetons chaque jour ? C'est la question que nous allons explorer dans cette section, en partant de ce que tu observes toi-même dans ta rue, ton quartier ou ton village.

\courssub{Des constats qui sautent aux yeux}

\courssub{Les sachets envahissent nos villes et nos villages}

Fais une petite expérience de la vie courante : sors de chez toi et compte les déchets plastiques visibles sur cent mètres de rue. Sachets noirs ou bleus, bouteilles vides, emballages de biscuits, restes de chaises en plastique cassées... Dans la plupart de nos quartiers, le résultat est impressionnant.

\begin{experience}[Observer la pollution plastique autour de soi]
\textbf{Matériel :} un cahier, un crayon, des gants (pour ne rien toucher directement).

\textbf{Protocole :} parcours un trajet de \SI{100}{m} dans ton quartier et note le nombre et la nature des déchets plastiques visibles (sachets, bouteilles, emballages). Observe aussi l'état du caniveau le plus proche.

\textbf{Observations typiques :} des dizaines de sachets et de bouteilles ; le caniveau est souvent partiellement bouché par un amas de sachets mêlés à de la boue ; de l'eau stagnante s'y accumule.

\textbf{Interprétation :} les plastiques jetés au sol ne disparaissent pas. Ils s'accumulent, obstruent les voies d'évacuation des eaux de pluie et se dégradent très lentement.

\textbf{Conclusion :} le rejet sauvage des déchets plastiques est la première étape d'une pollution qui touche ensuite toute la ville, les rivières et même la mer.
\end{experience}

\begin{exemplebox}[Douala et Yaoundé sous les eaux]
À Douala, la capitale économique, chaque saison des pluies transforme certains quartiers (Deïdo, New Bell, Bonabéri) en véritables marécages. Les enquêtes des services d'hygiène urbaine le confirment : les caniveaux et les buses d'évacuation sont colmatés par des tonnes de sachets plastiques mélangés aux boues. L'eau de pluie ne pouvant plus s'écouler, elle déborde dans les rues et les maisons. À Yaoundé, les quartiers comme Mokolo ou Briqueterie connaissent le même phénomène. Le plastique n'est pas la seule cause des inondations, mais il les aggrave considérablement.
\end{exemplebox}

\courssub{Des rivières et des plages souillées}

Le voyage du plastique ne s'arrête pas au caniveau. Les eaux de pluie entraînent les déchets vers les rivières : la Wouri à Douala, la Sanaga, le Mfoundi à Yaoundé. Ces rivières les transportent jusqu'à l'océan Atlantique. Résultat : les magnifiques plages de Kribi et de Limbé, pourtant destinations touristiques de premier plan, reçoivent à chaque marée des bouteilles, des sachets et des emballages rejetés par les vagues. Les pêcheurs de Limbé remontent parfois plus de plastique que de poissons dans leurs filets.

\begin{attention}[Ce n'est pas qu'une question de saleté]
On entend souvent : « Ce n'est que du plastique, ce n'est pas dangereux comme un produit chimique. » Erreur ! La pollution plastique est une pollution durable et cumulative : chaque sachet jeté s'ajoute à ceux déjà présents, car le plastique ne se biodégrade pratiquement pas. Un sachet jeté aujourd'hui polluera encore l'environnement lorsque tes arrière-petits-enfants seront nés.
\end{attention}

\courssub{Le mécanisme de la pollution : du sachet au microplastique}

\courssub{Une chaîne de pollution bien documentée}

Comment un sachet jeté dans la rue finit-il dans l'assiette ? La démarche scientifique permet de reconstituer ce parcours, étape par étape, grâce à des observations documentées par les chercheurs du monde entier :

\begin{enumerate}
\item \textbf{Rejet sauvage :} le sachet est jeté au sol, dans la rue ou dans un caniveau, au lieu d'être déposé dans une poubelle ou collecté.
\item \textbf{Transport par l'eau :} les pluies l'entraînent vers les rivières, puis vers l'océan.
\item \textbf{Fragmentation :} sous l'action du soleil (rayons ultraviolets), de la chaleur, des vagues et du frottement, le plastique se fragilise et se casse en morceaux de plus en plus petits. Attention : il ne disparaît pas, il se fragmente !
\item \textbf{Formation de microplastiques :} les fragments deviennent si petits qu'on ne les voit plus à l'œil nu.
\item \textbf{Ingestion par les animaux :} poissons, crevettes, zooplancton confondent ces particules avec de la nourriture et les avalent.
\item \textbf{Remontée dans la chaîne alimentaire :} le petit poisson est mangé par un plus gros, qui est pêché... et consommé par l'homme. Le plastique parcourt ainsi toute la chaîne alimentaire.
\end{enumerate}

\begin{definition}[Microplastique]
Un \cle{microplastique} est un fragment de plastique de taille inférieure à \SI{5}{mm}, issu de la dégradation des déchets plastiques sous l'effet du soleil, de la chaleur et de l'usure, ou rejeté directement sous forme de particules fines. Invisibles à l'œil nu pour la plupart, les microplastiques sont ingérés par les animaux aquatiques et se retrouvent dans la chaîne alimentaire, jusqu'à l'homme.
\end{definition}

\begin{popfigure}
\begin{tikzpicture}[
  font=\footnotesize,
  etape/.style={draw=popblue, fill=popblueL, rounded corners=3pt, minimum width=2.2cm, minimum height=1.1cm, align=center, text=popdark},
  fleche/.style={-{Stealth[length=2.5mm]}, very thick, poporange}
]
\node[etape, align=center] (s) at (0,0){Sachet\\jeté au sol};
\node[etape, align=center] (c) at (2.8,0){Caniveau\\colmaté};
\node[etape, align=center] (r) at (5.6,0){Rivière\\(Wouri, Sanaga)};
\node[etape, align=center] (o) at (8.4,0){Océan\\Atlantique};
\node[etape, fill=poporangeL, draw=poporange, align=center] (m) at (8.4,-2.4){Microplastiques\\($< \SI{5}{mm}$)};
\node[etape, fill=popgreenL, draw=popgreen, align=center] (p) at (5.6,-2.4){Poisson,\\crevette};
\node[etape, fill=poppinkL, draw=poppink, align=center] (h) at (2.8,-2.4){Homme\\(chaîne alimentaire)};
\draw[fleche] (s) -- (c);
\draw[fleche] (c) -- (r);
\draw[fleche] (r) -- (o);
\draw[fleche] (o) -- (m);
\draw[fleche] (m) -- (p);
\draw[fleche] (p) -- (h);
\draw[fleche, dashed, popmuted] (h.west) to[out=150,in=-150] node[left, text=popmuted, align=center, font=\scriptsize]{rejet\\sauvage} (s.west);
\end{tikzpicture}
\legende{La chaîne de la pollution plastique : du sachet jeté dans la rue jusqu'à l'assiette du consommateur. La flèche en pointillés rappelle que le cycle recommence tant que le rejet sauvage continue.}
\end{popfigure}

\begin{methode}[Démarche d'investigation : comment les scientifiques ont-ils établi cette chaîne ?]
\begin{enumerate}
\item \textbf{Observation :} des plastiques sont retrouvés dans l'estomac d'oiseaux marins, de tortues et de poissons morts échoués.
\item \textbf{Hypothèse :} les animaux confondent les fragments de plastique avec leurs proies naturelles.
\item \textbf{Expérience :} en laboratoire, on place des microplastiques dans l'eau de petits crustacés et de poissons ; on observe ensuite leurs organes au microscope.
\item \textbf{Résultat :} les particules de plastique sont bien ingérées et s'accumulent dans le tube digestif, voire passent dans le sang.
\item \textbf{Conclusion :} les microplastiques pénètrent la chaîne alimentaire dès la base (plancton) et se concentrent à chaque maillon, jusqu'à l'homme : c'est la \cle{bioaccumulation}.
\end{enumerate}
\end{methode}

\courssub{Les conséquences sanitaires de la pollution plastique}

\courssub{Des déchets qui abritent des maladies}

Regarde un sachet froissé ou une bouteille couchée dans la cour après la pluie : ils retiennent de petites poches d'eau stagnante. Or l'eau stagnante est le gîte larvaire idéal du moustique \emph{Anophèle}, qui transmet le \cle{paludisme}. Les déchets plastiques accumulés multiplient donc les sites de ponte des moustiques, juste à côté des habitations. Au Cameroun, où le paludisme reste l'une des premières causes de consultation médicale et de mortalité infantile, ce lien entre déchets plastiques et santé publique est très concret : moins de sachets traînant dans la cour, c'est moins d'eau stagnante, donc moins de moustiques.

\courssub{Brûler les plastiques : un faux remède, un vrai poison}

Face à l'accumulation des déchets, beaucoup de familles ont le réflexe de brûler les plastiques dans la cour ou au bord de la route. Cette pratique est très dangereuse.

\begin{attention}[Ne brûle jamais les plastiques !]
La combustion sauvage des plastiques, à basse température et à l'air libre, est incomplète. Elle libère des fumées noires chargées de gaz toxiques et de substances appelées \cle{dioxines}, reconnues comme \textbf{cancérigènes}. Ces fumées irritent les yeux et les bronches, aggravent l'asthme et contaminent le sol et les cultures voisines. Brûler un tas de sachets dans la cour, c'est transformer un problème de saleté en problème de santé pour toute la famille et les voisins. La bonne solution : trier, réduire, réutiliser et confier les déchets à la collecte ou au recyclage (section suivante).
\end{attention}

\begin{savaistu}[Que devient le plastique brûlé ou enfoui ?]
Les dioxines rejetées lors du brûlage des plastiques ne se dégradent presque pas : elles persistent dans le sol, contaminent les plantes, puis les animaux qui les mangent... et reviennent encore une fois vers l'homme par la chaîne alimentaire. Quant au plastique enfoui dans le sol, il le rend moins fertile en gênant la circulation de l'eau et des racines. Ni le feu ni la terre ne sont donc des solutions.
\end{savaistu}

\courssub{Les conséquences économiques}

La pollution plastique coûte cher, très cher, à l'État comme aux familles :

\begin{itemize}
\item \textbf{Mortalité du bétail :} dans les pâturages et les abords des villes, les chèvres, les moutons et les bœufs en liberté avalent des sachets en broutant. Les sachets s'accumulent dans leur estomac, bloquent la digestion et provoquent une mort lente. Pour un éleveur ou une famille, perdre un bœuf ou plusieurs chèvres est une perte financière considérable.
\item \textbf{Baisse de la pêche :} à Limbé, Kribi ou Douala, les pêcheurs constatent la diminution des captures dans les zones polluées et perdent du temps à démêler les sachets de leurs filets. Moins de poissons, c'est moins de revenus et une alimentation plus chère pour tous.
\item \textbf{Coût du curage :} les mairies et les entreprises d'hygiène (comme HYSACAM) dépensent chaque année des sommes énormes pour curer les caniveaux, déboucher les buses et ramasser les déchets sauvages : autant d'argent qui ne peut pas aller aux écoles, aux hôpitaux ou aux routes.
\item \textbf{Tourisme pénalisé :} des plages souillées à Kribi ou Limbé découragent les visiteurs, ce qui prive la région de revenus touristiques.
\end{itemize}

\begin{exoresolu}[Le coût caché d'un sachet]
\textbf{Énoncé.} Dans un quartier de Douala, une coopérative de 40 éleveurs perd en moyenne 3 chèvres par an à cause de l'ingestion de sachets plastiques. Une chèvre vaut en moyenne \SI{25000}{FCFA}. De plus, la mairie dépense \SI{1200000}{FCFA} par an pour curer les caniveaux du quartier, colmatés en grande partie par les plastiques. Calcule le coût annuel total de la pollution plastique pour ce quartier (pertes de bétail + curage).
\tcblower
\textbf{Solution détaillée.}
\begin{align*}
\text{Pertes de bétail} &= 40 \times 3 \times \SI{25000}{FCFA} \\
&= 120 \times \SI{25000}{FCFA} \\
&= \SI{3000000}{FCFA} \\
\text{Coût total} &= \SI{3000000}{FCFA} + \SI{1200000}{FCFA} \\
&= \SI{4200000}{FCFA}
\end{align*}
\textbf{Conclusion :} la pollution plastique coûte chaque année \SI{4200000}{FCFA} à ce seul quartier, soit bien plus que le prix de tous les sachets utilisés ! Ce calcul montre que « jeter n'importe où » a un prix collectif énorme : prévenir la pollution (trier, réutiliser, recycler) est un excellent investissement.
\end{exoresolu}

\courssub{Le cadre légal : un savoir citoyen}

Face à l'ampleur du problème, l'État du Cameroun a pris des mesures que tout citoyen — et donc tout candidat au BEPC — doit connaître.

\begin{savaistu}[Le Cameroun dit non aux sachets non biodégradables]
Depuis un décret de \textbf{2012} (renforcé par des arrêtés conjoints des ministères de l'Environnement et de l'Industrie dans les années suivantes), le Cameroun \textbf{interdit la fabrication, l'importation et la commercialisation des sachets plastiques non biodégradables de faible épaisseur (moins de 60 microns)}. Ce sont précisément ces sachets très fins, utilisés une seule fois, qui volent au vent, bouchent les caniveaux et se fragmentent en microplastiques. Des opérations de saisie et de destruction de stocks interdits sont régulièrement menées. Malgré cela, ces sachets continuent de circuler : la loi n'est efficace que si chaque citoyen la respecte et refuse ces emballages.
\end{savaistu}

\begin{aretenir}[L'essentiel de la section]
\begin{itemize}
\item Les déchets plastiques rejetés dans la nature s'accumulent car le plastique ne se biodégrade pratiquement pas.
\item La chaîne de pollution : rejet sauvage $\rightarrow$ caniveau $\rightarrow$ rivière $\rightarrow$ océan $\rightarrow$ fragmentation en \cle{microplastiques} (moins de \SI{5}{mm}) $\rightarrow$ ingestion par les animaux $\rightarrow$ chaîne alimentaire jusqu'à l'homme.
\item Conséquences sanitaires : eaux stagnantes dans les déchets = gîtes larvaires des moustiques (paludisme) ; brûlage sauvage = fumées toxiques et dioxines cancérigènes (pratique interdite de fait, à proscrire absolument).
\item Conséquences économiques : mort du bétail, baisse de la pêche, coût du curage des caniveaux, tourisme pénalisé.
\item Au Cameroun, les sachets plastiques non biodégradables de moins de 60 microns sont interdits depuis 2012 : c'est une règle de citoyenneté à connaître et à faire respecter.
\end{itemize}
\end{aretenir}

\begin{attention}[Pièges fréquents au BEPC]
\begin{itemize}
\item Ne dis pas que le plastique « disparaît » ou « fond » dans la nature : il se \textbf{fragmente} en morceaux de plus en plus petits, mais il reste du plastique.
\item Ne confonds pas \emph{biodégradable} et \emph{fragmentable} : un sachet qui se casse en petits morceaux n'est pas biodégradé, il est devenu des microplastiques.
\item Ne propose jamais le brûlage comme solution de gestion des déchets plastiques dans une copie : c'est une pratique toxique, à proscrire.
\item Retiens le chiffre clé : seuil de \SI{5}{mm} pour les microplastiques, et 60 microns / 2012 pour l'interdiction des sachets au Cameroun.
\end{itemize}
\end{attention}

Maintenant que tu comprends l'ampleur du problème, une question s'impose : que faire concrètement de tous ces déchets plastiques ? C'est l'objet de la section suivante, consacrée à la gestion des déchets : la règle des 3R et le recyclage, où tu découvriras que le plastique usagé peut aussi devenir une ressource.

\coursec{Gestion des déchets plastiques : les 3R et le recyclage}

Face à l'ampleur de la pollution plastique décrite dans la section précédente — sachets qui bouchent les caniveaux de Yaoundé ou de Douala, bouteilles qui jonchent les berges du Wouri, microplastiques dans nos poissons — l'inaction n'est pas une option. La solution ne viendra pas d'une seule mesure miracle, mais d'une stratégie combinée : \textbf{produire moins de déchets}, \textbf{prolonger la vie des objets} et \textbf{transformer ce qui ne peut plus servir}. C'est tout le sens de la règle des \cle{3R}, adoptée mondialement et inscrite dans la loi camerounaise de 2012 relative à la gestion des déchets.

\courssub{La règle des 3R : un ordre de priorité}

La règle des 3R — \textbf{Réduire}, \textbf{Réutiliser}, \textbf{Recycler} — n'est pas une simple liste : c'est une \cle{hiérarchie}. On agit d'abord sur la source (Réduire), puis sur l'usage (Réutiliser), enfin sur le déchet (Recycler). Cet ordre est crucial : recycler consomme de l'énergie et de l'eau ; réutiliser en consomme moins ; réduire n'en consomme aucune.

\begin{methode}[Appliquer la règle des 3R au quotidien]
\begin{enumerate}
    \item \textbf{Réduire (priorité n°1)} : Refuser le superflu. Au marché, dis « non » au sachet plastique pour trois tomates ; apporte ton panier ou ton sac en tissu. Préfère l'achat en vrac (riz, haricots, arachides) à la dose préemballée. Choisis des produits peu ou pas emballés (savon solide plutôt que gel douche en flacon).
    \item \textbf{Réutiliser (priorité n°2)} : Avant de jeter, demande-toi : « Cela peut-il resservir ? » Une bouteille d'eau vide devient une gourde pour l'école ou un arrosoir percé pour le jardin. Un seau de peinture propre sert à puiser l'eau. Une calebasse ou un bol en inox remplacent les gobelets jetables.
    \item \textbf{Recycler (priorité n°3)} : Quand l'objet est vraiment en fin de vie, trie-le pour qu'il devienne matière première. Dépose les bouteilles, sachets et flacons dans les bornes de tri ou remets-les aux collecteurs agréés. Ne mélange pas plastiques et ordures ménagères.
\end{enumerate}
\end{methode}

\begin{popfigure}
\begin{tikzpicture}[>=Stealth, font=\small]
% Trois flèches circulaires style logo 3R
\def\r{2.2} % rayon centre
\def\angle{120}
\foreach \i/\label/\col in {0/Réduire/popgreen, 1/Réutiliser/popblue, 2/Recycler/poporange} {
    \pgfmathsetmacro{\start}{90 - \i*\angle}
    \pgfmathsetmacro{\end}{90 - (\i+1)*\angle}
    % Flèche courbe (sens horaire)
    \draw[very thick, \col!80!black, ->] (\start:\r) arc (\start:\end:\r) node[midway, sloped, above=2pt, font=\bfseries\small] {\label};
}
% Centre
\node[circle, draw=popdark, thick, minimum size=2.5cm, fill=popcream, align=center] at (0,0){\textbf{3R}\\\textit{Ordre de priorité}};
% Lignes de séparation radiales (optionnel)
\foreach \i in {0,1,2} {
    \pgfmathsetmacro{\a}{90 - \i*\angle - 60}
    \draw[popmuted, dashed, very thin] (0,0) -- (\a:2.7);
}
\end{tikzpicture}
\legende{La règle des 3R : Réduire (vert) avant Réutiliser (bleu) avant Recycler (orange). C'est un cycle vertueux à appliquer dans cet ordre.}
\end{popfigure}

\begin{exemplebox}[Au lycée et à la maison]
\begin{itemize}
    \item \textbf{Réduire} : J'apporte ma gourde en inox (fini les sachets « pure water » jetés par terre). J'utilise une boîte hermétique pour mon goûter (arachides, beignets) au lieu de papier alu ou sachet.
    \item \textbf{Réutiliser} : Les vieux cahiers à pages blanches restantes deviennent brouillons. Les bidons d'huile vidés et lavés (attention : jamais pour l'alimentaire !) servent à stocker l'eau de pluie pour les fleurs.
    \item \textbf{Recycler} : À l'école, deux poubelles en classe : une jaune (plastiques, métaux, papier) et une verte (reste). À la maison, je donne les bouteilles PET au collecteur qui passe le samedi.
\end{itemize}
\end{exemplebox}

\courssub{Le recyclage des thermoplastiques : de la poubelle au nouvel objet}

Seuls les \cle{thermoplastiques} (PET, PEHD, PP, PVC, PS) peuvent être recyclés par \cle{refonte} : on les chauffe, ils fondent, on les remoule. Les thermodurcissables (bakélite, résines époxy, mélamine) ne fondent pas ; ils se décomposent en brûlant. Ils ne sont \textbf{pas recyclables} par cette voie (voir l'encadré \textbf{Attention} ci-dessous).

\begin{attention}[Thermodurcissables non recyclables par refonte]
Les \textbf{thermodurcissables} (bakélite, résines époxy, mélamine, polyesters insaturés) \textbf{ne fondent pas} : chauffés, ils se décomposent (charrent) irréversiblement. Ils ne peuvent \textbf{pas} être recyclés par refonte (recyclage matière). Leur valorisation passe par le recyclage énergétique (incinération avec récupération de chaleur) ou le broyage comme charge inerte (peu valorisant). Ne les mets pas dans le bac de tri plastiques (codes 1 à 6).
\end{attention}

\begin{definition}[Recyclage matière des plastiques]
Le \textbf{recyclage} (ou recyclage matière) est l'ensemble des opérations industrielles qui transforment un déchet plastique en \textbf{nouvelle matière première} (granulats, paillettes) ou en \textbf{nouvel objet}, sans changer sa nature chimique fondamentale. C'est une boucle : déchet $\to$ matière $\to$ produit $\to$ déchet...
\end{definition}

\begin{experience}[Observer le tri par code d'identification (TP ou démonstration)]
\textbf{Question de départ} : Comment l'usine de recyclage sait-elle quel plastique fondre avec quel autre ?
\textbf{Matériel} : Échantillons de déchets plastiques propres (bouteille d'eau, bidon de lait, bouchon, pot de yaourt, sachet, gobelet mousse). Loupe.
\textbf{Protocole} :
\begin{enumerate}
    \item Retourne chaque objet. Cherche le \textbf{triangle de Möbius} (trois flèches) avec un \textbf{chiffre de 1 à 7} et des \textbf{lettres} en dessous.
    \item Note le code et le nom du polymère dans un tableau.
    \item Regroupe les objets ayant le même chiffre.
\end{enumerate}
\begin{center}
\begin{tabularx}{\linewidth}{|c|X|c|X|}\hline
\textbf{Code} & \textbf{Polymère (Abréviation)} & \textbf{Exemple courant au Cameroun} & \textbf{Recyclabilité locale} \\ \hline
1 & PET (Polyéthylène téréphtalate) & Bouteilles d'eau (Tangui, Supermont), sodas & \cellcolor{popgreenL} Oui, très recherché \\ \hline
2 & PEHD (Polyéthylène haute densité) & Bouteilles de lait, bidons d'huile, lessive, shampoing & \cellcolor{popgreenL} Oui \\ \hline
3 & PVC (Polychlorure de vinyle) & Tuyaux, gaines électriques, certains flacons & \cellcolor{poporangeL} Difficile (chlore) \\ \hline
4 & PELD (Polyéthylène basse densité) & Sachets « pure water », films étirables, sacs poubelles & \cellcolor{poporangeL} Oui (si collectés propres) \\ \hline
5 & PP (Polypropylène) & Bouchons, pots de yaourt, seaux, chaises, pare-chocs moto & \cellcolor{popgreenL} Oui \\ \hline
6 & PS (Polystyrène) & Gobelets mousse, barquettes viande, emballages fragiles & \cellcolor{popredL} Rarement (volume/poids) \\ \hline
7 & Autres (PC, PLA, mélanges) & Biberons, lunettes, pièces techniques & \cellcolor{popredL} Quasiment jamais \\ \hline
\end{tabularx}
\end{center}
\textbf{Interprétation} : Le code permet le \textbf{tri par type}. Mélanger du PET (code 1) et du PP (code 5) donne un mélange infusible ou de mauvaise qualité. Le tri rigoureux \emph{à la source} (par toi, à la maison) est la condition n°1 du recyclage.
\textbf{Conclusion} : Chaque plastique a son code. Trie par code = recycler possible. Mélange = déchet final.
\end{experience}

\begin{popfigure}
\begin{tikzpicture}[>=Stealth, font=\small, node distance=1.8cm and 1.2cm]
% Nœuds du cycle
\node[draw, rounded corners, fill=popblueL, thick, minimum width=2.2cm, minimum height=1.1cm, align=center] (collecte) {\textbf{1. Collecte}\\\textit{Bornes, collecteurs, ramassage}};
\node[draw, rounded corners, fill=popgreenL, thick, minimum width=2.2cm, minimum height=1.1cm, align=center, right=of collecte] (tri) {\textbf{2. Tri}\\\textit{Séparation par code (1 à 7)}\\\textit{Optique, manuelle, densité}};
\node[draw, rounded corners, fill=poporangeL, thick, minimum width=2.2cm, minimum height=1.1cm, align=center, right=of tri] (broyage) {\textbf{3. Broyage}\\\textit{Paillettes / Flocons}\\\textit{Lavage + séchage}};
\node[draw, rounded corners, fill=poppurpleL, thick, minimum width=2.2cm, minimum height=1.1cm, align=center, below right=of broyage] (fusion) {\textbf{4. Fusion / Extrusion}\\\textit{Chauffage $\sim$ 200--280°C}\\\textit{Filtration impuretés}};
\node[draw, rounded corners, fill=poppinkL, thick, minimum width=2.2cm, minimum height=1.1cm, align=center, left=of fusion] (granulats) {\textbf{5. Granulats}\\\textit{Nouvelle matière première}\\\textit{Sacs de 25 kg}};
\node[draw, rounded corners, fill=popgoldL, thick, minimum width=2.2cm, minimum height=1.1cm, align=center, left=of granulats] (nouvel) {\textbf{6. Nouvel objet}\\\textit{Injection, soufflage,}\\\textit{extrusion : bouteilles, tuyaux,}\\\textit{pavés, fibres textiles}};
% Flèches principales
\draw[very thick, popblue, ->] (collecte) -- (tri) node[midway, above] {Tri à la source};
\draw[very thick, popgreen, ->] (tri) -- (broyage) node[midway, above] {Par type};
\draw[very thick, poporange, ->] (broyage) -- (fusion) node[midway, right] {Propre \& sec};
\draw[very thick, poppurple, ->] (fusion) -- (granulats) node[midway, below] {Refonte};
\draw[very thick, poppink, ->] (granulats) -- (nouvel) node[midway, below] {Moulage};
\draw[very thick, popgold, ->] (nouvel) -- ++(-2.5,0) |- (collecte) node[midway, left, align=center] {Usage $\to$ Déchet};
% Boucle interne : recyclage en boucle fermée (bouteille -> bouteille)
\draw[dashed, thick, popdark, ->] (granulats) to[out=150, in=30, looseness=1.5] node[above, align=center, font=\tiny] {Boucle fermée\\ (bouteille $\to$ bouteille)} (tri);
\end{tikzpicture}
\legende{Cycle du recyclage matière d'un thermoplastique (ex. PET). La boucle fermée (pointillés) est l'idéal : même usage. La boucle ouverte (plein trait) : usage différent (ex. bouteille $\to$ fibre polaire $\to$ pavé).}
\end{popfigure}

\begin{exoresolu}[Calculer le gain matière du recyclage]
\textbf{Énoncé} : Une usine de recyclage à Douala traite \SI{500}{tonnes} de bouteilles PET par an. Le rendement global (collecte $\to$ granulats) est de \SI{85}{\%}. Combien de tonnes de granulats PET sont produites ? Sachant que \SI{1}{tonne} de PET vierge nécessite \SI{1,9}{tonnes} de pétrole brut, quelle quantité de pétrole brut est économisée grâce à ce recyclage ?
\tcblower
\textbf{Solution détaillée} :
\begin{enumerate}
    \item \textbf{Masse de granulats PET obtenus} :
    \[ m_{\text{granulats}} = \SI{500}{t} \times 0,85 = \SI{425}{tonnes} \]
    \item \textbf{Pétrole brut économisé} :
    Pour produire \SI{425}{tonnes} de PET vierge, il aurait fallu :
    \[ m_{\text{pétrole}} = \SI{425}{t} \times \SI{1,9}{t/t} = \SI{807,5}{tonnes} \]
    \textbf{Conclusion} : L'usine produit \SI{425}{tonnes} de granulats et économise \SI{807,5}{tonnes} de pétrole brut par an.
\end{enumerate}
\end{exoresolu}

\newpage

\coursec{Activité d'intégration}

\coursec{Les matières plastiques}

\courssub{Objectifs de la leçon}

À la fin de cette leçon, tu seras capable de :
\begin{itemize}
\item identifier un plastique courant à partir de son \cle{code de recyclage} (symbole triangulaire numéroté de 1 à 7) et de ses propriétés physiques ;
\item distinguer expérimentalement un \cle{thermoplastique} d'un \cle{thermodurcissable} et prédire leur comportement au chauffage ;
\item trier des déchets plastiques selon leur recyclabilité effective dans le contexte camerounais ;
\item exploiter des documents chiffrés (durées de décomposition, articles de presse) pour évaluer l'impact de la pollution plastique sur l'environnement et la santé ;
\item proposer un plan d'action concret fondé sur la règle des \cle{3R} (Réduire, Réutiliser, Recycler) pour résoudre un problème environnemental réel.
\end{itemize}

\courssub{Situation déclenchante : le défi de la cité scolaire}

\textbf{Le problème de l'école bilingue de Bonabéri (Douala).}\par
L'établissement compte \num{1200} élèves. Chaque jour, la cantine et les vendeuses ambulantes distribuent boissons et biscuits dans des emballages plastiques. À la fin de chaque journée, on estime que \num{600} sachets et bouteilles vides sont jetés dans la cour ; une grande partie finit dans le caniveau qui longe l'école. Lors des fortes pluies d'octobre, ce caniveau déborde : l'eau stagnante envahit la cour pendant deux jours et les moustiques prolifèrent. Le principal lance un défi aux classes de 3ème : « Proposez un plan d'action scientifiquement fondé pour que notre école devienne un exemple de quartier \emph{zéro plastique sauvage}. »

\textbf{Les huit déchets collectés.}\par
Ta classe a ramassé huit objets représentatifs dans la cour et le caniveau. Chaque objet porte (ou non) un code de recyclage :

\begin{center}
\begin{tabularx}{\linewidth}{|c|X|c|l|}
\hline
\rowcolor{popblueL} \textbf{N°} & \textbf{Objet} & \textbf{Code} & \textbf{Observation au toucher} \\
\hline
1 & Sachet d'eau minérale vide & 4 & Souple, transparent \\
\hline
2 & Bouteille de boisson gazeuse & 1 & Rigide, transparente \\
\hline
3 & Gobelet de yaourt & 6 & Léger, cassant \\
\hline
4 & Morceau de tuyau d'eau & 3 & Rigide, gris \\
\hline
5 & Chaise cassée & 5 & Dure, légère \\
\hline
6 & Prise électrique brisée & aucun & Très dure, noire \\
\hline
7 & Bidon d'huile de vidange & 2 & Opaque, résistant \\
\hline
8 & Emballage de biscuits métallisé & 7 & Souple, plusieurs couches \\
\hline
\end{tabularx}
\end{center}

\textbf{Fiches de résultats expérimentaux (réalisés par l'enseignant sous hotte).}\par
Pour des raisons de sécurité, le chauffage des échantillons a été réalisé par l'enseignant. Voici les résultats consignés :

\begin{center}
\begin{tabularx}{\linewidth}{|c|X|X|}
\hline
\rowcolor{popblueL} \textbf{Échantillon} & \textbf{Chauffage doux (\SI{80}{\celsius})} & \textbf{Refroidissement puis nouveau chauffage} \\
\hline
Bouteille (code 1) & Se ramollit, se déforme & Se ramollit à nouveau \\
\hline
Bidon (code 2) & Se ramollit & Se ramollit à nouveau \\
\hline
Prise électrique & Aucune déformation & Noircit et se carbonise sans fondre \\
\hline
Gobelet (code 6) & Se ramollit & Se ramollit à nouveau \\
\hline
\end{tabularx}
\end{center}

\textbf{Document 1 : durées de décomposition dans la nature.}\par
Sachet plastique (PEBD) : \num{100} à \num{400} ans ; bouteille en PET : \num{450} ans ; gobelet en polystyrène : \num{500} ans ; mégot de cigarette : \num{2} ans ; pelure de banane : \num{6} mois.

\textbf{Document 2 : extrait d'article de presse (Douala, octobre 2025).}\par
« Après trois heures de pluie, plusieurs quartiers de la ville sont restés sous l'eau. Les services d'hygiène ont retiré plus de \num{12} tonnes de déchets plastiques des caniveaux du seul arrondissement de Douala IV. Selon les spécialistes, les sachets et bouteilles qui bouchent les canalisations aggravent chaque année les inondations et favorisent la prolifération des moustiques vecteurs du paludisme. »

\courssub{Consignes de l'activité}

Par groupes de quatre, traitez les tâches suivantes en rédigeant soigneusement vos réponses.

\begin{enumerate}
\item \textbf{Identification.} À l'aide du tableau des codes de recyclage, nomme le plastique (sigle et nom complet) correspondant à chacun des huit objets collectés.
\item \textbf{Prédiction.} Sans consulter les fiches de résultats, prédis pour chaque objet s'il s'agit d'un thermoplastique ou d'un thermodurcissable. Justifie ta prédiction pour la prise électrique.
\item \textbf{Vérification.} Exploite les fiches de résultats expérimentaux : quel critère permet de distinguer les deux familles ? Vérifie tes prédictions pour les quatre échantillons testés et conclus.
\item \textbf{Tri.} Sachant que la filière de recyclage de Douala n'accepte que les thermoplastiques purs (codes 1, 2, 4, 5), classe les huit objets en « recyclables » et « non recyclables », en justifiant chaque refus.
\item \textbf{Analyse de documents.} À l'aide des documents 1 et 2, explique en quelques phrases pourquoi les déchets plastiques de la cour du collège aggravent les inondations et les maladies. Calcule le nombre approximatif de sachets et bouteilles jetés par l'école en une année scolaire de \num{180} jours.
\item \textbf{Plan d'action.} Rédige un plan d'action en exactement cinq mesures concrètes pour l'école, dont au moins une mesure pour chacun des trois R (Réduire, Réutiliser, Recycler), une mesure de sensibilisation et une mesure de collecte. Pour chaque mesure, précise le résultat attendu.
\item \textbf{Restitution.} Présente ton plan à la classe en cinq minutes ; le meilleur plan sera transmis au principal.
\end{enumerate}

\courssub{Ressources à mobiliser : savoirs de référence}

\begin{definition}[Matière plastique et polymère]
Une \cle{matière plastique} est un matériau solide, organique, constitué principalement de \cle{polymères}. Un polymère est une molécule géante (macromolécule) obtenue par l'assemblage répétitif de petites molécules appelées \cle{monomères}, le plus souvent issues du \cle{pétrole} (raffiné à la SONARA à Limbé).
\end{definition}

\begin{propriete}[Les deux grandes familles de plastiques]
\begin{itemize}
\item \textbf{Thermoplastique :} se ramollit à chaque chauffage et durcit en refroidissant. Cette transition est \textbf{réversible}. Il est \textbf{recyclable} par fusion (ex. : PET, PEHD, PEBD, PP, PVC, PS).
\item \textbf{Thermodurcissable :} durcit définitivement après sa première mise en forme (réticulation irréversible). Chauffé à nouveau, il ne fond pas : il se décompose et se carbonise. Il est \textbf{non recyclable} par fusion (ex. : bakélite des prises électriques, résines époxy, mélamine).
\end{itemize}
\end{propriete}

\begin{methode}[Identifier le code de recyclage d'un objet plastique]
\begin{enumerate}
\item Repère le symbole triangulaire (anneau de Möbius) gravé ou moulé sur l'objet (souvent sous les bouteilles, au fond des gobelets).
\item Lis le chiffre au centre (de 1 à 7).
\item Reporte-toi au tableau ci-dessous pour connaître le sigle et le nom du polymère.
\end{enumerate}
\end{methode}

\begin{center}
\begin{tabularx}{\linewidth}{|c|l|X|}
\hline
\rowcolor{popblueL} \textbf{Code} & \textbf{Sigle} & \textbf{Nom complet \& Usages courants} \\
\hline
1 & PET & Polyéthylène téréphtalate — bouteilles d'eau, de soda, d'huile \\
\hline
2 & PEHD & Polyéthylène haute densité — bidons de lait, d'huile moteur, de lessive, bouchons \\
\hline
3 & PVC & Polychlorure de vinyle — tuyaux d'évacuation, gaines électriques, cartes bancaires \\
\hline
4 & PEBD & Polyéthylène basse densité — sachets d'eau, sacs poubelles, films étirables \\
\hline
5 & PP & Polypropylène — chaises, seaux, bouchons, pièces auto, contenants micro-ondables \\
\hline
6 & PS & Polystyrène — gobelets yaourt, barquettes, emballages fragiles (blanc cassant) \\
\hline
7 & Autres & Tous les autres plastiques (PC, PA, bioplastiques) et \textbf{mélanges multicouches} (sachets chips, café) \\
\hline
\end{tabularx}
\end{center}

\begin{experience}[Distinction thermoplastique / thermodurcissable par chauffage]
\textbf{Matériel :} échantillons de plastiques (bouteille PET, bidon PEHD, gobelet PS, prise électrique), pinces à creuset, bec Bunsen ou plaque chauffante, hotte aspirante, lunettes de protection, gants thermiques.\par
\textbf{Protocole :}
\begin{enumerate}
\item Prendre un petit morceau de l'échantillon avec la pince.
\item Le chauffer doucement (\SI{80}{\celsius} ou flamme douce) pendant \SI{30}{s}.
\item Observer : ramollissement / déformation ? ou rien / noircissement ?
\item Laisser refroidir, puis chauffer à nouveau le même morceau.
\item Noter si le comportement est réversible (ramollit à nouveau) ou irréversible (carbonisation).
\end{enumerate}
\textbf{Sécurité :} Ne jamais inhaler les fumées (risque de chlorure d'hydrogène pour PVC, styrène pour PS, phénols pour bakélite). Manipulation réservée à l'enseignant sous hotte.
\textbf{Interprétation attendue :} Ramollissement réversible $\rightarrow$ thermoplastique. Carbonisation sans fusion $\rightarrow$ thermodurcissable.
\end{experience}

\begin{savaistu}[Le plastique au Cameroun : de la SONARA au recyclage]
Le Cameroun produit une partie de ses polymères via la SONARA (Limbé) et l'industrie pétrochimique (Kribi). Cependant, la majorité des objets finis sont importés. À Douala et Yaoundé, des entreprises comme \emph{Redplast}, \emph{Cameroon Plastic Recycling} ou des coopératives de trieurs (souvent des femmes au quartier \emph{New Bell} ou \emph{Bonabéri}) collectent le PET (code 1) et le PEHD (code 2) pour les broyer en paillettes exportées ou transformées localement en tuyaux, sacs, pavés. Le PVC (code 3) est refusé car sa combustion dégage du chlorure d'hydrogène (toxique, corrosif). Les sachets PEBD (code 4) sont peu collectés car trop légers et sales (non rentables). Les thermodurcissables (prises, poignées de casserole) et les multicouches (code 7) ne trouvent aucune filière : ils finissent en décharge ou brûlés à l'air libre, polluant l'air et les sols.
\end{savaistu}

\begin{aretenir}[Essentiel pour le BEPC]
\begin{itemize}
\item Un \cle{plastique} = polymères (macromolécules) issues mostly du pétrole.
\item \cle{Thermoplastique} : fond \textbf{réversiblement} $\rightarrow$ \textbf{recyclable} (codes 1, 2, 3, 4, 5, 6).
\item \cle{Thermodurcissable} : \textbf{ne fond pas}, carbonise $\rightarrow$ \textbf{non recyclable} (prise électrique, bakélite).
\item \cle{Codes de recyclage} : 1=PET, 2=PEHD, 3=PVC, 4=PEBD, 5=PP, 6=PS, 7=Autres/Multicouches.
\item \cle{3R} (hiérarchie) : 1. \textbf{Réduire} (ne pas produire), 2. \textbf{Réutiliser} (prolonger la vie), 3. \textbf{Recycler} (transformer la matière). Le recyclage consomme de l'énergie.
\item Un plastique classique n'est \textbf{ni biodégradable, ni compostable} : il persiste \num{100} à \num{500} ans, bouche les caniveaux (inondations), fragmente en microplastiques, pollue sols, eaux et chaîne alimentaire.
\item \textbf{Brûler du plastique à l'air libre est interdit et dangereux} : fumées toxiques (dioxines, furanes, HCl, styrène).
\end{itemize}
\end{aretenir}

\begin{attention}[Pièges fréquents au BEPC]
\begin{itemize}
\item Confondre « code de recyclage » et « ordre de tri » : le code 3 (PVC) et 6 (PS) sont des thermoplastiques mais souvent \textbf{non recyclés} localement (filière absente ou non rentable).
\item Croire qu'un objet « plastique » est automatiquement recyclable : les thermodurcissables (code absent ou 7) et les multicouches (code 7) ne le sont pas.
\item Oublier que le tri exige des plastiques \textbf{purs} : un seul objet PVC dans un lot de PET gâche le recyclage (acide chlorhydrique lors de l'extrusion).
\item Écrire « biodégradable » pour un plastique classique : \textbf{faux}. Seuls certains bioplastiques (PLA, PHA) le sont, en conditions industrielles.
\end{itemize}
\end{attention}

\courssub{Démarche d'investigation et corrigé commenté}

\begin{exoresolu}[Corrigé détaillé de l'activité]
Énoncé : identifier les huit déchets, déterminer leur famille, les trier, analyser les documents et bâtir un plan d'action en cinq mesures.
\tcblower
\textbf{Étape 1 : Identification des plastiques (consigne 1).}\par
On lit le code triangulaire et on utilise le tableau de référence :
\begin{center}
\begin{tabularx}{\linewidth}{|c|X|l|}
\hline
\rowcolor{popblueL} \textbf{N°} & \textbf{Objet} & \textbf{Plastique (sigle + nom)} \\
\hline
1 & Sachet d'eau (code 4) & PEBD — Polyéthylène basse densité \\
\hline
2 & Bouteille soda (code 1) & PET — Polyéthylène téréphtalate \\
\hline
3 & Gobelet yaourt (code 6) & PS — Polystyrène \\
\hline
4 & Tuyau eau (code 3) & PVC — Polychlorure de vinyle \\
\hline
5 & Chaise (code 5) & PP — Polypropylène \\
\hline
6 & Prise électrique (sans code) & Thermodurcissable (type bakélite / résine phénolique) \\
\hline
7 & Bidon huile (code 2) & PEHD — Polyéthylène haute densité \\
\hline
8 & Emballage biscuits (code 7) & Multicouche (mélange PE / Alu / PET / colle) — « Autres » \\
\hline
\end{tabularx}
\end{center}

\textbf{Étape 2 : Prédiction de la famille (consigne 2).}\par
Les objets 1, 2, 3, 4, 5, 7 sont des emballages ou objets moulés courants, souples ou rigides, qui fondent à l'usine : on prédit \textbf{thermoplastiques}.\\
La prise électrique (objet 6) doit supporter l'échauffement par effet Joule sans se déformer : on prédit \textbf{thermodurcissable} (sinon court-circuit, incendie).\\
L'emballage multicouche (objet 8) est un mélange complexe : comportement incertain, mais non recyclable par fusion classique.

\textbf{Étape 3 : Vérification expérimentale (consigne 3).}\par
\textbf{Critère de distinction :} un thermoplastique ramollit de façon \textbf{réversible} à chaque chauffage ; un thermodurcissable \textbf{ne fond pas} et se carbonise (décomposition thermique).\\
\textbf{Exploitation des fiches :}
\begin{itemize}
\item Bouteille (1), Bidon (2), Gobelet (6) : ramollissent aux deux chauffages $\rightarrow$ \textbf{thermoplastiques confirmés}.
\item Prise électrique : ne se déforme pas, noircit, carbonise $\rightarrow$ \textbf{thermodurcissable confirmé}.
\end{itemize}
Prédictions validées.

\textbf{Étape 4 : Tri recyclables / non recyclables à Douala (consigne 4).}\par
\textbf{Recyclables (filière locale : codes 1, 2, 4, 5 purs) :}
\begin{itemize}
\item Objet 2 (Bouteille PET, code 1) — accepté.
\item Objet 7 (Bidon PEHD, code 2) — accepté.
\item Objet 1 (Sachet PEBD, code 4) — accepté \emph{théoriquement}, mais souvent refusé car trop souple, sale, non rentable.
\item Objet 5 (Chaise PP, code 5) — accepté si collecté séparément (gros volume).
\end{itemize}
\textbf{Non recyclables (justifications) :}
\begin{itemize}
\item Objet 3 (Gobelet PS, code 6) : thermoplastique mais filière absente (faible densité, cassant, non rentable).
\item Objet 4 (Tuyau PVC, code 3) : thermoplastique mais \textbf{refusé} (combustion $\rightarrow$ HCl toxique, corrosion des machines).
\item Objet 6 (Prise électrique) : thermodurcissable, ne fond pas, impossible à broyer/fondre.
\item Objet 8 (Emballage code 7) : multicouche indissociable, impossible à trier chimiquement.
\end{itemize}

\textbf{Étape 5 : Analyse des documents (consigne 5).}\par
\textbf{Explication :} Les sachets (PEBD, \num{100}--\num{400} ans) et bouteilles (PET, \num{450} ans) jetés dans la cour sont emportés par le ruissellement vers le caniveau. Ils s'y accumulent (Doc 1 : persistance siècles) et forment des barrages. L'eau de pluie ne s'écoule plus $\rightarrow$ inondations (Doc 2). L'eau stagnante devient gîte larvaire pour moustiques \emph{Anopheles} $\rightarrow$ paludisme.\\
\textbf{Calcul :} $\num{600} \text{ emballages/jour} \times \num{180} \text{ jours} = \num{108000} \text{ emballages/an}$.
Soit plus de \textbf{100 000 objets} par an pour une seule école, menaçant directement le caniveau et la santé du quartier.

\textbf{Étape 6 : Plan d'action en cinq mesures (consigne 6).}\par
\begin{enumerate}
\item \textbf{Réduire (source) :} Interdire la vente de boissons en sachets/gobelets jetables à la cantine ; imposer gourde et gobelet réutilisables pour chaque élève. \textit{Résultat : suppression de \SI{50}{\%} des emballages quotidiens (\num{300} objets/jour en moins).}
\item \textbf{Réutiliser :} Atelier « Bouteilles utiles » : transformer les bouteilles PET collectées en arrosoirs (bouchon percé), pots de pépinière (coupées), briques écologiques (remplies de sable tassé) pour murets du jardin scolaire. \textit{Résultat : \num{200} bouteilles/an détournées de la poubelle, végétalisation de l'école.}
\item \textbf{Recycler :} Installer trois bacs de tri colorés (Jaune : PET/PEHD codes 1 \& 2 ; Bleu : PP/PEBD codes 4 \& 5 ; Gris : autres). Convention avec un collecteur agréé de Bonabéri (achat PET \SI{150}{FCFA/kg}). \textit{Résultat : \SI{50}{kg} PET/an vendus, fonds reversés au club environnement.}
\item \textbf{Collecte :} Opération « Samedi Propre » mensuelle : point de collecte ouvert aux riverains devant l'école, pesée publique, affichage des tonnes détournées du caniveau. \textit{Résultat : mobilisation quartier, \SI{200}{kg} plastiques/an collectés hors école.}
\item \textbf{Sensibilisation :} Affiches élèves « Mon code, mon geste », sketch « Le caniveau qui tousse » au conseil de classe, porte-à-porte familles pour expliquer tri et danger brûlage. \textit{Résultat : \SI{80}{\%} familles du quartier informées, changement de comportement visible.}
\end{enumerate}
\end{exoresolu}

\courssub{Bilan synthétique de la leçon}

Cette leçon a permis de construire plusieurs compétences clés :
\begin{itemize}
\item \textbf{Savoir identifier :} le code de recyclage est l'outil du citoyen pour nommer le polymère (PET, PEHD, PVC, PEBD, PP, PS, Autres).
\item \textbf{Savoir classer :} l'expérience de chauffage (ramollissement réversible vs carbonisation) distingue sans ambiguïté thermoplastiques et thermodurcissables. Cette distinction décide du devenir du déchet : seul le thermoplastique pur peut être refondu.
\item \textbf{Comprendre l'impact :} la persistance des plastiques (\num{100}--\num{500} ans) transforme un geste individuel (jeter un sachet) en catastrophe collective (caniveaux bouchés, inondations, paludisme à Douala).
\item \textbf{Agir responsable :} la règle des 3R est une \textbf{hiérarchie} : \textbf{Réduire} (le plus efficace) $>$ \textbf{Réutiliser} $>$ \textbf{Recycler} (coûte de l'énergie). Le tri rigoureux (plastiques purs) est la condition sine qua non du recyclage.
\end{itemize}
Transformer ces savoirs en plan d'action concret pour son école, c'est exercer sa citoyenneté scientifique.

\begin{popfigure}
\begin{tikzpicture}[scale=0.9, every node/.style={font=\small}]
% Recycling symbol triangle with number
\def\drawTriangle#1#2#3#4{
  \begin{scope}[shift={(#1,#2)}]
    \draw[popdark, thick] (0,0) -- (1.5,0) -- (0.75,1.3) -- cycle;
    \draw[popblue, thick, ->] (0.75,0.2) arc (90:210:0.5);
    \draw[popblue, thick, ->] (0.75,0.2) arc (90:-30:0.5);
    \draw[popblue, thick, ->] (0.75,0.2) arc (90:330:0.5);
    \node at (0.75,0.5) {\textbf{\large #3}};
    \node[align=center, text width=2cm] at (0.75,-0.6) {#4};
  \end{scope}
}
\drawTriangle{0}{0}{1}{PET\\Bouteilles}
\drawTriangle{2.2}{0}{2}{PEHD\\Bidons}
\drawTriangle{4.4}{0}{3}{PVC\\Tuyaux}
\drawTriangle{6.6}{0}{4}{PEBD\\Sachets}
\drawTriangle{8.8}{0}{5}{PP\\Chaises}
\drawTriangle{11}{0}{6}{PS\\Gobelets}
\drawTriangle{13.2}{0}{7}{Autres\\Multicouches}
\end{tikzpicture}
\legende{Les 7 codes de recyclage (anneau de Möbius) : chiffre central = famille de polymère. Seuls les codes 1, 2, 4, 5 sont couramment recyclés à Douala.}
\end{popfigure}

\begin{popfigure}
\begin{tikzpicture}[scale=1.2, every node/.style={font=\small}]
% Polymer chains schematic
% Thermoplastic: linear chains
\begin{scope}[shift={(0,0)}]
\draw[popdark, thick] (-2,1.5) rectangle (2,2.5);
\node[popdark] at (0,2.1) {\textbf{Thermoplastique (linéaire / ramifié)}};
\foreach \y in {1.0, 0.5, 0.0, -0.5} {
  \draw[popblue, thick] (-1.8,\y) -- (1.8,\y);
  \foreach \x in {-1.5,-1.0,-0.5,0,0.5,1.0,1.5} {
    \fill[popblue!60] (\x,\y) circle (0.08);
  }
}
\draw[poporange, <->, thick] (-2.3,1.0) -- (-2.3,-0.5) node[midway, left, rotate=90] {Forces de Van der Waals (faibles)};
\node[align=center] at (0,-1.2) {Chauffage : chaînes glissent $\rightarrow$ ramollit\\Refroidissement : chaînes se figent $\rightarrow$ durcit};
\end{scope}
% Thermosetting: crosslinked network
\begin{scope}[shift={(5,0)}]
\draw[popdark, thick] (3,1.5) rectangle (7,2.5);
\node[popdark] at (5,2.1) {\textbf{Thermodurcissable (réticulé)}};
\foreach \y in {1.0, 0.5, 0.0, -0.5} {
  \foreach \x in {3.2,3.7,4.2,4.7,5.2,5.7,6.2,6.7} {
    \fill[poporange!80] (\x,\y) circle (0.08);
  }
}
% Crosslinks
\draw[poppurple, thick] (3.2,1.0) -- (3.7,0.5) -- (3.2,0.0) -- (3.7,-0.5);
\draw[poppurple, thick] (4.2,1.0) -- (4.7,0.5) -- (4.2,0.0) -- (4.7,-0.5);
\draw[poppurple, thick] (5.2,1.0) -- (5.7,0.5) -- (5.2,0.0) -- (5.7,-0.5);
\draw[poppurple, thick] (6.2,1.0) -- (6.7,0.5) -- (6.2,0.0) -- (6.7,-0.5);
\draw[poppurple, thick] (3.7,1.0) -- (4.2,0.5) -- (3.7,0.0);
\draw[poppurple, thick] (4.7,1.0) -- (5.2,0.5) -- (4.7,0.0);
\draw[poppurple, thick] (5.7,1.0) -- (6.2,0.5) -- (5.7,0.0);
\node[poppurple, align=center] at (5,1.8) {Ponts covalents (forts)};
\node[align=center] at (5,-1.2) {Chauffage : réseau rigide $\rightarrow$ pas de fusion\\Décomposition $\rightarrow$ carbonisation};
\end{scope}
\end{tikzpicture}
\legende{Structure moléculaire : thermoplastique (chaînes libres, forces faibles) vs thermodurcissable (réseau 3D réticulé, liaisons covalentes fortes).}
\end{popfigure}

\courssub{Exercices d'entraînement (type BEPC)}

\begin{exercice}[Identification et tri]
On ramasse dans la cour de l'école cinq objets : une bouteille d'eau vide (code 1), un sac de courses (code 4), une vieille sandale en plastique (code 7), un pot de fleur cassé (code 5), une poignée de casserole (sans code, noire, dure).
\begin{enumerate}
\item Donne le nom complet du polymère pour les trois objets codés.
\item Classe ces cinq objets en thermoplastiques et thermodurcissables. Justifie pour la poignée de casserole.
\item Seuls les codes 1, 2, 4, 5 sont acceptés par le recycleur local. Quels objets mets-tu dans le bac jaune ? Justifie le refus des autres.
\end{enumerate}
\end{exercice}

\begin{exercice}[Calcul et impact environnemental]
Une école de \num{800} élèves jette en moyenne \num{400} emballages plastiques par jour. L'année scolaire compte \num{180} jours.
\begin{enumerate}
\item Calcule le nombre total d'emballages jetés par an.
\item Sachant qu'un sachet PEBD met \num{200} ans à se décomposer, explique pourquoi l'expression « jeter » est trompeuse.
\item Relie la présence de ces déchets dans les caniveaux à l'apparition de maladies vectorielles (paludisme, dengue).
\end{enumerate}
\end{exercice}

\begin{exercice}[Argumentation 3R]
Le principal propose d'acheter une machine à broyer le plastique pour recycler sur place. Un élève répond : « Mieux vaut interdire les sachets à la cantine. »
\begin{enumerate}
\item Quelle règle des 3R illustre la proposition de l'élève ? Celle du principal ?
\item Pourquoi la proposition de l'élève est-elle scientifiquement plus pertinente (énergie, pollution, efficacité) ?
\item Donne un exemple concret de « Réutiliser » pour une bouteille PET dans l'école.
\end{enumerate}
\end{exercice}

\courssub{Corrigés des exercices}

\begin{corrige}[Exercice 1 : Identification et tri]
\begin{enumerate}
\item Bouteille d'eau (code 1) : \textbf{PET} (Polyéthylène téréphtalate). Sac courses (code 4) : \textbf{PEBD} (Polyéthylène basse densité). Pot de fleur (code 5) : \textbf{PP} (Polypropylène).
\item \textbf{Thermoplastiques :} bouteille (1), sac (4), pot (5), sandale (7 — souvent PE/PP/EVA). \textbf{Thermodurcissable :} poignée de casserole (bakélite/résine phénolique). \textit{Justification :} la poignée doit résister à la chaleur de la casserole (\SI{>100}{\celsius}) sans ramollir ; un thermoplastique fondrait et provoquerait des brûlures.
\item \textbf{Bac jaune (recyclables) :} bouteille (code 1), sac (code 4), pot (code 5). \textbf{Refusés :} sandale (code 7 = mélange multicouche ou caoutchouc/plastique, non triable), poignée (thermodurcissable, ne fond pas).
\end{enumerate}
\end{corrige}

\begin{corrige}[Exercice 2 : Calcul et impact]
\begin{enumerate}
\item $\num{400} \times \num{180} = \num{72000}$ emballages par an.
\item « Jeter » suggère une disparition. En réalité, le sachet persiste \num{200} ans : il fragmente en microplastiques, pollue le sol, l'eau, entre dans la chaîne alimentaire. Il n'y a pas de « dehors » : la Terre est un système fermé pour la matière.
\item Déchets dans caniveaux $\rightarrow$ bouchage $\rightarrow$ eau stagnante $\rightarrow$ gîte larvaire moustiques (\emph{Anopheles} paludisme, \emph{Aedes} dengue) $\rightarrow$ piqûres $\rightarrow$ transmission pathogènes $\rightarrow$ épidémies. Le déchet plastique devient déterminant environnemental de santé.
\end{enumerate}
\end{corrige}

\begin{corrige}[Exercice 3 : Argumentation 3R]
\begin{enumerate}
\item Élève : \textbf{Réduire} (agir à la source, ne pas produire le déchet). Principal : \textbf{Recycler} (traiter le déchet déjà produit).
\item Réduire est prioritaire car : (1) Zéro déchet = zéro énergie pour le collecter, transporter, broyer, laver, fondre, regranuler. (2) Zéro risque de fuite dans la nature. (3) Le recyclage dégrade le polymère (downcycling) : une bouteille ne redevient pas bouteille indéfiniment. (4) Au Cameroun, l'électricité pour la machine (groupe électrogène ou ENEO) a un coût carbone et financier.
\item Exemple Réutiliser : bouteille PET coupée $\rightarrow$ entonnoir pour verser l'huile moteur / arrosoir (bouchon percé) / pot de pépinière (fond percé) / brique écologique (remplie de sable tassé) pour bancs ou murets.
\end{enumerate}
\end{corrige}

\newpage

\coursec{Méthodes et exercices résolus}

\courssub{Méthode 1 : Identifier et définir une matière plastique}

\begin{methode}[Reconnaître et définir une matière plastique]
Pour répondre à une question du type « Qu'est-ce qu'une matière plastique ? », suis ces étapes :
\begin{enumerate}
\item \textbf{Définir} : une matière plastique est un matériau composé d'un \cle{polymère} (très grande molécule formée par l'enchaînement de petites molécules appelées \cle{monomères}) auquel on ajoute souvent des additifs (colorants, plastifiants, charges).
\item \textbf{Préciser l'origine} : la plupart des plastiques sont fabriqués à partir du \cle{pétrole} ou du gaz naturel, après raffinage et transformation chimique (polymérisation).
\item \textbf{Citer la propriété clé} : un plastique peut être \cle{moulé} sous l'action de la chaleur et de la pression.
\item \textbf{Donner un exemple concret} : sachet, bouteille d'eau minérale, cuvette, chaise en plastique.
\end{enumerate}
\textbf{Réflexe} : une définition complète contient toujours trois éléments : la nature (polymère), l'origine (pétrole le plus souvent) et la propriété (moulage à chaud).
\end{methode}

\begin{exoresolu}[Définir une matière plastique]
Énoncé : Au marché de Mokolo, Amina achète des légumes que le vendeur emballe dans un sachet plastique. Son petit frère lui demande : « C'est quoi exactement le plastique ? » Définis une matière plastique en précisant sa nature chimique, son origine et sa propriété caractéristique.
\tcblower
Solution détaillée :
\begin{itemize}
\item \textbf{Nature chimique} : une matière plastique est un matériau constitué essentiellement d'un \cle{polymère}, c'est-à-dire une très grande molécule (macromolécule) obtenue par l'enchaînement de centaines ou de milliers de petites molécules identiques appelées monomères. On y ajoute des additifs pour lui donner couleur, souplesse ou rigidité.
\item \textbf{Origine} : les plastiques sont fabriqués industriellement à partir du \cle{pétrole} (ou du gaz naturel). Au Cameroun, le pétrole raffiné à la SONARA fournit des composés de base qui servent à fabriquer les monomères.
\item \textbf{Propriété caractéristique} : un plastique se \cle{déforme} et se \cle{moule} sous l'action de la chaleur et de la pression, ce qui permet de lui donner toutes les formes (sachets, bouteilles, chaises).
\end{itemize}
\textbf{Conclusion} : une matière plastique est un matériau à base de polymère, généralement issu du pétrole, qui peut être moulé sous l'action de la chaleur et de la pression. Le sachet d'Amina est donc un objet en matière plastique.
\end{exoresolu}

\courssub{Méthode 2 : Classer un plastique (thermoplastique ou thermodurcissable)}

\begin{methode}[Distinguer thermoplastique et thermodurcissable]
Pour classer un plastique, raisonne sur son \textbf{comportement à la chaleur} :
\begin{enumerate}
\item \textbf{Chauffe (mentalement ou en TP)} l'objet en plastique.
\item \textbf{Observe} :
\begin{itemize}
\item s'il \textbf{ramollit et fond} à la chaleur, puis durcit en refroidissant (transformation \cle{réversible}) : c'est un \cle{thermoplastique} ;
\item s'il \textbf{ne fond pas} mais se décompose ou brûle sans ramollir (transformation irréversible) : c'est un \cle{thermodurcissable}.
\end{itemize}
\item \textbf{Conclus sur le recyclage} : seuls les thermoplastiques peuvent être refondus et \cle{recyclés} facilement.
\end{enumerate}
\textbf{Exemples à retenir} : thermoplastiques = PE, PP, PVC, PS, PET ; thermodurcissables = bakélite (poignées de marmite), résines, colles époxy.
\end{methode}

\begin{exoresolu}[Classer deux objets en plastique]
Énoncé : Dans la cuisine de maman, on trouve un seau en plastique et la poignée noire de la marmite. En approchant (avec précaution) une flamme du seau, celui-ci ramollit et se déforme. La poignée de la marmite, elle, noircit et brûle sans fondre.
\begin{enumerate}
\item À quelle famille de plastiques appartient chacun de ces objets ? Justifie.
\item Lequel de ces deux objets peut être recyclé par fusion ? Explique pourquoi.
\end{enumerate}
\tcblower
Solution détaillée :
\begin{enumerate}
\item \textbf{Le seau} ramollit et se déforme à la chaleur : sa transformation est réversible (il durcirait en refroidissant). C'est donc un \cle{thermoplastique} (par exemple du polyéthylène PE ou du polypropylène PP). \textbf{La poignée de marmite} ne fond pas : elle noircit et brûle, sa transformation est irréversible. C'est un \cle{thermodurcissable} (par exemple de la bakélite), dont les chaînes de polymère sont liées entre elles en réseau, ce qui empêche la fusion.
\item \textbf{Le seau (thermoplastique)} peut être recyclé par fusion : on peut le refondre et le remouler pour fabriquer un nouvel objet. La poignée thermodurcissable ne peut pas être refondue : la chauffer davantage la détruit au lieu de la ramollir. Elle n'est donc pas recyclable par fusion.
\end{enumerate}
\textbf{Conclusion} : le seau est un thermoplastique recyclable ; la poignée de marmite est un thermodurcissable non recyclable par fusion.
\end{exoresolu}

\courssub{Méthode 3 : Identifier un plastique courant et choisir son usage}

\begin{methode}[Associer un plastique à son usage]
Pour identifier un plastique courant et justifier son emploi :
\begin{enumerate}
\item \textbf{Repère les indices} : aspect (souple, rigide, transparent), code de recyclage éventuel (triangle avec un chiffre : 1 = PET, 2 = PEHD, 3 = PVC, 5 = PP).
\item \textbf{Rappelle les propriétés} du plastique : légèreté, imperméabilité, transparence, résistance, isolation électrique.
\item \textbf{Relie propriété et usage} : c'est la propriété qui justifie l'usage (ex. PVC isolant $\Rightarrow$ gaine des fils électriques).
\end{enumerate}
\textbf{Couples à retenir} : PET $\Rightarrow$ bouteilles d'eau ; PE $\Rightarrow$ sachets et bidons ; PVC $\Rightarrow$ tuyaux et gaines électriques ; PP $\Rightarrow$ chaises et cuvettes ; PS $\Rightarrow$ emballages légers.
\end{methode}

\begin{exoresolu}[Choisir le bon plastique pour un usage]
Énoncé : Un électricien de Douala doit installer des fils électriques dans une maison. Il dispose de trois plastiques : le PET (transparent, utilisé pour les bouteilles), le PVC (souple, bon isolant électrique) et le polystyrène (très léger et cassant).
\begin{enumerate}
\item Quel plastique doit-il choisir pour la gaine de protection des fils ? Justifie par deux propriétés.
\item Pourquoi le polystyrène est-il inadapté à cet usage ?
\end{enumerate}
\tcblower
Solution détaillée :
\begin{enumerate}
\item L'électricien doit choisir le \cle{PVC}. Justification : (a) le PVC est un excellent \cle{isolant électrique}, il empêche le passage du courant et protège contre les électrocutions ; (b) il est \cle{souple} et résistant, il peut donc suivre les coudes des installations sans se fissurer, et il dure longtemps.
\item Le polystyrène est inadapté car il est \cle{cassant} : il se fissurerait lors de l'installation, laissant les fils nus, ce qui créerait un danger d'électrocution et d'incendie. Sa légèreté n'apporte aucun avantage pour une gaine.
\end{enumerate}
\textbf{Conclusion} : c'est le PVC qui convient pour les gaines électriques, grâce à ses propriétés d'isolation électrique et de souplesse. On retient la règle : \emph{l'usage d'un plastique dépend de ses propriétés}.
\end{exoresolu}

\courssub{Méthode 4 : Réaliser et interpréter un test simple d'identification}

\begin{methode}[Interpréter un test de combustion ou de chauffage]
En TP, on peut distinguer des plastiques par un test de chauffage (réalisé par le professeur, avec les consignes de sécurité) :
\begin{enumerate}
\item \textbf{Consignes de sécurité d'abord} : blouse, lunettes, hotte ou espace aéré, pince pour tenir l'échantillon, ne jamais respirer les fumées (elles sont \cle{toxiques}).
\item \textbf{Chauffer} un petit échantillon et \textbf{observer} : ramollit-il ? fond-il ? brûle-t-il ? quelle odeur, quelle fumée ?
\item \textbf{Consigner les observations} dans un tableau : matériau, comportement à la chaleur, conclusion (famille).
\item \textbf{Conclure} en reliant l'observation à la classification : fusion réversible $\Rightarrow$ thermoplastique ; pas de fusion $\Rightarrow$ thermodurcissable.
\end{enumerate}
\textbf{Réflexe rédaction} : au BEPC, rédige toujours en trois temps : observation $\rightarrow$ interprétation $\rightarrow$ conclusion.
\end{methode}

\begin{exoresolu}[TP : identifier deux plastiques inconnus]
Énoncé : Au laboratoire du collège, le professeur chauffe doucement deux échantillons de plastique A et B, tenus avec une pince au-dessus d'une flamme. Observations : l'échantillon A ramollit, fond et coule en gouttelettes ; l'échantillon B noircit, dégage une fumée âcre, mais ne fond pas.
\begin{enumerate}
\item Présente les résultats sous forme d'un tableau.
\item Identifie la famille de chaque échantillon en justifiant.
\item Cite deux consignes de sécurité respectées pendant ce test.
\end{enumerate}
\tcblower
Solution détaillée :
\begin{enumerate}
\item Tableau des résultats :
\begin{center}
\begin{tabular}{|l|l|l|}
\hline
\rowcolor{popblueL} Échantillon & Observation à la chaleur & Famille \\
\hline
A & ramollit, fond, coule en gouttes & thermoplastique \\
\hline
B & noircit, fumée âcre, ne fond pas & thermodurcissable \\
\hline
\end{tabular}
\end{center}
\item \textbf{Échantillon A} : il fond à la chaleur, sa transformation est réversible (en refroidissant, il redeviendrait solide). C'est donc un \cle{thermoplastique}. \textbf{Échantillon B} : il ne fond pas mais se décompose en brûlant, sa transformation est irréversible. C'est donc un \cle{thermodurcissable}, dont les chaînes de polymère forment un réseau rigide qui ne peut pas fondre.
\item Consignes de sécurité : tenir l'échantillon avec une \cle{pince} (et non à la main) pour éviter les brûlures ; ne pas \cle{respirer les fumées} dégagées car elles sont toxiques (travailler dans un espace aéré ou sous la hotte) ; porter blouse et lunettes.
\end{enumerate}
\textbf{Conclusion} : le test de chauffage permet de distinguer un thermoplastique (A, qui fond) d'un thermodurcissable (B, qui brûle sans fondre), à condition de respecter les consignes de sécurité.
\end{exoresolu}

\courssub{Méthode 5 : Argumenter sur la pollution plastique et proposer des solutions}

\begin{methode}[Traiter une question sur la pollution et la gestion des plastiques]
Pour une question ouverte du type « Explique pourquoi les plastiques polluent et propose des solutions » :
\begin{enumerate}
\item \textbf{Explique la cause} : les plastiques sont \cle{non biodégradables} (ou très lentement dégradables : des centaines d'années) ; ils s'accumulent dans les sols, les rues, les rivières et l'océan.
\item \textbf{Cite les conséquences} : bouchage des caniveaux et inondations (ex. quartiers de Yaoundé et Douala en saison des pluies), étouffement des animaux, pollution des sols et de l'eau, fumées toxiques lors des brûlages sauvages.
\item \textbf{Propose les solutions des 3R} : \cle{Réduire} la consommation (refuser les sachets inutiles), \cle{Réutiliser} (bidons, sacs durables), \cle{Recycler} (collecte, tri, fonte des thermoplastiques).
\item \textbf{Conclus} par un comportement citoyen concret.
\end{enumerate}
\textbf{Réflexe} : une bonne réponse cite toujours au moins une cause, une conséquence et deux solutions.
\end{methode}

\begin{exoresolu}[La pollution plastique dans mon quartier]
Énoncé : Après chaque forte pluie, les caniveaux du quartier Briqueterie à Yaoundé débordent et l'eau envahit les maisons. On constate que les caniveaux sont remplis de sachets plastiques et de bouteilles vides.
\begin{enumerate}
\item Explique pourquoi les déchets plastiques provoquent ces inondations.
\item Pourquoi ces plastiques ne disparaissent-ils pas d'eux-mêmes avec le temps ?
\item Propose trois actions concrètes, inspirées de la règle des 3R, que les habitants peuvent mener.
\end{enumerate}
\tcblower
Solution détaillée :
\begin{enumerate}
\item Les sachets et bouteilles plastiques jetés dans la rue sont entraînés par l'eau de pluie vers les caniveaux. Ils s'y accumulent et \cle{bouchent} les canalisations : l'eau ne peut plus s'écouler, elle monte et envahit les maisons. C'est l'accumulation des déchets plastiques qui provoque les inondations.
\item Les plastiques sont des polymères \cle{non biodégradables} : les micro-organismes du sol ne parviennent pas à les décomposer. Un sachet plastique peut mettre des \cle{centaines d'années} à disparaître. Ils s'accumulent donc au lieu de disparaître.
\item Actions concrètes (règle des 3R) :
\begin{itemize}
\item \textbf{Réduire} : refuser les sachets plastiques inutiles au marché et utiliser un sac en tissu ou un panier ;
\item \textbf{Réutiliser} : se servir plusieurs fois des bouteilles et des sacs solides au lieu de les jeter après un seul usage ;
\item \textbf{Recycler} : trier les bouteilles et sachets, les vendre ou les remettre aux filières de recyclage qui les fondent pour fabriquer de nouveaux objets (pavés, arrosoirs, chaises).
\end{itemize}
\end{enumerate}
\textbf{Conclusion} : les plastiques, non biodégradables, bouchent les caniveaux et provoquent des inondations ; la solution passe par la règle des 3R : Réduire, Réutiliser, Recycler.
\end{exoresolu}

\courssub{Méthode 6 : Exploiter des données chiffrées sur les plastiques}

\begin{methode}[Calculer avec des données sur les déchets plastiques]
Pour un exercice chiffré (durées de dégradation, masses, pourcentages) :
\begin{enumerate}
\item \textbf{Repère les données} et note-les avec leurs unités.
\item \textbf{Identifie la grandeur cherchée} et choisis l'opération (multiplication, proportionnalité, pourcentage).
\item \textbf{Pose le calcul} en ligne, avec les unités à chaque étape.
\item \textbf{Conclus par une phrase} qui redonne du sens au résultat (ex. durée de vie d'un plastique, masse recyclée).
\end{enumerate}
\textbf{Formule utile} : pour un pourcentage, $\text{partie} = \text{total} \times \dfrac{t}{100}$.
\end{methode}

\begin{exoresolu}[Calcul sur le recyclage des bouteilles]
Énoncé : Un quartier de Douala produit chaque mois \SI{2500}{kg} de déchets plastiques. Une coopérative de recyclage parvient à collecter et recycler $40\,\%$ de cette masse.
\begin{enumerate}
\item Calcule la masse de plastique recyclée chaque mois.
\item Calcule la masse de plastique qui reste encore abandonnée dans la nature chaque mois.
\item Sachant qu'une bouteille plastique met environ 450 ans à se dégrader, explique en une phrase pourquoi le recyclage est indispensable.
\end{enumerate}
\tcblower
Solution détaillée :
\begin{enumerate}
\item Masse recyclée : on applique le pourcentage à la masse totale :
\[ m_{\text{recyclée}} = 2500 \times \frac{40}{100} = 2500 \times 0{,}40 = \SI{1000}{kg}. \]
La coopérative recycle donc \SI{1000}{kg} de plastique par mois.
\item Masse abandonnée : c'est le reste de la masse totale :
\[ m_{\text{abandonnée}} = 2500 - 1000 = \SI{1500}{kg}. \]
Il reste donc \SI{1500}{kg} de plastique abandonnés dans la nature chaque mois.
\item Une bouteille plastique met environ 450 ans à se dégrader : si on ne la recycle pas, elle pollue l'environnement pendant plus de quatre siècles, soit plusieurs générations. Le recyclage est donc indispensable pour éviter l'accumulation de ces déchets presque éternels.
\end{enumerate}
\textbf{Conclusion} : sur \SI{2500}{kg} de déchets plastiques mensuels, \SI{1000}{kg} sont recyclés et \SI{1500}{kg} restent dans la nature ; vu la durée de dégradation des plastiques, il faut augmenter la part recyclée.
\end{exoresolu}

\begin{attention}[Les 5 erreurs qui coûtent des points au Bac]
\begin{enumerate}
\item \textbf{Confondre thermoplastique et thermodurcissable} : c'est le comportement à la chaleur qui tranche. Thermoplastique = \emph{fond} à la chaleur (recyclable) ; thermodurcissable = \emph{brûle sans fondre} (non recyclable par fusion). Ne jamais inverser.
\item \textbf{Définir le plastique de façon incomplète} : écrire « c'est un matériau léger » n'est pas une définition. Exige-toi les trois éléments : polymère + origine (pétrole) + moulage à chaud.
\item \textbf{Oublier les consignes de sécurité} dans un test de combustion : ne jamais tenir l'échantillon à la main (pince obligatoire) et ne jamais respirer les fumées, qui sont toxiques. Une question sur la sécurité est presque toujours posée.
\item \textbf{Affirmer que « le plastique se décompose vite »} : c'est faux. Les plastiques sont non biodégradables et mettent des centaines d'années à disparaître ; c'est précisément la cause de la pollution.
\item \textbf{Erreurs de calcul sur les pourcentages} : pour trouver $40\,\%$ de \SI{2500}{kg}, on calcule $2500 \times \dfrac{40}{100}$, et non $2500 \times 40$. Et n'oublie jamais l'\cle{unité} dans le résultat final : une masse sans unité fait perdre des points.
\end{enumerate}
\end{attention}

\newpage

\coursec{Exercices}

\courssub{Je vérifie mes connaissances}

\begin{exercice}[QCM : à chaque objet son plastique]
Pour chaque objet ci-dessous, choisis la bonne réponse et justifie ton choix par une propriété du plastique.
\begin{enumerate}
\item Une bouteille d'eau minérale transparente est fabriquée en : \\
\textbf{a)} PVC \quad \textbf{b)} PET \quad \textbf{c)} PS \quad \textbf{d)} bakélite
\item Un tuyau d'arrosage rigide est fabriqué en : \\
\textbf{a)} PET \quad \textbf{b)} PS \quad \textbf{c)} PVC \quad \textbf{d)} PE expansé
\item Une chaise plastique de salon est fabriquée en : \\
\textbf{a)} PP \quad \textbf{b)} PET \quad \textbf{c)} PS \quad \textbf{d)} PVC souple
\item Un sachet d'emballage fin et souple est fabriqué en : \\
\textbf{a)} PS \quad \textbf{b)} PE \quad \textbf{c)} bakélite \quad \textbf{d)} PET
\item La prise électrique de ta maison, dure et infusible, est en : \\
\textbf{a)} thermoplastique \quad \textbf{b)} thermodurcissable \quad \textbf{c)} PE \quad \textbf{d)} PET
\end{enumerate}
\end{exercice}

\begin{exercice}[Vrai ou faux ? Justifie ta réponse]
Réponds par \textbf{vrai} ou \textbf{faux}, puis justifie chaque réponse par une phrase.
\begin{enumerate}
\item Toutes les matières plastiques fondent quand on les chauffe.
\item Un thermodurcissable peut être recyclé par simple refonte.
\item Le PET des bouteilles vides peut être recyclé pour fabriquer des fibres textiles.
\item Les matières plastiques sont des polymères, c'est-à-dire de très grosses molécules formées par la répétition d'un motif appelé monomère.
\item Brûler les sachets plastiques à l'air libre est une bonne méthode de gestion des déchets.
\end{enumerate}
\end{exercice}

\begin{exercice}[Définitions à compléter]
Recopie et complète les phrases suivantes avec les mots ou expressions de la liste : \emph{monomère ; polymère ; thermoplastique ; thermodurcissable ; polymérisation ; recyclage}.
\begin{enumerate}
\item Une matière plastique est un \ldots\ldots\ldots\ldots, c'est-à-dire une macromolécule obtenue par l'assemblage de nombreuses petites molécules appelées \ldots\ldots\ldots\ldots.
\item La réaction chimique qui assemble les monomères s'appelle la \ldots\ldots\ldots\ldots.
\item Un plastique \ldots\ldots\ldots\ldots se ramollit à la chaleur et peut être remodelé : il est recyclable par refonte.
\item Un plastique \ldots\ldots\ldots\ldots durcit définitivement lors de sa fabrication et se décompose sans fondre quand on le chauffe.
\item Le \ldots\ldots\ldots\ldots consiste à transformer un déchet plastique en nouvelle matière première.
\end{enumerate}
\end{exercice}

\begin{exercice}[Petit calcul : la masse de sachets]
Un quartier de Yaoundé consomme en moyenne \num{2500} sachets plastiques par jour. Un sachet a une masse moyenne de \SI{4}{g}.
\begin{enumerate}
\item Calcule la masse totale de sachets utilisés chaque jour dans ce quartier, en kilogrammes.
\item Calcule la masse de sachets accumulée en une année ($365$ jours), en tonnes.
\item On estime que seulement $10\,\%$ de ces sachets sont collectés et recyclés. Quelle masse de plastique reste abandonnée dans la nature chaque année ?
\end{enumerate}
\end{exercice}

\courssub{Je m'entraîne}

\begin{exercice}[Expérience à interpréter : le test de chauffe]
Au laboratoire, on chauffe doucement deux échantillons de plastiques inconnus, A et B, dans deux tubes à essai.
\begin{itemize}
\item Échantillon A : il se ramollit puis fond vers \SI{110}{\celsius} ; en refroidissant, il durcit à nouveau.
\item Échantillon B : il ne fond pas ; il noircit et se décompose en dégageant une odeur âcre.
\end{itemize}
\begin{enumerate}
\item À quelle famille de plastiques (thermoplastique ou thermodurcissable) appartient chaque échantillon ? Justifie.
\item Lequel des deux échantillons peut être recyclé par refonte ? Explique pourquoi.
\item Cite un objet de la vie courante qui pourrait être fabriqué avec le plastique A, puis un objet fabriqué avec le plastique B.
\item Pourquoi ne faut-il jamais réaliser ce test de chauffe sans hotte ni ventilation ?
\end{enumerate}
\end{exercice}

\begin{exercice}[Durées de décomposition : construis un diagramme]
Le tableau ci-dessous donne la durée estimée de décomposition de quelques déchets dans la nature.

\begin{center}
\begin{tabularx}{\linewidth}{|l|X|}
\hline
\rowcolor{popblueL} \textbf{Déchet} & \textbf{Durée de décomposition estimée} \\
\hline
Pelure de banane & 2 mois \\
\hline
Mouchoir en papier & 3 mois \\
\hline
Sachet plastique (PE) & 100 à 400 ans \\
\hline
Bouteille plastique (PET) & 100 à 1000 ans \\
\hline
Canette en aluminium & 200 à 500 ans \\
\hline
\end{tabularx}
\end{center}

\begin{enumerate}
\item Construis un diagramme en barres représentant la durée de décomposition (en années, en prenant la valeur maximale de chaque intervalle) en fonction du type de déchet. Échelle conseillée : \SI{1}{cm} pour 100 ans.
\item Compare la durée de vie d'un sachet plastique à celle d'un mouchoir en papier.
\item Déduis de ce diagramme deux conclusions sur la pollution plastique.
\item Propose un geste simple qu'un élève de 3ème peut adopter pour réduire cette pollution.
\end{enumerate}
\end{exercice}

\begin{exercice}[Étude de document : l'interdiction des sachets au Cameroun]
\textbf{Document.} Depuis 2014, le Cameroun a interdit la fabrication, l'importation et la commercialisation des sachets plastiques non biodégradables de faible épaisseur. Ces sachets, distribués gratuitement dans les marchés, bouchent les caniveaux, provoquent des inondations en saison des pluies, étouffent le bétail qui les avale et polluent les sols pendant des siècles. Des alternatives existent : sachets biodégradables, emballages en papier, paniers et sacs réutilisables.

\begin{enumerate}
\item Dégage du texte trois problèmes causés par les sachets plastiques non biodégradables.
\item Explique, à l'aide de tes connaissances, pourquoi un sachet en polyéthylène persiste longtemps dans la nature.
\item Cite deux solutions de remplacement proposées dans le texte et classe-les selon les « 3R » (Réduire, Réutiliser, Recycler).
\item Rédige un slogan de sensibilisation (une phrase) pour convaincre les commerçants de ton quartier de respecter cette interdiction.
\end{enumerate}
\end{exercice}

\begin{exercice}[Le tri des déchets de la famille Essomba]
Pendant une semaine, la famille Essomba, à Douala, a pesé ses déchets plastiques :

\begin{center}
\begin{tabularx}{\linewidth}{|l|X|X|}
\hline
\rowcolor{popblueL} \textbf{Déchet} & \textbf{Nature du plastique} & \textbf{Masse} \\
\hline
Bouteilles d'eau vides & PET (thermoplastique) & \SI{1,8}{kg} \\
\hline
Sachets d'emballage & PE (thermoplastique) & \SI{0,9}{kg} \\
\hline
Vieille prise électrique cassée & Bakélite (thermodurcissable) & \SI{0,3}{kg} \\
\hline
Cuvettes et calebasses plastiques & PP (thermoplastique) & \SI{1,0}{kg} \\
\hline
\end{tabularx}
\end{center}

\begin{enumerate}
\item Calcule la masse totale de déchets plastiques produite par la famille en une semaine.
\item Calcule la masse de plastiques recyclables par refonte. Justifie ta réponse.
\item Quel pourcentage de la masse totale représente la bakélite ? Pourquoi ne peut-elle pas être fondue pour être recyclée ?
\item Propose un circuit de valorisation réaliste pour les bouteilles en PET (de la poubelle au produit fini), en citant un acteur possible (collecteur, entreprise de recyclage, artisan).
\end{enumerate}
\end{exercice}

\courssub{Je me prépare au Bac}

\begin{exercice}[Sujet type BEPC n° 1 : les matières plastiques et leur classification]
\textbf{Partie A — Connaissances.}
\begin{enumerate}
\item Définis les termes suivants : \emph{monomère}, \emph{polymère}, \emph{matière plastique}.
\item L'éthylène a pour formule \ce{CH2=CH2}. Écris le motif (monomère) du polyéthylène et représente une chaîne contenant trois motifs.
\item Définis les termes \emph{thermoplastique} et \emph{thermodurcissable}, et donne un exemple de plastique pour chaque famille.
\end{enumerate}
\textbf{Partie B — Expérience.} On dispose de trois échantillons notés X, Y et Z. Chauffés modérément, X fond à \SI{130}{\celsius}, Y fond à \SI{250}{\celsius}, Z noircit sans fondre.
\begin{enumerate}
\setcounter{enumi}{3}
\item Classe les échantillons X, Y et Z dans les familles de plastiques étudiées. Justifie.
\item Parmi ces trois échantillons, lequel conviendrait à la fabrication d'une poignée de casserole ? Justifie.
\item Lesquels peuvent être recyclés par refonte ? Décris en deux étapes le procédé de recyclage par refonte.
\end{enumerate}
\textbf{Partie C — Citoyenneté.}
\begin{enumerate}
\setcounter{enumi}{6}
\item Explique en trois phrases maximum pourquoi la gestion des déchets plastiques est un enjeu de santé publique dans les villes camerounaises.
\end{enumerate}
\end{exercice}

\begin{exercice}[Sujet type BEPC n° 2 : les plastiques dans la vie quotidienne]
\textbf{Partie A — Identification.} Voici une liste d'objets : bouteille de boisson gazeuse, tuyau de plomberie, gobelet jetable, chaise plastique, sachet de marché, interrupteur électrique.
\begin{enumerate}
\item Recopie et complète le tableau suivant en identifiant pour chaque objet le plastique utilisé (PET, PVC, PS, PP, PE ou thermodurcissable) et la propriété qui justifie ce choix (transparence, rigidité, légèreté, souplesse, résistance à la chaleur\ldots).

\begin{center}
\begin{tabularx}{\linewidth}{|l|X|X|}
\hline
\rowcolor{popblueL} \textbf{Objet} & \textbf{Plastique} & \textbf{Propriété justifiant le choix} \\
\hline
Bouteille de boisson gazeuse & & \\
\hline
Tuyau de plomberie & & \\
\hline
Gobelet jetable & & \\
\hline
Chaise plastique & & \\
\hline
Sachet de marché & & \\
\hline
Interrupteur électrique & & \\
\hline
\end{tabularx}
\end{center}
\end{enumerate}
\textbf{Partie B — Avantages et limites.}
\begin{enumerate}
\setcounter{enumi}{1}
\item Cite trois avantages des matières plastiques qui expliquent leur succès dans la vie courante.
\item Cite deux limites (inconvénients) des matières plastiques pour l'environnement et la santé.
\end{enumerate}
\textbf{Partie C — Calcul.} Une entreprise de recyclage de Yaoundé collecte \SI{12}{tonnes} de bouteilles en PET par mois. Le procédé de recyclage permet de récupérer $85\,\%$ de la masse sous forme de granulés réutilisables.
\begin{enumerate}
\setcounter{enumi}{3}
\item Calcule la masse de granulés de PET obtenue chaque mois.
\item La fabrication de \SI{1}{kg} de PET neuf consomme \SI{2}{kWh} d'énergie de plus que la production de \SI{1}{kg} de PET recyclé. Calcule l'énergie économisée chaque mois grâce au recyclage.
\item Conclus en une phrase sur l'intérêt énergétique du recyclage.
\end{enumerate}
\end{exercice}

\begin{exercice}[Sujet type BEPC n° 3 : pollution plastique et gestion des déchets]
\textbf{Partie A — Étude de document.} Un reportage indique : « Chaque année, environ 8 millions de tonnes de plastiques finissent dans les océans du monde. Au Cameroun, les déchets plastiques abandonnés bouchent les caniveaux de Douala et de Yaoundé, favorisent les inondations et la prolifération des moustiques vecteurs du paludisme. »
\begin{enumerate}
\item Explique le lien entre les déchets plastiques, les inondations et le paludisme.
\item Pourquoi dit-on qu'une bouteille en PET abandonnée dans la nature est un polluant « durable » ? Appuie ta réponse sur un ordre de grandeur de durée de décomposition.
\end{enumerate}
\textbf{Partie B — Les 3R.}
\begin{enumerate}
\setcounter{enumi}{2}
\item Donne la signification des trois « R » de la gestion des déchets.
\item Pour chacun des gestes suivants, indique le « R » correspondant : (a) utiliser une gourde au lieu d'acheter de l'eau en bouteille ; (b) transformer les bouteilles vides en arrosoirs ; (c) fondre les bouchons en PE pour fabriquer des pavés ; (d) refuser un sachet plastique au marché quand on a son panier.
\end{enumerate}
\textbf{Partie C — Projet de quartier.} Ton quartier produit \SI{500}{kg} de déchets plastiques par semaine, dont $60\,\%$ de thermoplastiques recyclables.
\begin{enumerate}
\setcounter{enumi}{4}
\item Calcule la masse de plastiques recyclables produite par semaine, puis par mois (4 semaines).
\item Une coopérative de jeunes propose de collecter ces plastiques et de les vendre \SI{50}{FCFA} le kilogramme à une usine de recyclage. Calcule la recette mensuelle de la coopérative.
\item Propose deux actions concrètes, en plus de la collecte, pour que ce projet réduise durablement la pollution dans ton quartier.
\end{enumerate}
\end{exercice}

\newpage

\coursec{Corrigés des exercices}

\begin{corrige}[Exercice 1 — QCM : à chaque objet son plastique]
\begin{enumerate}
\item \textbf{b) PET.} Le PET est transparent, léger, imperméable aux gaz : il convient parfaitement aux bouteilles de boissons.
\item \textbf{c) PVC.} Le PVC est rigide, résistant à l'eau et aux produits chimiques : c'est le plastique des tuyaux d'arrosage et de plomberie.
\item \textbf{a) PP.} Le polypropylène est rigide, solide et résistant aux chocs : il sert à fabriquer les chaises, cuvettes et calebasses plastiques.
\item \textbf{b) PE.} Le polyéthylène est souple, fin et léger : c'est le plastique des sachets d'emballage.
\item \textbf{b) thermodurcissable.} Une prise électrique doit rester dure et ne pas fondre même en cas d'échauffement : on utilise un thermodurcissable (bakélite), infusible et isolant.
\end{enumerate}
\end{corrige}

\begin{corrige}[Exercice 2 — Vrai ou faux ?]
\begin{enumerate}
\item \textbf{Faux.} Seuls les thermoplastiques se ramollissent et fondent à la chaleur ; les thermodurcissables se décomposent (ils noircissent) sans fondre.
\item \textbf{Faux.} Un thermodurcissable ne peut pas être refondu : une fois durci, sa structure est définitive ; chauffé, il se décompose. Il n'est donc pas recyclable par refonte.
\item \textbf{Vrai.} Le PET est un thermoplastique : fondu puis étiré, il donne des fibres textiles (utilisées pour des polaires, des coussins, des sacs).
\item \textbf{Vrai.} C'est la définition même : une matière plastique est un polymère, macromolécule formée par la répétition de nombreuses petites molécules appelées monomères.
\item \textbf{Faux.} Brûler les plastiques à l'air libre dégage des fumées toxiques (gaz irritants et cancérigènes) dangereuses pour la santé et l'environnement. Il faut trier, réduire, réutiliser et recycler.
\end{enumerate}
\end{corrige}

\begin{corrige}[Exercice 3 — Définitions à compléter]
\begin{enumerate}
\item Une matière plastique est un \textbf{polymère}, c'est-à-dire une macromolécule obtenue par l'assemblage de nombreuses petites molécules appelées \textbf{monomères}.
\item La réaction chimique qui assemble les monomères s'appelle la \textbf{polymérisation}.
\item Un plastique \textbf{thermoplastique} se ramollit à la chaleur et peut être remodelé : il est recyclable par refonte.
\item Un plastique \textbf{thermodurcissable} durcit définitivement lors de sa fabrication et se décompose sans fondre quand on le chauffe.
\item Le \textbf{recyclage} consiste à transformer un déchet plastique en nouvelle matière première.
\end{enumerate}
\end{corrige}

\begin{corrige}[Exercice 4 — Petit calcul : la masse de sachets]
\begin{enumerate}
\item Masse journalière :
\[ m = \num{2500} \times \SI{4}{g} = \SI{10000}{g} = \SI{10}{kg}. \]
Le quartier utilise \textbf{\SI{10}{kg} de sachets par jour}.
\item Masse annuelle :
\[ m_{\text{an}} = 10 \times 365 = \SI{3650}{kg} = \SI{3,65}{tonnes}. \]
Le quartier accumule \textbf{\SI{3,65}{tonnes} de sachets par an}.
\item Seulement $10\,\%$ sont recyclés, donc $90\,\%$ sont abandonnés :
\[ m_{\text{abandonnée}} = \num{3,65} \times \frac{90}{100} = \SI{3,285}{tonnes}. \]
Environ \textbf{\SI{3,29}{tonnes}} de plastique restent dans la nature chaque année, rien que pour ce quartier.
\end{enumerate}
\end{corrige}

\begin{corrige}[Exercice 5 — Expérience à interpréter : le test de chauffe]
\begin{enumerate}
\item \textbf{Échantillon A : thermoplastique.} Il se ramollit, fond à \SI{110}{\celsius} puis durcit en refroidissant : ce comportement réversible est la signature des thermoplastiques. \textbf{Échantillon B : thermodurcissable.} Il ne fond pas mais noircit et se décompose : c'est le comportement des thermodurcissables.
\item Seul l'échantillon \textbf{A} peut être recyclé par refonte : sa fusion est réversible, on peut le fondre, le remettre en forme dans un moule, puis le laisser refroidir. L'échantillon B se détruirait au lieu de fondre.
\item Plastique A (thermoplastique) : une bouteille d'eau, un sachet, une cuvette. Plastique B (thermodurcissable) : une prise électrique, un interrupteur, une poignée de casserole.
\item La chauffe des plastiques dégage des \textbf{fumées toxiques} (gaz irritants, odeurs âcres). Sans hotte ni ventilation, on risque de les inhaler : il faut travailler sous hotte aspirante, avec des lunettes et des gants.
\end{enumerate}
\end{corrige}

\begin{corrige}[Exercice 6 — Durées de décomposition : construis un diagramme]
\begin{enumerate}
\item On prend les valeurs maximales : pelure de banane et mouchoir en papier (moins d'un an, barres quasi nulles), sachet PE : 400 ans, bouteille PET : \num{1000} ans, canette : 500 ans.

\begin{popfigure}
\begin{tikzpicture}
\begin{axis}[
    ybar, bar width=0.55cm,
    width=\linewidth, height=6.5cm,
    ylabel={Dur\'ee de d\'ecomposition (ann\'ees)},
    symbolic x coords={Banane, Papier, Sachet PE, Bouteille PET, Canette},
    xtick=data, x tick label style={font=\small, rotate=15, anchor=east},
    ymin=0, ymax=1100,
    ytick={0,200,400,600,800,1000},
    axis lines=left, grid=major, grid style={popline!50},
    every axis plot/.append style={fill=popblue, draw=popdark}
]
\addplot coordinates {(Banane,0.2) (Papier,0.25) (Sachet PE,400) (Bouteille PET,1000) (Canette,500)};
\end{axis}
\end{tikzpicture}
\legende{Diagramme en barres des durées de décomposition (valeurs maximales). Les déchets organiques disparaissent en quelques mois, les plastiques persistent des siècles.}
\end{popfigure}

\item Un sachet plastique met jusqu'à 400 ans à disparaître, contre environ 3 mois pour un mouchoir en papier : le sachet dure environ \textbf{\num{1600} fois plus longtemps} ($400 \times 12 / 3 = \num{1600}$).
\item Conclusions : (1) les plastiques sont des polluants très persistants, ils s'accumulent dans la nature pendant des siècles ; (2) contrairement aux déchets biodégradables, ils ne disparaissent pas seuls : seule leur collecte et leur recyclage peuvent limiter la pollution.
\item Geste simple : utiliser un sac ou un panier réutilisable pour les courses et refuser les sachets plastiques au marché (ou boire à la gourde plutôt qu'acheter des bouteilles jetables).
\end{enumerate}
\end{corrige}

\begin{corrige}[Exercice 7 — Étude de document : l'interdiction des sachets au Cameroun]
\begin{enumerate}
\item Trois problèmes cités dans le texte : les sachets \textbf{bouchent les caniveaux} et provoquent des \textbf{inondations} ; ils \textbf{étouffent le bétail} qui les avale ; ils \textbf{polluent les sols} pendant des siècles.
\item Un sachet en polyéthylène est un polymère à très longues chaînes carbonées, \textbf{non biodégradable} : les micro-organismes du sol ne savent pas le décomposer. Il persiste donc 100 à 400 ans dans la nature.
\item Solutions du texte : les \textbf{paniers et sacs réutilisables} relèvent de \textbf{Réutiliser} ; les \textbf{sachets biodégradables et emballages en papier} relèvent de \textbf{Réduire} (on remplace le plastique persistant par un matériau qui se décompose vite). Le \textbf{recyclage} (3\ieme{} R) consisterait à fondre les sachets collectés pour fabriquer de nouveaux objets.
\item Exemple de slogan : « Un panier à la main, un quartier propre et sans inondation : dis non au sachet plastique ! »
\end{enumerate}
\end{corrige}

\begin{corrige}[Exercice 8 — Le tri des déchets de la famille Essomba]
\begin{enumerate}
\item Masse totale :
\[ m_{\text{totale}} = 1,8 + 0,9 + 0,3 + 1,0 = \SI{4,0}{kg}. \]
\item Seuls les \textbf{thermoplastiques} sont recyclables par refonte : PET, PE et PP.
\[ m_{\text{recyclable}} = 1,8 + 0,9 + 1,0 = \SI{3,7}{kg}. \]
\item Pourcentage de bakélite :
\[ \frac{0,3}{4,0} \times 100 = 7,5\,\%. \]
La bakélite est un \textbf{thermodurcissable} : chauffée, elle se décompose sans fondre ; on ne peut donc pas la refondre pour lui donner une nouvelle forme.
\item Circuit de valorisation du PET : les bouteilles vides sont triées à la maison puis vendues à un \textbf{collecteur} (ou déposées à un point de collecte) ; une \textbf{entreprise de recyclage} les lave, les broie en paillettes et les fond en granulés ; ces granulés servent à fabriquer des fibres textiles, des sacs ou de nouveaux objets plastiques, vendus ensuite par des \textbf{artisans} ou des industriels.
\end{enumerate}
\end{corrige}

\begin{corrige}[Exercice 9 — Sujet type BEPC n° 1 : les matières plastiques et leur classification]
\textbf{Partie A — Connaissances.}
\begin{enumerate}
\item \textbf{Monomère} : petite molécule qui, par assemblage avec de nombreuses molécules identiques, forme un polymère. \textbf{Polymère} : macromolécule (très grosse molécule) formée par la répétition d'un grand nombre de monomères. \textbf{Matière plastique} : matériau constitué d'un polymère, auquel on ajoute souvent des additifs, et que l'on peut mettre en forme.
\item Le monomère du polyéthylène est l'éthylène \ce{CH2=CH2} ; après ouverture de la double liaison, le motif s'écrit \ce{-CH2-CH2-}. Une chaîne de trois motifs :
\[ \ce{-CH2-CH2-CH2-CH2-CH2-CH2-} \quad \text{soit} \quad \ce{(-CH2-CH2-)_{3}} \]
\item \textbf{Thermoplastique} : plastique qui se ramollit et fond quand on le chauffe, puis durcit en refroidissant (exemple : PE, PET, PVC, PP, PS). \textbf{Thermodurcissable} : plastique qui durcit définitivement lors de sa fabrication et se décompose sans fondre quand on le chauffe (exemple : bakélite).
\end{enumerate}
\textbf{Partie B — Expérience.}
\begin{enumerate}
\setcounter{enumi}{3}
\item X et Y fondent sans se décomposer : ce sont des \textbf{thermoplastiques} (de températures de fusion différentes : \SI{130}{\celsius} et \SI{250}{\celsius}). Z noircit sans fondre : c'est un \textbf{thermodurcissable}.
\item La poignée de casserole doit résister à la chaleur sans fondre : c'est \textbf{Z} (thermodurcissable) qui convient, car il reste rigide et ne fond pas.
\item \textbf{X et Y} peuvent être recyclés par refonte. Procédé : (1) on trie, lave et broie les déchets thermoplastiques ; (2) on les fond puis on les moule (ou extrude) en granulés ou en nouveaux objets, qui durcissent en refroidissant.
\end{enumerate}
\textbf{Partie C — Citoyenneté.}
\begin{enumerate}
\setcounter{enumi}{6}
\item Les plastiques abandonnés bouchent les caniveaux, ce qui provoque des inondations en saison des pluies. Les eaux stagnantes favorisent la prolifération des moustiques, vecteurs du paludisme. De plus, brûler les plastiques à l'air libre dégage des fumées toxiques qui provoquent des maladies respiratoires.
\end{enumerate}
\textbf{Barème indicatif (sur 20)} : définitions A1 : 3 pts ; motif et chaîne A2 : 3 pts ; définitions et exemples A3 : 3 pts ; classification B4 : 3 pts ; choix justifié B5 : 2 pts ; recyclage B6 : 3 pts ; enjeu sanitaire C7 : 3 pts.

\textbf{Points de vigilance.} Ne pas écrire que le thermodurcissable « fond » : il se \emph{décompose}. Pour A2, la double liaison \ce{C=C} doit disparaître dans le motif (simples liaisons \ce{C-C}). Toute classification doit être \emph{justifiée} par l'observation expérimentale (fond / ne fond pas).
\end{corrige}

\begin{corrige}[Exercice 10 — Sujet type BEPC n° 2 : les plastiques dans la vie quotidienne]
\textbf{Partie A — Identification.}
\begin{enumerate}
\item Tableau complété :
\begin{center}
\begin{tabularx}{\linewidth}{|l|X|X|}
\hline
\rowcolor{popblueL} \textbf{Objet} & \textbf{Plastique} & \textbf{Propriété justifiant le choix} \\
\hline
Bouteille de boisson gazeuse & PET & Transparence, étanchéité aux gaz, légèreté \\
\hline
Tuyau de plomberie & PVC & Rigidité, résistance à l'eau et aux produits chimiques \\
\hline
Gobelet jetable & PS & Légèreté, faible coût, facilité de moulage \\
\hline
Chaise plastique & PP & Rigidité, résistance aux chocs \\
\hline
Sachet de marché & PE & Souplesse, finesse, légèreté \\
\hline
Interrupteur électrique & Thermodurcissable & Résistance à la chaleur, infusibilité, propriété isolante \\
\hline
\end{tabularx}
\end{center}
\end{enumerate}
\textbf{Partie B — Avantages et limites.}
\begin{enumerate}
\setcounter{enumi}{1}
\item Trois avantages : les plastiques sont \textbf{légers}, \textbf{bon marché} et \textbf{faciles à mettre en forme} ; on peut ajouter : imperméables à l'eau, isolants électriques, résistants à la corrosion.
\item Deux limites : ils sont \textbf{non biodégradables} et persistent des siècles dans la nature (pollution durable) ; leur \textbf{combustion dégage des gaz toxiques} dangereux pour la santé.
\end{enumerate}
\textbf{Partie C — Calcul.}
\begin{enumerate}
\setcounter{enumi}{3}
\item Masse de granulés récupérée :
\[ m = 12 \times \frac{85}{100} = \SI{10,2}{tonnes} \text{ de PET recyclé par mois.} \]
\item Énergie économisée : $\SI{10,2}{t} = \SI{10200}{kg}$, donc
\[ E = \num{10200} \times 2 = \SI{20400}{kWh} \text{ économisés chaque mois.} \]
\item Conclusion : le recyclage du PET permet d'économiser une quantité considérable d'énergie (ici \SI{20400}{kWh} par mois), tout en réduisant la pollution : c'est un double bénéfice pour la société.
\end{enumerate}
\textbf{Barème indicatif (sur 20)} : tableau A1 : 6 pts (1 pt par ligne) ; avantages B2 : 3 pts ; limites B3 : 2 pts ; calcul de masse C4 : 3 pts ; calcul d'énergie C5 : 4 pts ; conclusion C6 : 2 pts.

\textbf{Points de vigilance.} Dans le tableau, la propriété citée doit correspondre à l'usage de l'objet (ex. transparence pour la bouteille, pas « souplesse »). Dans C5, bien convertir les tonnes en kilogrammes avant de multiplier par \SI{2}{kWh/kg} : c'est l'erreur la plus fréquente.
\end{corrige}

\begin{corrige}[Exercice 11 — Sujet type BEPC n° 3 : pollution plastique et gestion des déchets]
\textbf{Partie A — Étude de document.}
\begin{enumerate}
\item Les déchets plastiques jetés dans les rues bouchent les caniveaux ; l'eau de pluie ne s'écoule plus, ce qui provoque des inondations. Les eaux stagnantes deviennent des lieux de ponte des moustiques, qui transmettent le paludisme en piquant les habitants.
\item Une bouteille en PET est un polluant « durable » car elle est non biodégradable : sa décomposition dans la nature prend de \textbf{100 à \num{1000} ans}. Elle pollue donc les sols et les cours d'eau pendant plusieurs siècles, soit plusieurs générations.
\end{enumerate}
\textbf{Partie B — Les 3R.}
\begin{enumerate}
\setcounter{enumi}{2}
\item Les trois « R » signifient : \textbf{Réduire} (consommer moins de plastique), \textbf{Réutiliser} (se servir plusieurs fois du même objet), \textbf{Recycler} (transformer le déchet en nouvelle matière première).
\item (a) Gourde au lieu de bouteilles jetables : \textbf{Réduire}. (b) Bouteilles vides transformées en arrosoirs : \textbf{Réutiliser}. (c) Bouchons fondus pour fabriquer des pavés : \textbf{Recycler}. (d) Refuser un sachet au marché : \textbf{Réduire}.
\end{enumerate}
\textbf{Partie C — Projet de quartier.}
\begin{enumerate}
\setcounter{enumi}{4}
\item Masse recyclable par semaine :
\[ m_{\text{sem}} = 500 \times \frac{60}{100} = \SI{300}{kg}. \]
Par mois (4 semaines) :
\[ m_{\text{mois}} = 300 \times 4 = \SI{1200}{kg} = \SI{1,2}{t}. \]
\item Recette mensuelle de la coopérative :
\[ R = \num{1200} \times 50 = \SI{60000}{FCFA} \text{ par mois.} \]
\item Deux actions complémentaires : organiser des \textbf{campagnes de sensibilisation} (portes-à-porte, affiches, slogans) pour que les habitants trient leurs plastiques et refusent les sachets ; installer des \textbf{poubelles de tri} et des points de collecte dans le quartier, et plaider auprès de la mairie pour le curage régulier des caniveaux.
\end{enumerate}
\textbf{Barème indicatif (sur 20)} : lien pollution--inondation--paludisme A1 : 3 pts ; durabilité du PET A2 : 2 pts ; signification des 3R B3 : 3 pts ; classement des gestes B4 : 4 pts ; calculs de masses C5 : 3 pts ; recette C6 : 3 pts ; actions C7 : 2 pts.

\textbf{Points de vigilance.} Pour B4, ne pas confondre \emph{réutiliser} (même objet, nouvel usage, sans transformation industrielle) et \emph{recycler} (transformation en nouvelle matière par fonte ou broyage). Pour C6, vérifier l'unité monétaire et poser clairement la multiplication ; le résultat doit être accompagné de son unité (FCFA).
\end{corrige}

\newpage

\coursec{Fiche bilan}

\begin{popfigure}
\begin{center}
\begin{tikzpicture}[mindmap, grow cyclic,
  every node/.style={concept, font=\small},
  root concept/.append style={concept color=popblue, font=\bfseries},
  level 1 concept/.append style={sibling angle=60, level distance=3.4cm, font=\footnotesize},
  level 2 concept/.append style={level distance=2.2cm, font=\scriptsize}]
\node[root concept]{Matières plastiques}
  child[concept color=poporange]{ node {Origine et définition}
    child { node {Pétrole (SONARA)} }
    child { node {Macromolécules : polymères} }
  }
  child[concept color=popgreen]{ node {Classification}
    child { node {Thermoplastiques : ramollissent à chaud} }
    child { node {Thermodurcissables : durcissent à jamais} }
  }
  child[concept color=poppurple]{ node {Plastiques courants}
    child { node {PE, PP, PVC, PET, PS} }
    child { node {Usages : sachets, tuyaux, bouteilles} }
  }
  child[concept color=popgold]{ node {Avantages et limites}
    child { node {Légers, isolants, peu chers} }
    child { node {Non biodégradables, inflammables} }
  }
  child[concept color=poppink]{ node {Pollution}
    child { node {Sachets, bouteilles : sols, drains, mer} }
    child { node {Durée : des centaines d'années} }
  }
  child[concept color=popdark]{ node {Gestion : les 3R}
    child { node {Réduire, Réutiliser, Recycler} }
    child { node {Tri, collecte, valorisation} }
  };
\end{tikzpicture}
\end{center}
\legende{Carte mentale de la leçon : les six idées clés sur les matières plastiques.}
\end{popfigure}

\begin{bilan}
\textbf{Définitions clés}
\begin{itemize}
  \item Une \cle{matière plastique} est un matériau organique, fabriqué par l'homme, constitué de très longues molécules appelées \cle{polymères} (macromolécules formées par la répétition d'une même unité, le \cle{monomère}).
  \item La matière première des plastiques est le \cle{pétrole} (ou le gaz naturel), transformé par raffinage puis par \cle{polymérisation}.
  \item Un \cle{thermoplastique} se ramollit à la chaleur et durcit en refroidissant, de façon \textbf{réversible} : il est recyclable par fusion (ex. : PE, PP, PVC, PET).
  \item Un \cle{thermodurcissable} durcit \textbf{définitivement} lors de sa première mise en forme à chaud : il ne peut plus être refondu (ex. : bakélite des poignées de marmites).
\end{itemize}

\textbf{Plastiques courants et usages}
\begin{center}
\begin{tabularx}{\linewidth}{|l|X|X|}
\hline
\rowcolor{popblueL} \textbf{Plastique} & \textbf{Propriété remarquable} & \textbf{Usages courants} \\
\hline
PE (polyéthylène) & Souple, léger & Sachets, bidons, films \\
\hline
PP (polypropylène) & Résistant à la chaleur & Chaises, seaux, marmites \\
\hline
PVC & Rigide ou souple & Tuyaux, gaines électriques \\
\hline
PET & Transparent, imperméable & Bouteilles d'eau et de boissons \\
\hline
PS (polystyrène) & Léger, isolant & Barquettes, emballages \\
\hline
\end{tabularx}
\end{center}

\textbf{Méthodes express}
\begin{itemize}
  \item \cle{Identifier} un thermoplastique : il se déforme et fond quand on le chauffe (sachet qui se rétracte près d'une flamme).
  \item \cle{Reconnaître} un thermodurcissable : il noircit et se carbonise sans fondre (poignée de casserole).
  \item \cle{Gérer} les déchets plastiques avec la règle des \textbf{3R} : \textbf{Réduire} la consommation, \textbf{Réutiliser} les objets, \textbf{Recycler} par tri, broyage, fusion et nouvelle mise en forme.
  \item \cle{Sécurité} : ne jamais brûler les plastiques à l'air libre (gaz toxiques) ; ne pas jeter les sachets dans la nature (bouchage des caniveaux, inondations).
\end{itemize}
\end{bilan}

\begin{aretenir}[Checklist avant le BEPC]
\begin{itemize}
  \item $\square$ Je sais définir une matière plastique, un polymère et un monomère.
  \item $\square$ Je sais que la matière première des plastiques est le pétrole (raffinage, SONARA à Limbé).
  \item $\square$ Je distingue thermoplastique (ramollit à chaud, recyclable) et thermodurcissable (durcit à jamais).
  \item $\square$ Je cite au moins quatre plastiques courants (PE, PP, PVC, PET) et un usage de chacun.
  \item $\square$ Je sais identifier un thermoplastique par une expérience simple de chauffage.
  \item $\square$ Je cite trois avantages des plastiques (légers, isolants, bon marché, résistants à la corrosion).
  \item $\square$ Je cite deux limites : non biodégradables, inflammables, gaz toxiques à la combustion.
  \item $\square$ Je connais la durée de vie d'un plastique dans la nature : des centaines d'années.
  \item $\square$ Je connais la règle des 3R : Réduire, Réutiliser, Recycler.
  \item $\square$ Je sais décrire les étapes du recyclage : tri, broyage, lavage, fusion, nouvelle mise en forme.
  \item $\square$ Je respecte les consignes de sécurité : pas de combustion des plastiques, pas de jet dans la nature.
  \item $\square$ Je sais relier la pollution plastique aux inondations et aux maladies au Cameroun.
\end{itemize}
\end{aretenir}

\findecours

\end{document}
